\documentclass[format=acmsmall, review=false]{acmart} % EC formatting
\usepackage{acm-ec-26} % EC formatting
\usepackage{style} % My style macros
\newcommand{\timeind}{t}

\newcommand{\parahead}[1]{\vspace{0.1in}\noindent\textbf{#1. }}
\newcommand{\subparahead}[1]{\vspace{0.07in}\noindent\emph{#1. }}
\newcommand{\paraheadwnogap}[1]{\vspace{0.0in}\noindent\textbf{#1. }}
\newcommand{\subparaheadwnogap}[1]{\vspace{0.0in}\noindent\emph{#1. }}


\newcommand{\competitora}{a}
\newcommand{\competitorb}{b}
\newcommand{\arbitrarycompetitor}{z}

% Valuations and utilities
\newcommand{\numdataa}{n}
\newcommand{\numdatab}{m}

\newcommand{\numdatafora}{i}
\newcommand{\numdataforatt}[1]{\numdatafora_{#1}}
\newcommand{\numdataforat}{\numdataforatt{t}}

\newcommand{\numdataforb}{j}
\newcommand{\numdataforbtt}[1]{\numdataforb_{#1}}
\newcommand{\numdataforbt}{\numdataforbtt{t}}

\newcommand{\vala}{V^\competitora}
\newcommand{\valaii}[1]{\vala_{#1}}
\newcommand{\valai}{\valaii{i}}

\newcommand{\valb}{V^\competitorb}
\newcommand{\valbjj}[1]{\valb_{#1}}
\newcommand{\valbj}{\valbjj{j}}


\newcommand{\setexchanges}{\Scal}
\newcommand{\exchange}{s}
\newcommand{\exchangett}[1]{s_{#1}}
\newcommand{\exchanget}{\exchangett{t}}

\newcommand{\exchangea}{i}
\newcommand{\exchangeb}{j}

\newcommand{\utila}{U^\competitora}
\newcommand{\utilass}[1]{\utila_{#1}}
\newcommand{\utilas}{\utilass{s}}

\newcommand{\utilb}{U^\competitorb}
\newcommand{\utilbss}[1]{\utilb_{#1}}
\newcommand{\utilbs}{\utilbss{s}}

\newcommand{\prior}{D}
\newcommand{\priora}{\prior^a}
\newcommand{\priorb}{\prior^b}

\newcommand{\compcoeff}{\alpha}

\newcommand{\welfarefunc}{W}

\newcommand{\reward}{Y}
\newcommand{\rewardss}[1]{\reward_{#1}}
\newcommand{\rewards}{\rewardss{s}}

% Algorithm formulations
% \newcommand{\sharingalg}{\pi}
\newcommand{\algstatemap}{f}
\newcommand{\algstatemaptt}[1]{f_{#1}}
\newcommand{\algstatemapt}{\algstatemaptt{t}}

\newcommand{\algstopmap}{\delta}
\newcommand{\algstopmaptt}[1]{\algstopmap_{#1}}
\newcommand{\algstopmapt}{\algstopmaptt{t}}

\newcommand{\altalgstopmap}{\delta'}
\newcommand{\altalgstopmaptt}[1]{\altalgstopmap_{#1}}
\newcommand{\altalgstopmapt}{\altalgstopmaptt{t}}

\newcommand{\algendtime}{T} % Used to be \tau
\newcommand{\altalgendtime}{T'}

\newcommand{\algstate}{s} % Used to be \sigma
\newcommand{\algstatett}[1]{\algstate_{#1}}
\newcommand{\algstatet}{\algstatett{t}}

\newcommand{\timefilt}{\Gcal}
\newcommand{\timefilttt}[1]{\timefilt_{#1}}
\newcommand{\timefiltt}{\timefilttt{t}}

\newcommand{\statefilt}{\Fcal}
\newcommand{\statefiltss}[1]{\statefilt_{#1}}
\newcommand{\statefilts}{\statefiltss{s}}

\newcommand{\alg}{A}
\newcommand{\altalg}{A'}
\newcommand{\classalgs}{\Acal}
\newcommand{\classicalgs}{\Acal^{\text{IC}}}

\newcommand{\strategy}{\sigma}
\newcommand{\strategytt}[1]{\strategy_{#1}}
\newcommand{\strategyt}{\strategytt{t}}

\newcommand{\altstrategy}{\sigma'}
\newcommand{\altstrategytt}[1]{\altstrategy_{#1}}
\newcommand{\altstrategyt}{\altstrategytt{\timeind}}

\newcommand{\stoptime}{\tau}
\newcommand{\altstoptime}{\tau'}


\newcommand{\apolicy}{\pi} % for some reason \policy seems taken?
\newcommand{\setpolicies}{\Pi}

\newcommand{\exchangepolicy}{\pi}
\newcommand{\altexchangepolicy}{\pi'}

\newcommand{\setexchangepolicies}{\setpolicies^{\text{ex}}}
\newcommand{\seticexchangepolicies}{\setpolicies^{\text{IC}}}

% osas - stopping time paper
\newcommand{\osasreward}{Z}
\newcommand{\osasrewardss}[1]{\osasreward_{#1}}
\newcommand{\osasrewards}{\osasrewardss{s}}

\newcommand{\osasstrategy}{\sigma}
\newcommand{\osasstrategytt}[1]{\osasstrategy_{#1}}
\newcommand{\osasstrategyt}{\osasstrategytt{t}}

\newcommand{\osasstoptime}{\tau}

\newcommand{\osaspolicy}{\pi}
\newcommand{\osasstate}{s}
\newcommand{\osasstates}{S}

\newcommand{\osaspolicies}{\Pi}
\newcommand{\osaspoliciesss}[1]{\osaspolicies_{#1}}
\newcommand{\osaspoliciess}{\osaspoliciesss{s}}

\newcommand{\osassnellenv}{X}
\newcommand{\osassnellenvss}[1]{\osassnellenv_{#1}}
\newcommand{\osassnellenvs}{\osassnellenvss{s}}

\newcommand{\osasstatefilt}{\Fcal}
\newcommand{\osasstatefiltss}[1]{\osasstatefilt_{#1}}
\newcommand{\osasstatefilts}{\osasstatefiltss{s}}

\newcommand{\osasdirectsuccessorss}[1]{U(#1)}

\newcommand{\snellenv}{X}
\newcommand{\snellenvss}[1]{\snellenv_{#1}}
\newcommand{\snellenvs}{\snellenvss{s}}


\newcommand{\nextexchange}{u}

\newcommand{\directsuccessorss}[1]{U(#1)}

% \newcommand{\optalg}{\tilde{\alg}}
\newcommand{\optalg}{\alg}
% \newcommand{\optexpolicy}{\tilde{\apolicy}}
\newcommand{\optexpolicy}{\apolicy}

% \newcommand{\optstrategy}{\tilde{\sigma}}
\newcommand{\optstrategy}{\sigma}
\newcommand{\optstrategytt}[1]{\optstrategy_{#1}}
\newcommand{\optstrategyt}{\optstrategytt{t}}

% \newcommand{\optstoptime}{\tilde{\tau}}
\newcommand{\optstoptime}{\tau}

\newcommand{\zerotuple}{(0,0)}
% \newcommand{\zeroabrev}{}

\newcommand{\mbleexistsalgstatemap}{\tilde{\algstatemap}}
\newcommand{\mbleexistsalgstatemapt}{\tilde{\algstatemapt}}

% Notation for truthfulness proof
\newcommand{\opttruthexpolicy}{\pi'}
\newcommand{\opttruthexpolicykk}[1]{\opttruthexpolicy_{#1}}
\newcommand{\opttruthexpolicyk}{\opttruthexpolicykk{k}}

\newcommand{\opttruthstrategy}{\sigma'}
\newcommand{\opttruthstrategytt}[1]{\opttruthstrategy_{#1}}
\newcommand{\opttruthstrategyt}{\opttruthstrategytt{t}}

\newcommand{\opttruthstoptime}{\tau'}

\newcommand{\stopstateforpolicy}[1]{\alpha\rbr{#1}}

\newcommand{\fixedpath}{P}

\newcommand{\icsnellenv}{X}
\newcommand{\icsnellenvpp}[1]{\icsnellenv_{#1}}
\newcommand{\icsnellenvp}{\icsnellenvpp{\fixedpath}}

\newcommand{\rewardbound}{C}

% \newcommand{\policyfollowspath}[2]{#1\rbr{#2}}
\newcommand{\policyfollowspath}[2]{#2\rbr{#1}}
% \newcommand{\policyfollowspathpp}[2]{#2\rbr{#1}}

\newcommand{\exchangeindexrange}[2]{\exchangett{#1:#2}}
\newcommand{\strategyindexrange}[2]{\strategytt{#1:#2}}

\newcommand{\existspolicymatchingpath}[1]{E\rbr{#1}}

\newcommand{\validicextensions}{\Ucal}
\newcommand{\validicextensionspp}[1]{\validicextensions_{#1}}

\newcommand{\expolicymatchicopt}{\pi}
% \newcommand{\expolicymatchicopt}{\pi}

\newcommand{\lastfollowtime}{\xi}
\newcommand{\lastfollowtimeii}[1]{\lastfollowtime_{#1}}
\newcommand{\lastfollowtimei}{\lastfollowtimeii{i}}

\newcommand{\lastcommonfollowtime}{\chi}

\newcommand{\chosenpath}{\iota}

% Choose a citation style by commenting/uncommenting the appropriate line:
\setcitestyle{acmnumeric}
% \setcitestyle{authoryear}

\newcommand{\kkcomment}[1]{\textcolor{magenta}{(#1 --KK)}}

% Title. Note the optional short title for running heads. In the interest of anonymization, please do not include any acknowledgements.
% \title[Exchange of Resources Between Two Uncertain Competitors]{Exchange of Resources Between Two Uncertain Competitors}
% \title[Exchanging Data Between Two Competitors with Uncertain Valuations]{Exchanging Data Between Two Competitors with Uncertain Valuations}
% \title[Exchanging Freely-replicable Goods Between Two Competitors with Uncertain Valuations]{Exchanging Freely-replicable Goods Between Two Competitors with Uncertain Valuations}
% Need to be careful when we just say 'resources', since resources could mean anything. We want to convey that it is the amount of goods that matter (not just which resources)
\title[Incremental Resource Exchange between Two Competitors under Valuation Uncertainty]{Incremental Resource Exchange between Two Competitors under Valuation Uncertainty}





% Anonymized submission.
\author{Submission 385}

\begin{abstract}
% Competing entities can often benefit mutually by sharing resources, such as datasets  or access to experimental equipment for research and 
% development. However, uncertainty about an agent's valuation of another's resources can make participants reluctant to share 
% even when mutual benefit is likely, due to fears of imbalanced realized benefits. For example, the value of access to a competitor's 
% dataset or time on experimental equipment is typically revealed only through use and is difficult to assess \emph{ex ante}. In this 
% work, we study the design of exchange algorithms for two competing agents whose valuations are uncertain and revealed incrementally 
% through interaction. We focus on resources whose value depends on the amount shared (e.g. volume of data, hours of equipment use) 
% and is inherently non-retractable (i.e. once data is transmitted or time on equipment is used, it cannot be reclaimed).

% We introduce multi-round \emph{exchange algorithms} that propose sequences of increasing exchanges, leveraging information 
% from prior rounds to improve expected welfare. Our first and contribution is an optimal multi-exchange algorithm under minimal 
% assumptions on valuations, utilities, and welfare functions, obtained by reducing the design problem to an optimal stopping 
% problem over a lattice. The resulting algorithm is also individually rational for reasonable choices of welfare functions.
% %
% We then study a significantly more difficult setting in which agents may exit early. Under this model,
% we define \emph{incentive-compatible} exchange algorithms, to be
% those in which no agent can benefit from leaving prematurely.
% Under added assumptions on the environment, we construct an exchange algorithm which is optimal
% among all incentive-compatible exchange algorithms.


% ---- new

Competing entities can often mutually benefit by sharing resources, such as datasets or access to experimental equipment. 
However, agents are often uncertain about the value they will realize from the others' resources.
For example, the value of access to a competitor's dataset or time on experimental equipment is typically revealed only 
through use and is difficult to assess \emph{ex ante}. This uncertainty can make agents reluctant to share even when mutual 
benefit is likely, due to fears of imbalance in realized benefits; an agent may be concerned that a competitor will 
derive more value from her resources than she would from theirs.

In this work, we study the design of exchange algorithms for two competing agents whose valuations for the other's 
resources are uncertain and only revealed incrementally through interaction. We focus on resources whose value 
depends on the amount shared (e.g. volume of data, hours of equipment use) and is inherently non-retractable 
(i.e. once data is transmitted or time on equipment is used, it cannot be reclaimed). We introduce multi-round 
\emph{exchange algorithms} that propose sequences of increasing exchanges, leveraging information from prior rounds 
to improve expected welfare. Our first contribution is an optimal and individually rational multi-round exchange 
algorithm, obtained by reducing this problem to an optimal stopping problem over the lattice $\NN^2_{\geq 0}$. 
We then study the 
setting in which agents may exit early, defining \emph{incentive-compatible} exchange algorithms to be those
in which no agent can benefit from leaving prematurely.
Under additional mild assumptions on the environment, we construct an exchange algorithm which is optimal
among incentive-compatible exchange algorithms.

\end{abstract}

% Abstract. Note that this must come before \maketitle.
% \begin{abstract}
%     Competition often discourages companies, labs, and governments from sharing valuable resources like 
%     data or equipment. Yet, collaboration can still be mutually beneficial and even socially
%     desirable.
%     However, it may be unclear 
%     if collaborating is worthwhile when
%     competitors do not know their value for each others' resources until using them.
%     For instance, the value of a dataset may only become apparent after training a machine learning model on it.

%     We formalize the problem of determining how competitors should exchange resources given
%     uncertainty
%     as designing a multi-round resource exchange algorithm.
%     We reduce the task of designing an optimal exchange algorithm to solving an optimal stopping 
%     problem over the 
%     lattice $\NN^2_{\geq 0}$, adapting ideas from optimal stopping over partially ordered sets.
%     This algorithm is individually rational for fair choices of welfare functions.
%     We then extend our model to a significantly more challenging setting: when competitors
%     can exit the sharing algorithm early 
%     (as opposed to opting in or out at the beginning). 
%     Algorithms for which
%     it is never beneficial for either competitor to leave early are called incentive compatible (IC).
%     Under reasonable conditions on competitors' valuations, 
%     we determine the performance of an optimal IC algorithm. 
%     % if it exists?
%     % we give a construction for an optimal IC algorithm.

% \end{abstract}

%\newcommand{\kkcomment}[1]{\textcolor{magenta}{(#1 --KK)}}
\newcommand{\note}[1]{\textcolor{blue}{(#1)}}
\newcommand{\mycomment}[1]{\textsf{\textcolor{red}{(#1)}}}
\newcommand{\todo}[1]{\textsf{\textcolor{purple}{(Todo: #1)}}}

\begin{document}

% Title page for title and abstract only.
\begin{titlepage}

\maketitle

% Optionally include a table of contents
\vspace{1cm}
\setcounter{tocdepth}{1} % adjust to 1 if desired
\tableofcontents

\end{titlepage}


\section{Introduction}
\label{sec:intro}

Competition is a defining feature of many multi-agent environments in which valuable resources are held by firms, research laboratories, and government institutions. Even in competitive settings, resource sharing can generate substantial mutual and social benefits. For instance, in data-driven scientific domains such as materials design, drug discovery, and chip design, pooling datasets across organizations increases statistical power, accelerates the identification of promising designs, and reduces redundant experimentation~\cite{DataReuseOpenDataAdvantage, Villanueva2021,resonac_joint3_2025}. Similarly, open collaboration on clinical research data has been shown to facilitate secondary analyses, improve public health insights, and inform new hypotheses beyond the scope of any single study~\cite{NCBI_SharingClinicalData, NatGeneticsDataSharing}. In academic and scientific contexts, shared access to specialized experimental facilities and high-performance computing infrastructure can enable studies that would be infeasible for a single group to carry out independently~\citep{ambulos2025greater,cuboulder_greenlabs_2024,orduna2014sharing}. More broadly, many competitive resource-sharing environments are nonzero-sum: when agents possess complementary assets, joint use can increase total surplus, allowing all participants to mutually benefit despite competition, with these benefits extending to society at large.

\parahead{Uncertainty in valuations}
However, a central challenge in incentivizing agents to share in such settings is uncertainty about how much an agent values another agent’s resources. For example, the economic value of access to a competitor’s dataset or time on specialized equipment is difficult to assess \emph{ex ante} and is typically revealed only through use.
For instance, the utility of a dataset is determined by the performance of models trained using this dataset in downstream analyses; this information is generally unavailable without direct access to the dataset itself. Similarly, the value of access to a competitor’s laboratory equipment often hinges on whether the resulting experiments produce novel academic insights or commercially viable outcomes.

Hence, a key challenge in resource sharing is the possibility of imbalance in realized values. Precisely, due to this uncertainty, an agent may be concerned that competitors will derive more value from her resources than she will from theirs. Standard market mechanisms and contractual approaches are ill-suited to such environments, as they generally cannot accommodate settings in which one's value for a competitor's resource is learned only through interaction.



\parahead{This work}
We study the design of exchange algorithms, \emph{without monetary transfers}, that incentivize resource sharing between \emph{two} 
competing agents. We focus on resources whose valuation can be expressed as a function of the \emph{amount} shared 
(e.g. volume of data or number of hours of access to laboratory equipment) and for which sharing is inherently \emph{non-retractable} 
(i.e. once data have been transmitted or time on a machine has been used, the original agent cannot reclaim those units).
The latter feature creates a fundamental tension: valuations are revealed only through sharing, yet any over-sharing 
that makes an agent worse off cannot be reversed.
%
We focus on algorithms that enable \emph{incremental} collaboration (e.g. sharing small amounts of data, granting access 
to specialized equipment for a limited time), which allows agents to learn each other's valuations over multiple rounds. 
Our goal is to design exchange algorithms that achieve socially desirable outcomes by sustaining cooperation when it is mutually beneficial
and terminating interaction when it is not.


\subsection{Summary of Our Contributions}

We formalize this problem via a novel model and study resource exchange in two settings: first, when agents cannot withdraw from 
a multi-round algorithm until the algorithm terminates, and second, when agents can withdraw whenever they wish. 
We now summarize our model and results.


\parahead{Model}
There are two competing agents, $\competitora$ and $\competitorb$, each possessing 
collections of $\numdatab$ and $\numdataa$ 
different goods
respectively (e.g. quantity of data, number of hours an equipment can be lent).
Both agents derive value from each others' goods based on the quantity
received, quantified by sequences $\{\valai\}_{i\in[\numdataa]}$ and $\{\valbj\}_{j\in[\numdatab]}$,
% $\vala:\cbr{0,\ldots,\numdataa}\rightarrow\RR$ and 
% $\valb:\cbr{0,\ldots,\numdatab}\rightarrow\RR$,
where $\valai$ is the value  of agent $a$ when she receives $i$ of $b$'s goods, and likewise for $\valbj$.
An exchange 
\(
    \rbr{\exchangea,\exchangeb}
    =
    \exchange\in\setexchanges
    =
    \cbr{0,\ldots,n}\times\cbr{0,\ldots,m}
\)
corresponds to $\competitora$ receiving $\numdatafora$ of $\competitorb$'s resources and 
$\competitorb$ receiving $\numdataforb$ of $\competitora$'s resources.
For a particular exchange 
\(
    \exchange = \rbr{\exchangea,\exchangeb}
    \in\setexchanges
    % =
    % \cbr{0,\ldots,n}\times\cbr{0,\ldots,m}
\)
let
\(
    \utilas = \utila\big(\valai,\valbj\big)
\)
and
\(
    \utilbs = \utilb\big(\valai,\valbj\big)
\)
% let $\utila:\setexchanges\rightarrow\RR$ and $\utilb:\setexchanges\rightarrow\RR$ denote
be the utilities for both agents,
where
\(
    \utila,\utilb:\RR^2\rightarrow\RR
\)
are functions mapping the value obtained by both competitors to their utilities.
These interactions may be non-zero-sum (i.e. $\utilas \neq -\utilbs$), 
and the utility gains for the agents may be asymmetric (one agent could benefit more from collaboration than the other).
%Moreover, while the canonical setting of interest involves negative externalities due to competition (i.e., $\utila$ decreases in $\valbj$), our model does not preclude positive externalities (i.e., $\utila$ may increase with $\valbj$).

We capture uncertainty in valuations by modeling $\vala$ and $\valb$ as independent random draws from \emph{known} distributions $\priora, \priorb$. Neither player observes their or the others' realized valuations $\valai$ and $\valbj$ for any $i,j$ until the corresponding goods have been exchanged.




\subparahead{Welfare maximization under uncertainty}
Let $\welfarefunc:\RR^2\rightarrow\RR$ be a known welfare function used to quantify the quality of an exchange
$\exchange\in\setexchanges$ via $\welfarefunc\rbr{\utilas,\utilbs}$.
Common examples include the egalitarian welfare~\citep{mas1995microeconomic} 
$\welfarefunc(x,y)=\min(x,y)$, the 
Nash welfare~\citep{nash1950bargaining} $W(x,y)=xy$, or the social welfare subject to individual rationality 
$\welfarefunc(x,y)=(x+y)\cdot 1_{\cbr{x,y\geq 0}}$.
We are interested in designing  
algorithms in which a central planner (a trusted and neutral third party) facilitates an exchange $\exchange\in\setexchanges$ to maximize the welfare, while ensuring that both agents benefit from the exchange.

A straightforward option is to choose the exchange $s$ which maximizes the ex ante welfare, i.e.
$\argmax_{
        \exchange\in\setexchanges
    }
    \EE[\welfarefunc(\utilas,\utilbs)      
    ] $.
However, instead of proposing a one time exchange, an alternative setting is where the planner proposes a sequence of incremental multi-round exchanges, with both competitors reporting their observed values after each round.
By accounting for this feedback, it is possible to outperform the above ex ante baseline. This is demonstrated via the following example.


\begin{restatable}{example}{multi-vs-single-round-alg}
    \label{ex:multi-vs-single-round-alg}
    \emph{
    Let $\numdataa=\numdatab=1$, so $\setexchanges=\cbr{0,1}^2$. Suppose that each competitor's value for the other's good is 
$\valaii{1},\valbjj{1}\drawniid\Bern(1/2)$ and assume $\valaii{0}=\valbjj{0}=0$.
    Let the utilities and welfare be given by
    \begin{align*}
        \utilas
        =
        \valai-\frac{1}{4}\valbj,
        \quad 
        \utilbs
        =
        \valbj-\frac{1}{4}\valai,
        \quad
        \text{and}
        \quad
        \welfarefunc(x,y)=\min(x,y) .
    \end{align*}  
    A direct computation shows that 
    $(1,1)\in\setexchanges$ is the single exchange that maximizes the expected welfare:
    \begin{align*}
        \argmax_{
            \exchange\in\cbr{0,1}^2
        }
        \EE\sbr{
            \min\rbr{
                \utilas,\utilbs
            }
        } =
                \rbr{1,1} 
        \quad
        \text{and}
        \quad
        \EE\sbr{
            \min\rbr{
                \utilass{(1,1)},\utilbss{(1,1)}
            }
        } 
        =
        \frac{1}{16} .
    \end{align*}  
    }

    \emph{
    Now consider a multi-exchange algorithm that works as follows. It first gives one agent the other's resource; without loss of generality assume 
    $\competitora$ receives $\competitorb$'s resource. 
    Competitor $\competitora$ is then asked to report their value of the received resource
    (either $0$ or $1$). If $\competitora$ (truthfully) reports $1$, then $\competitorb$ is 
    given $\competitora$'s resource. 
    However, if $\competitora$ reports a $0$, then $\competitorb$ is given nothing.
    Let $\exchange'=(i',j')\in\cbr{0,1}^2$
    denote the final exchange in this algorithm. A direct computation shows that this algorithm achieves a higher expected welfare than the above static policy. We have:
    \begin{align*}
        \EE\sbr{
            \min\rbr{
                \utilass{s'},\utilbss{s'}
            }
        } 
        &=
        \EE\sbr{
            \min\rbr{
                \valaii{i'}
                -
                \frac{1}{4}
                \valbjj{j'}
                ,
                \valbjj{j'}
                -
                \frac{1}{4}
                \valaii{i'}
            }
        } 
        \\
        &=
        \frac{1}{2}
        \rbr{
            \frac{1}{2}
            \min\rbr{
                1-\frac{1}{4}
                ,
                1-\frac{1}{4}
            }
            +
            \frac{1}{2}
            \min\rbr{
                1-\frac{1}{4}
                ,
                -\frac{1}{4}
            }
        }
        +
        \frac{1}{2}\min(0,0)
        % \\
        % &=
        =
        \frac{1}{8} .
    \end{align*}  
    }
\end{restatable}  


\subparahead{An algorithm for resource exchange}
Motivated by the above example, we study multi-round exchange algorithms that share goods incrementally over time. 
An algorithm consists of a \emph{stopping rule} and an \emph{exchange recommendation rule}. In each round $t$, 
the stopping rule determines whether the interaction continues or terminates. If it continues, the exchange 
recommendation rule proposes an exchange $(i_{t}, j_{t}) \in \setexchanges$ such that $i_{t} \ge i_{t-1}$ and $j_{t} \ge j_{t-1}$, 
with at least one of these inequalities holding strictly, i.e. at least one of the two agents receives an additional unit of the resource. 
The agents realize the valuations $\valaii{i_t}$ and $\valbjj{j_t}$, report them (truthfully) to the exchange algorithm which 
then uses this feedback to guide subsequent decisions.



\parahead{Setting 1: Agents cannot withdraw early}
In Section~\ref{sec:opt-non-ic-alg}, we first consider a setting where agents may choose to participate, but once they do, they cannot withdraw until the algorithm terminates. This restriction can be enforced via legal agreements or monetary penalties for early withdrawal.



\subparahead{Challenges}
Our setting introduces several challenges that do not arise in conventional resource-sharing protocols:
\emph{(i)}
In each round, the algorithm must make its decisions using only the information revealed by previous exchanges.
\emph{(ii)}
Decisions are irreversible: once goods have been shared, the original agent cannot reclaim those units.
\emph{(iii)}
The algorithm must account for competition between agents, as an exchange that improves one agent's outcome may reduce the other's.
\emph{(iv)}
These challenges are further compounded by the minimal assumptions we impose on valuation, utility, and welfare functions. For example, the random variables
\(
\{\valai\}_{\numdatafora=0}^{\numdataa}
\)
may be continuous, discrete, or arbitrarily correlated, and likewise for
\(
\{\valbj\}_{\numdataforb=0}^{\numdatab}
\).
The only assumption we impose on these random variables, as well as on the utility and welfare functions $\utilas$, $\utilbs$, and $\welfarefunc$, is boundedness.

% As a result, predicting the welfare impact of future exchanges is inherently difficult, making it challenging to design algorithms that repeatedly propose sequences of ex ante optimal exchanges based solely on information revealed in prior rounds.




\subparahead{Method and techniques}
We can view an exchange algorithm as a policy for traversing the lattice 
\(
    \setexchanges
    =
    \cbr{0,\ldots,\numdataa}
    \times
    \cbr{0,\ldots,\numdatab}
\), 
starting at $(0,0)$ and moving only up/right at each step, using the information 
reported from previous exchanges.
Each such traversal corresponds to an increasing sequence of exchanges, representing that competitors
cannot un-share resources.
Our key insight is that we can use this reformulation to reduce the design of an optimal exchange
algorithm to solving an optimal stopping time problem over the lattice $\setexchanges$~\citep{mandelbaum1981optimal}.
At a high level, the algorithm works as follows. On each round, we determine which of the following three options would lead to 
a higher expected welfare: (1) terminating the exchange, (2) giving $\competitora$ an additional one of 
$\competitorb$'s resources and then ``proceeding optimally'', or (3) giving $\competitorb$ an additional one of $\competitora$'s
resources and then ``proceeding optimally''.
We compute the expected welfare of ``proceeding optimally'' after a given exchanged via a dynamic program.
We show that under a simple assumption for $\welfarefunc$, that a non-negative expected welfare implies
non-negative expected utilities,
% is non-decreasing in agent valuations, 
the resulting optimal algorithm is 
also individually rational.



\subparahead{Impossibility under strategic reporting of values}
We assume that agents' valuations are reported truthfully or can be publicly verified by the central planner. 
It is intuitive that this assumption is necessary: without it, an agent could under-report her realized value to influence future allocations, 
either to receive more goods herself or to limit goods to her competitor. For example, in Example~\ref{ex:multi-vs-single-round-alg}, 
the agent who receives the good first could always report $0$, even if she realized a value of $1$, to prevent the other from receiving the good. 
Proposition~\ref{prop:opt-truthful-vs-opt-untruthful} formalizes this intuition, showing that the gap between the expected welfare of the optimal 
multi-exchange algorithm and that of the optimal truthful algorithm can be unbounded.

% Theorem XXX 

% The following theorem shows this assumption is necessary:

% \begin{restatable}{proposition}{opt-truthful-vs-opt-untruthful}
% \label{prop:untruthfulreporting}
% For any $M > 0$, there exists an instance 
% such the ratio between the expected welfare of the optimal 
% multi-exchange algorithm and that of the optimal truthful exchange algorithm is larger than $M$.
% \end{restatable}

% \kkcomment{Maybe, briefly outline proof idea here.}


\parahead{Setting 2: Agents can withdraw early}
Next, in Section~\ref{sec:opt-ic-alg}, we study a significantly more challenging setting in which agents can withdraw at any round, after which no further exchanges occur. We define \emph{incentive-compatible} (IC) exchange algorithms as those in which neither agent can benefit from leaving the exchange algorithm early.
IC does not come for free: an agent may choose to exit if she is fortunate in an early round and anticipates that future exchanges are more 
likely to benefit her competitor. For example, in Example~\ref{ex:multi-vs-single-round-alg}, the agent who receives the first good 
could leave early after realizing a value of $1$, preventing her competitor from obtaining a good. Consequently, the only IC policies 
are either to share nothing at all, yielding an expected utility of $0$, or to share both goods with both agents in the first round, yielding 
an expected utility of $1/16$, both of which are strictly worse than the optimal algorithm when participants cannot withdraw early.


% We explore the design of an optimal IC algorithm, and present conditions on the model under which techniques from optimal
% stopping and martingale theory can be applied to do so.
% We study how to design an optimal IC algorithm, and identify conditions on the model that allow us to use optimal 
% stopping and martingale methods to characterize it.

% We explore the design an optimal IC algorithm, and identify conditions on the model that allow us to use 
% optimal stopping and martingale methods to characterize such an algorithm.
We study designing an optimal IC algorithm and identify model conditions under which optimal 
stopping and martingale methods can be used to characterize it.
% We then give examples, illustrating the different challenges 
% \kkcomment{@Alex: finish this part up.} 
We also give examples showing the 
different challenges at play that complicate the design of an optimal incentive compatible algorithm and why
additional model assumptions are required to apply the aforementioned techniques. 


\vspace{0.1in}

\noindent
Finally, we note that our model and results in both settings can be straightforwardly extended to multi-resource exchanges.  In this case, 
competitor $a$ has $B_a$ ``buckets'' of resources with sizes $[m_1, m_2, \dots,$ $ m_{B_a}]$, and similarly for competitor $b$. 
An exchange $(i_t, j_t)$ then corresponds to a pair of vectors specifying the number of goods exchanged from each bucket.

% Finally, we note that our model and results extend naturally to multi-resource exchange settings, albeit with additional technical overhead. In this extension, each agent $a$ has $B_a$ resource ``buckets'' with capacities $[m_1, m_2, \dots, m_{B_a}]$, and similarly for agent $b$. An exchange $(i_t, j_t)$ is then specified by a pair of vectors indicating the number of goods transferred from each bucket.



% \kkcomment{@Alex: is it possible to guarantee the following `stability' condition, where, when the algorithm terminates, the agents will have no incentive to exchange among themselves. If so, it would be good to prove it and also mention it in the summary.}
% \mycomment{no, in the case that the welfare is misaligned, but I think yes for reasonable choices like the egalitatian welfare}

\subsection{Related work}

To the best of our knowledge, there is no prior work that directly addresses the problem we study. 
However, our work is inspired by and contributes to several related lines of research on resource sharing and 
exchange under uncertainty.

\paragraph{Classical models for exchange and allocation.} 
The allocation and exchange of resources among self-interested agents has a long history in economics. 
Early work on matching and exchange problems investigates mechanisms for assigning or trading goods without monetary transfers~\citep{varian1973equity,walras1900elements,fisher1891mathematical}. 
More broadly, mechanism design studies how to structure incentives so that agents reveal their true preferences, 
with applications in resource allocation and matching markets~\citep{gale1962college,shapley1974cores,abdulkadiroglu1999house}.

\paragraph{Exchange under uncertainty and learning.} 
A recent line of work extends classical exchange models to settings where agent utilities or preferences 
are uncertain and must be learned over time. These studies often leverage bandit and online learning techniques 
to approach efficient equilibria in matching markets~\citep{guo2021learning,jagadeesan2021learning,liu2020competing}. 
While these models typically do not consider direct competition and employ techniques distinct from ours, 
they underscore the importance of incorporating uncertainty into the design of exchange mechanisms.


\paragraph{Data sharing.} 
Recent literature on data markets examines how firms share or trade data, 
and how externalities influence outcomes. Some models explore optimal data auctions and mechanisms 
where allocations create positive or negative externalities among buyers~\citep{agarwal2024towards,gu2022data,bonatti2024selling}. 
Empirical work documents the emergence of data exchange agreements and the challenges of valuing data assets in practice~\citep{elsaify2021data}. 
Other studies focus on designing protocols for sharing data among non-competitive agents, often emphasizing 
fair contribution-based rewards~\citep{clintoncollaborative,sim2020collaborative,karimireddy2022mechanisms}.

\paragraph{Optimal stopping in auction and mechanism design.} 
Optimal stopping principles are widely used in  mechanism design, particularly in dynamic 
and sequential allocation settings. In dynamic auctions, designers often face decisions about when to allocate 
or price goods as agents' valuations or participation opportunities evolve; such problems naturally reduce 
to stopping or recursive decision frameworks~\citep{ashlagi2016sequential,samuel1984comparison,lucier2017economic}. 
Automated online mechanism design leverages optimal stopping and prophet inequalities to design approximately 
optimal online auctions, where the choice to accept a bid or wait for future offers is analogous to a stopping rule~\citep{hajiaghayi2007automated}. 
More broadly, our work connects to the study of \emph{dynamic mechanisms}, where goods are allocated across multiple stages. 
\citet{pavan2014dynamic} established results for incentive compatibility in these settings, showing that 
the stopping rule itself must be carefully designed to prevent manipulation of the timing of information disclosure. 
However, unlike all of the above works, our setting requires accounting for competition between agents and designing optimal stopping over a  multi-dimensional state space: 
the lattice of possible exchange quantities across competing agents.





% \begin{enumerate}
%     \item Intro
%     \begin{enumerate}
%         \item approx 3/4 pages describing real-world use cases
%         \item 1.1 Model: environment
%         \item 
%     \end{enumerate}
%     \item Model
%     \begin{enumerate}
%         \item Environment
%         \item Examples to motivate the policies
%         \item SEtting 1, desiderata
%         \item Setting 2, desiderata
%         \item Impossibility results (truthful reporting)
%     \end{enumerate}
%     \item Non-IC algo
%     \item IC algo
% \end{enumerate}
%\section{Introduction}

% Competition is a strong deterant for companies, research labs, and governmental insitutions to exchange
% resources. For example, pharaceutical companies may be hesitant to share data on different classes of 
% drugs with their competitors for fear of losing intellectual property. Similarly, research labs 
% may avoid sharing their equipment or findings,
% fearing that other labs would reap the rewards of publishing first. 
% However, it is often socially desirable for these insitutions to collaborate. In the aforementioned 
% examples, collaboration could expidite
% the development of novel pharaceuticals and scientific discoveries, both of which are beneficial
% to society.
% Moreover, while competition may deter the exchange of resources, it can be still be 
% mutually beneficial in non-zero sum settings.
% An important question is thus how to incentivize two competitors to share resources?

Competition often deters companies, research labs, and government institutions from sharing resources. 
For example, pharmaceutical companies may hesitate to share data on drug classes with competitors 
for feat of losing intellectual property. Likewise, research labs might avoid sharing equipment or findings to prevent 
others from publishing first. Despite this, collaboration is often socially beneficial. In the aforementioned cases, 
it could accelerate the development of new drugs and scientific discoveries. 
Moreover, while competition can hinder cooperation, resource sharing may still be mutually advantageous 
in nonzero sum settings. 

An obstacle to resource sharing between two competitors is that each may be uncertain of the value of the 
other's resources until after using them. For example, the utility of a dataset might only become 
clear once it's used to train a machine learning model. Similarly, the value of access to compute or equipment 
may depend on whether the resulting experiments or simulations produce novel findings. 
% This work explores how to incentivize 
The central question this work answers is how to incentivize
two competing agents to share resources when mutual benefit is likely, but their valuations for each 
others' resources are uncertain.

% A practical barrier to two competitors sharing resources is that they 
% may not know the value for 
% each others' resources until recieving and using them.
% For example if two competitors are exchanging different datasets, a competitor may not know their 
% value for the other's data until receiving it and attempting to use it to train a machine learning model.
% In the case of research labs exchanging access to compute or equipment, the value for this access 
% could be contingent upon the results of experiments or simulations run on these resources. If an
% experiment/simulation produced novel results then access was valuable but if it did not, then the access 
% was substantially less valuable.
% The main question of this work is how can we incentivize competing agents to share resources
% when it is likely mutually beneficial to do so, given uncertainty about their valuations for each
% others' resources.

% goal is to fin

% \begin{enumerate}

%     \item 
%     Give data examples

%     \item 
%     You don't know the value of compute time until you try to use it to do your experiments or 
%     simulations. If the simulations or experiments finish early, great, if not then need more compute.
%     Or find a negative vs positive result from running scientific experiments or simulations. 
%     For example your access to a piece of equipment was supper valuable if you got great results
%     but not very valuable if you didn't.

%     \item 

% \end{enumerate}  

% \begin{enumerate}
%     \item Paragraph about rising importance of data, unique from econ perspective that
%     once collected it can be freely shared or replecated.
    
%     \item 
%     Data sharing platforms exist but make unrealistic assumptions, agents willing 
%     to give away data for free (in same company) because don't assume competition.
%     Also assume things like truthful reporting.

%     \item 
%     Prior work has studied incentivizes in data sharing (truthfulness, fairness, )

%     \item In this work we study competition for strategic agents
    
% \end{enumerate}  

\subsection{Problem Motivation}

We will suppose that there are two competitors, $\competitora$ and $\competitorb$, each possessing 
collections of $\numdatab$ and $\numdataa$ 
different goods
respectively. Suppose that both agents would derive value from each others' goods based on the quantity
received.
Let this be quantified by random value functions
$\vala:\cbr{0,\ldots,\numdataa}\rightarrow\RR$ and 
$\valb:\cbr{0,\ldots,\numdatab}\rightarrow\RR$, where $\vala$ and $\valb$ are independent.


\begin{restatable}{example}{nme-val-functions}
    \label{ex:nme-val-functions}
    Suppose $\competitora$ and $\competitorb$ possess datasets
    \(
        \cbr{X^\competitora_i}_{i=1}^\numdatab
        \drawniid
        % \Ncal\rbr{\mu_\competitora,\sigma^2}
        D_a
    \)
    and
    \(
        % \quad
        % \text{and}
        % \quad
        \cbr{X^\competitorb_i}_{i=1}^\numdatab
        \drawniid
        % \Ncal\rbr{\mu_\competitorb,\sigma^2}
        D_b
    \)
    drawn from bounded distributions
    % where $D_a$ and $D_b$ are bounded.
    % respectively
    but wish to estimate the mean of each others' distributions. In this case, the value 
    functions may be defined as
    \begin{align*}
    % \(
        \vala_{\numdatafora}
        =
        % \rbr{
        \bigg(
            \frac{1}{k}
            \sum_{k=1}^\numdatafora
            X^b_k
            -
            \mu_b
        \bigg)^2
        % ~\forall \numdatafora\geq 1
        \quad
        \text{and}
        \quad
        % }^2
    % \),
    % \(
        \valb_{\numdataforb}
        =
        % \rbr{
        \bigg(
            \frac{1}{k}
            \sum_{k=1}^\numdataforb
            X^\competitora_k
            -
            \mu_\competitora
        \bigg)^2
        % ~\forall \numdataforb\geq 1
        % }^2
    % \),
    \end{align*}
    for $\numdatafora,\numdataforb\geq 1$, and $\vala_0=-C_a$, $\valb_0=-C_b$ 
    for sufficiently large $C_a,C_b\in\RR$,
    representing agents' values when they lack data.
\end{restatable}
% Beyond Example~\ref{ex:nme-val-functions} 

More generally, in machine learning, the value for data may be unknown until using it to train a model and
observing its performance. 

However, the agents may be hesitant to exchange resources if they are in competition. 
A exchange 
\(
    \rbr{\exchangea,\exchangeb}
    =
    \exchange\in\setexchanges
    =
    \cbr{0,\ldots,n}\times\cbr{0,\ldots,m}
\)
corresponds to $\competitora$ receiving $\numdatafora$ of $\competitorb$'s resources and 
$\competitorb$ receiving $\numdataforb$ of $\competitora$'s resources.
For a particular exchange 
\(
    \rbr{\exchangea,\exchangeb}
    =
    \exchange\in\setexchanges
    % =
    % \cbr{0,\ldots,n}\times\cbr{0,\ldots,m}
\)
let
\(
    \utilas = \utila\big(\valai,\valbj\big)
\)
and
\(
    \utilbs = \utilb\big(\valai,\valbj\big)
\)
% let $\utila:\setexchanges\rightarrow\RR$ and $\utilb:\setexchanges\rightarrow\RR$ denote
be the utilities for both agents,
where
\(
    \utila,\utilb:\RR^2\rightarrow\RR
\)
are functions mapping the value obtained by both competitors to their utilites.
In other words, $\utila,\utilb$ model how competitors care about either value 
derived from collaborating.

For example, a natural model is when competitors' utilies are linear both part parties'
valuations and are given by
\(
    \utilas
    =
    \valai-\compcoeff\valbj
\)
and
\(
    \utilbs
    =
    \valbj-\compcoeff\valai
\)
for $\compcoeff\in(0,1)$.
In this model, $\compcoeff$ quantifies how directly the agents are in competiton.
When $\compcoeff=1$ agents are in a zero-sum game whereas $\compcoeff=0$
signifies no competition. 

Let $\welfarefunc:\RR^2\rightarrow\RR$ be a selected welfare function used to quantify the quality of an exchange
$\exchange\in\setexchanges$ via $\welfarefunc\rbr{\utilas,\utilbs}$. 
Common examples of welfare functions include the egalitarian welfare~\citep{mas1995microeconomic} 
$\welfarefunc(x,y)=\min(x,y)$, the 
Nash welfare~\citep{nash1950bargaining} $W(x,y)=xy$, or the social welfare subject to individual rationality 
$\welfarefunc(x,y)=(x+y)\cdot 1_{\cbr{x,y\geq 0}}$. We are interested in the problem of designing  
algorithm which facilitates an exchange $\exchange\in\setexchanges$ to maximize the welfare.

If the valuation functions $\vala,\valb$ are deterministic, this is not an interesting problem as it simply reduces to 
the discrete optimization problem
\(
    \argmax_{
        \exchange\in\setexchanges
    }
    \welfarefunc\rbr{
        \utilas,\utilbs
    }
\). 
When $\vala$ and $\valb$ are random,
% However, agents may not know their value for a particular quantity of goods until after 
% receiving and using them? 
% For example, an agent may value data based on the accuracy of a model trained from it, which is 
% unknown until the data is used. In such cases what exchanges should a data sharing algorithm propose 
% to maximize the welfare?
% To model this, we assume that $\vala$ and $\valb$ are independent random variables, each drawn from known but not 
% necessarily identical distributions. 
% Under this model 
a natural approach would be to propose an
exchange that maximizes the expected welfare, i.e. to solve for 
% \begin{align*}
\(
    \argmax_{
        \exchange\in\setexchanges
    }
    \EE\sbr{
        \welfarefunc\rbr{
            \utilas,\utilbs
        }
    } .
\)
% \end{align*}

However, it turns out that if instead of proposing a one time exchange, we propose a sequence of increasing
exchanges (with both competitors reporting their observed values after each round), then we can obtain a 
higher expected welfare. This is demonstrated in the following example.

\begin{restatable}{example}{multi-vs-single-round-alg}
    \label{ex:multi-vs-single-round-alg}
    Let $\numdataa=\numdatab=1$, so $\setexchanges=\cbr{0,1}^2$. Suppose that each competitor's value for the other's good is 
    $\valaii{1},\valbjj{1}\drawniid\Bern(1/2)$ and assume $\valaii{0}=\valbjj{0}=0$.
    Let the utilities and welfare be given by
    \begin{align*}
        \utilas
        =
        \valai-\frac{1}{4}\valbj,
        \quad 
        \utilbs
        =
        \valbj-\frac{1}{4}\valai,
        \quad
        \text{and}
        \quad
        \welfarefunc(x,y)=\min(x,y) .
    \end{align*}  
    A direct computation shows that 
    \begin{align*}
        \rbr{1,1} 
        =
        \argmax_{
            \exchange\in\cbr{0,1}^2
        }
        \EE\sbr{
            \min\rbr{
                \utilas,\utilbs
            }
        } 
        \quad
        \text{and}
        \quad
        \EE\sbr{
            \min\rbr{
                \utilass{(1,1)},\utilbss{(1,1)}
            }
        } 
        =
        \frac{1}{16} .
    \end{align*}  
    Thus $(1,1)\in\setexchanges$ is the single exchange that maximizes the expected welfare.

    Now consider a multi-exchange algorithm that works as follows. The algorithm first gives one of 
    $\competitora$ or $\competitorb$ the other's resouce (without loss of generality assume 
    $\competitora$ receives $\competitorb$'s resource). 
    Competitor $\competitora$ is then asked to report their value of the received resource
    (either $0$ or $1$). If the $\competitora$ reports a $1$, then $\competitorb$ is 
    given $\competitora$'s resource. 
    However, if $\competitora$ reports a $0$, then $\competitorb$ is given nothing.
    Let $\exchange'=(i',j')\in\cbr{0,1}^2$
    denote the final exchange in this algorithm. A direct computation shows that 
    \begin{align*}
        \EE\sbr{
            \min\rbr{
                \utilass{s'},\utilbss{s'}
            }
        } 
        &=
        \EE\sbr{
            \min\rbr{
                \valaii{i'}
                -
                \frac{1}{4}
                \valbjj{j'}
                ,
                \valbjj{j'}
                -
                \frac{1}{4}
                \valaii{i'}
            }
        } 
        \\
        &=
        \frac{1}{2}
        \rbr{
            \frac{1}{2}
            \min\rbr{
                1-\frac{1}{4}
                ,
                1-\frac{1}{4}
            }
            +
            \frac{1}{2}
            \min\rbr{
                1-\frac{1}{4}
                ,
                -\frac{1}{4}
            }
        }
        +
        \frac{1}{2}\min(0,0)
        \\
        &=
        \frac{1}{8} .
    \end{align*}  
\end{restatable}  

Example~\ref{ex:multi-vs-single-round-alg} demonstrates that algorithms which can use competitors'
reported valuations to propose an increasing sequence of exchanges can obtain a strictly higher
expected welfare than a single exchange algorithm. Thus our goal will be to find the
optimal multi-exchange algorithm.

\textbf{Truthfulness assumptions.}
% Multi-exchange algorithms like the one given in example 1.1 rely on competitors truthfully reporting their value
% or their value being publicly observable/validated 
% for the data they have received. This is because in order to fairly determine future exchanges in both
% competitors best interests, they need to know the utilities of both agents. 
% To capture this difficulty, we first define the price of stability in this problem to be the 
% ratio of the expected welfare of the optimal multi-exchange algorithm to the expected welfare of the 
% optimal truthful mult-exchange algorithm. We then present a hardness result in Prop(\mycomment{fill later}),
% showing that the price of stability is unbounded. In other words, truthful algorithms cannot compete 
% with multi-exchange algorithms.
Multi-exchange algorithms, like Example~\ref{ex:multi-vs-single-round-alg}, rely on competitors truthfully 
reporting (or publicly validating) the value of the resources they receive. This is essential for 
fairly determining future exchanges based on both agents' utilities. To formalize this challenge, 
we will define the price of stability as the ratio between the expected welfare of the optimal 
multi-exchange algorithm and that of the optimal truthful one. 
In Proposition (\mycomment{fill later}), we show this ratio is unbounded, i.e. truthful algorithms cannot 
match the performance of unrestricted multi-exchange algorithms.


\mycomment{say other person's utility is a function of their valuation???}

\mycomment{add an unbounded PoS result}
\mycomment{does PoS even make sense cause welfare can be 0 or negative?}
\mycomment{abs difference unbounded}
% \mycomment{put in here something saying in order to create multi round }

\textbf{When participants can exit early.}


\subsection{Results and Technical Contributions.}
Our main contributions are formulating the problem of incentivizing competing agents to exchange resources
given uncertainty about their valuations for each others' resources and the design of an optimal 
algorithm to do so.

\textbf{Our novel setting.} 
\mycomment{edit this to be a mini prior work of why our problem setting is different and not captured
by existing work}
\mycomment{argue that existing models in other papers don't capture increasingly important usecases like
data sharing for ml}
\mycomment{while doing this, introduce our notation}
In Section~\ref{sec:model} we start by formalizing the interaction protocol
between competitors when exchanging goods, defining their valuations, utilities, and the welfare under
a particular exchange. 
Motivated by Example~\ref{ex:multi-vs-single-round-alg}, we define an exchange algorithm, to be a sequence
of increasing proposed exchanges, each based on reported valuations from the prior exchanges.
From this, we define an optimal exchange algorithm to be one which maximizes the
expected welfare. 
Finding an algorithm is our first goal.
We also define what it means for an exchange algorithm to be individually rational: in expectation 
competitors are no worse off participating than forgoing collaboration.
\mycomment{make sure we introduce utility, welfare here to, cause mentioned in following section}

\textbf{Reducing algorithm design to optimal stopping.} 

\emph{Challenges:}
A central challenge in designing an optimal exchange algorithm is proposing a sequence 
of optimal exchanges using only the information revealed by the previous exchange about 
future exchanges.
This is further complicated by the lack of assumptions we require of 
the valuation functions,
utilty functions, and welfare function. For example, 
% for $\arbitrarycompetitor\in\cbr{\competitora,\competitorb}$, 
the random variables 
\(
    \big\{
        \valai
    \big\}_{\numdatafora=0}^\numdataa
\)
could be continuous, discrete, or arbitrarily correlated
as is also the case with
\(
    \big\{
    % \{
        \valbj
    \big\}_{\numdataforb=0}^\numdatab
    % \big\}_{\numdataforb=0}^\numdatab
\).
Furthermore, the only assumption we make on these random variables (and the utilty and welfare funcitons 
$\utilas,\utilbs,\welfarefunc$) is boundedness.
Because of this, projecting the welfare of future exchanges is complicated.
Thus, it is difficult to design an algorithm which repeatedly proposes
a sequence of exactly optimal exchanges given the available information from the previous exchanges.
Competition adds another layer of complication as exhanges that are in one competitor's best interest may not 
be in the interest of the other's.

\emph{Techniques and insights:}
Our first insight is to reframe an algorithm as a policy for traversing the lattice 
\(
    \setexchanges
    =
    \cbr{0,\ldots,\numdataa}
    \times
    \cbr{0,\ldots,\numdatab}
\), 
which encodes the set of all valid exchanges,
starting at $(0,0)$ and moving only up/right at each step, using the information 
reported from previous exchanges.
Each such traversal corresponds to an increasing sequence of exchanges, representing that competitors
cannot unshare resources.
Our core insight is that we can use this reformulation to reduce the design of an optimal exchange
algorithm to solving an optimal stopping time problem over the partially ordered lattice $\setexchanges$,
leveraging techniques from~\citep{mandelbaum1981optimal}.
At a high level, the algorithm works as follows. At each time, determine which of the following three options would lead to 
a higher expected welfare: (1) terminating the exchange, (2) giving $\competitora$ an additional one of 
$\competitorb$'s resources and then ``proceeding optimally'', or (3) giving $\competitorb$ an additonal one of $\competitorb$'s
resources and then ``proceeding optimally''.
We compute the expected welfare of ``proceeding optimally'' after a given exchanged, using a dynamic programming 
type result from~\citep{mandelbaum1981optimal} which is a common approach in optimal stopping.
We show that under a reasonable assumption on the welfare function $\welfarefunc$,
the optimal algorithm we obtain is also automatically individually rational.

% \begin{enumerate}
%     \item A main challenge in designing an opt alg is that the space of all possible algorithms is extremely complicated
%     and that we do not put \textbf{any} restrictions on the valuation functions, for example they could be discrete support
%     or they could be continues and markovian, give example of trending upwards markov normal dist model? We also do not 
%     make any assumptions on the welfare functions, utility funcitons, other than normalization so that when no data is shared
%     welfares and utilities are 0. Competiton also inhererntly makes this complicated because exchanges that are in one agents' 
%     best interest may not be in the others'. 
    
%     \item
%     Our first technical insight is to reframe the design of an exchange algorithm to the design of an algorithm which traverses
%     the lattice (S= ..., moving up/right) which encodes the possible exchanges between agents. Our core insight is that 
%     we can use this reformulation to reduce the design of an optimal exchange algorithm to solving an optimal stopping time 
%     over the partially ordered lattice using techniques from~(). 

%     \item
%     To perform this reduction we consider the state space ... along with the filtration ... to capture the information know
%     to the algorithm designer after both agents report their valuations after an exchange. 
%     The result is an algorithm which achieves the optimal expected revenue among all algorithms. At a high level, this algorithm
%     works as follows. At each time t give only one of either a or b the others' next good and ask them to report their valuation
%     based on this the algorithm can decide which agent is 
    
%     \mycomment{mention dynamic programming}
% \end{enumerate}  


% In Section~\ref{sec:opt-non-ic-alg} we show how the problem of 
% designing an optimal exchange algorithm can be reduced to solving an optimal stopping time problem over 
% the partially ordered lattice $\NN^2_{\geq 0}$ using the framework of~\citet{mandelbaum1981optimal}. 
% % stopping time problem over the partially ordered set $\NN^2_{\geq 0}$ (with $(x_1,y_1)\leq (x_2,y_2) \iff$ 
% % $x_1\leq x_2$ and $y_1\leq y_2$).
% % More concretely, we leverage an extension of the classical optimal stopping 
% Because an exchange algorithm is an increasing sequence of exchanges (we use the standard partial ordering
% over $\NN^2$ to compare exchanges: $(i_1,j_1)\leq(i_2,j_2)\iff i_1\leq i_2$ and $j_1\leq j_2$), solving 
% for the optimal exchange algorithm is equivalent to determinining how to traverse $\NN^2_{\geq 0}$ in an
% increasing manner and determining when to stop traversing (i.e. when to settle on a final exchange).
% In addition to proving optimality, we show that for any choice of welfare function that respects individual
% rationality, our optimal exchange algorithm is also individually rational.
% As our definition of individual rationality is ex-ante, we supplement this with high probability guarantees
% that our optimal exchange algorithm is ex-post individually rational.

\mycomment{sep line}

\mycomment{put in IC motivating example here? or is it weird to have an example in the contributions section}
\textbf{An optimal incentive compatible algorithm.}
Our model thus far has assumed that once an agent signs on to participate in an algorithm, it will 
adhere to the sequence of exchanges proposed until the algorithm terminates. 
However, a competitor may be motivated to exit the exchange process early if she gets lucky with the value
she obtains from early on exchanges. In Section~\ref{sec:opt-ic-alg} we define incentive compatible
exchange algorithms to be those where neither competitor ever benefits from exiting the 
exchange algorithm early. We then give simple examples where the optimal incentive compatible exchange algorithm 
achieves a strictly lower expected welfare than the optimal exchange algorithm, demonstrating that 
incentive compatiblility is a property that cannot be obtained for free. We will then give examples showing the 
different challenges at play that complicate the design of an optimal incentive compatible algorithm. 

A priori,
the answer to whether an optimal incentive compatible algorithm even exists is not clear. We answer this question
in the affirmative, by constructing such an algorithm, providing novel adaptations of concepts from the theory 
of optimal stopping. 
We show that in addition to being incentive compatible,
this algorithm satisfies the same ex-ante and ex-post individual rationality guarantees as our first algorithm.
% One caveat is that 


% \paragraph{Truthfulness impossibility results}


\subsection{Related Work}

\begin{enumerate}
    \item \mycomment{Is there existing econ literature on the exchange of goods under competition?}
    \item Look at peter's? for info
    \item related-work-planning.md
\end{enumerate}  
\section{Model}
\label{sec:model}

% We state our model in two parts. First, we define our environment to capture how competitors value
% each others' resources, how competition influences their willingness to exchange, and how to quantify
% the quality of an exchange. In the second section, we give a precise definition of a resource 
% exchange algorithm under a particular environment and its potential properties
% (e.g. optimality, individual rationality, and incentive compatibility).

We now present our model. In Section~\ref{sec:modelenv} we define the environment, which describes how competitors value each other's 
resources, how competition affects their willingness to trade, and how to evaluate the quality of an exchange. 
In Section~\ref{sec:modelalgo} we formally define a resource exchange algorithm within this environment and outline its desirable properties, 
such as optimality, individual rationality, and incentive compatibility.

\subsection{Environment}
\label{sec:modelenv}

\paraheadwnogap{Agents, valuations, and utilities}
There are two competitors $\competitora$ and $\competitorb$, each possessing 
a collection of $\numdatab$ and $\numdataa$ goods respectively, that the other competitor values.
Let \(
    \cbr{
        \valai
    }_{\numdatafora=0}^\numdataa,
    \big\{
        \valbj
    \big\}_{\numdataforb=0}^\numdatab
\) be a sequence of random variables, drawn independently from distributions $\priora, \priorb$, 
where $\valai$ denotes how much player $a$ values $i$ goods from $b$ and likewise for $\valbj$.
Although we assume the sequences are independent of each other, i.e. $\valai \perp \valbj$ for all $i,j$, 
each sequence could be correlated, $\valaii{i_1} \not\perp \valaii{i_2}$ for all $i_1,i_2$.
For convenience, we will take $\valaii{0} = \valbjj{0} = 0$.
Finally, while the distributions
$\priora,\priorb$,
%  $\prior$ 
are public information, neither player knows the realizations of their valuations
\(
    \cbr{
        \valai
    }_{\numdatafora=1}^\numdataa,
    \big\{
        \valbj
    \big\}_{\numdataforb=1}^\numdatab
\)
until the corresponding number of goods have been exchanged.

% \[
% \vala:\cbr{0,\ldots,\numdataa}\rightarrow\RR
% \quad
% \text{and} 
% \quad
% \valb:\cbr{0,\ldots,\numdatab}\rightarrow\RR 
% \]
% be two independent random valuation functions, capturing how much each competitor values a particular number of 
% resources received from the other.
% Equivalently, we can view these two functions as two sequence of random variables
% \(
%     \cbr{
%         \valai
%     }_{\numdatafora=0}^\numdataa,
%     \big\{
%         \valbj
%     \big\}_{\numdataforb=0}^\numdatab
% \)
% that are independent from 
% each other. For convenience, assume that 
% \(
%     \valaii{0}
%     =
%     \valbjj{0}
%     =
%     0
% \), 
% i.e. competitors derive no value from each others' resources when none are shared.

% Suppose that there are two competitors $\competitora$ and $\competitorb$, each possessing 
% a collection of $\numdatab$ and $\numdataa$ different resources
% respectively, that the other competitor values. Let 
% \[
% \vala:\cbr{0,\ldots,\numdataa}\rightarrow\RR
% \quad
% \text{and} 
% \quad
% \valb:\cbr{0,\ldots,\numdatab}\rightarrow\RR 
% \]
% be two independent random valuation functions, capturing how much each competitor values a particular number of 
% resources received from the other.
% Equivalently, we can view these two functions as two sequence of random variables
% \(
%     \cbr{
%         \valai
%     }_{\numdatafora=0}^\numdataa,
%     \big\{
%         \valbj
%     \big\}_{\numdataforb=0}^\numdatab
% \)
% that are independent from 
% each other. For convenience, assume that 
% \(
%     \valaii{0}
%     =
%     \valbjj{0}
%     =
%     0
% \), 
% i.e. competitors derive no value from each others' resources when none are shared.

Let $\setexchanges=\cbr{0,\ldots,\numdataa}\times\cbr{0,\ldots,\numdatab}$
denote the set of possible exchanges, where
\(
    \rbr{\numdatafora,\numdataforb}\in\setexchanges
\)
represents $\competitora$ receiving $\numdatafora$ of $\competitorb$'s resources and 
$\competitorb$ receiving $\numdataforb$ of $\competitora$'s resources.
Because $\competitora$ and $\competitorb$ are in competition, we assume their utilities under 
an exchange $\exchange = (i,j)\in\setexchanges$ are functions of both their own value and the value the other agent derives from the exchange.
Concretely,
let
\(
% \(
    \utilas = \utila\big(\valai,\valbj\big)
% \)
%\quad
\text{ and }
%\quad
% \(
    \utilbs = \utilb\big(\valai,\valbj\big)
% \)
\)
be the utilities for both agents under an exchange $\exchange\in\setexchanges$
where
\(
    \utila,\utilb:\RR^2\rightarrow\RR
\)
are bounded functions. 
For notational convenience,
we will denote $\utilass{0} = \utilass{(0,0)}=0$ and
$\utilbss{0} = \utilbss{(0,0)}=0$.
% we 
% assume
% \(
%     \utilass{0}
%     =
%     \utilbss{0}
%     =0
% \)
% and let the subscript $0$ refer to $(0,0)$.

This interaction between the two agents may be non-zero-sum (i.e. $\utilas \neq -\utilbs$).
Moreover, utility gains for the agents may be asymmetric;
for example, we could have
% \[
\(
    \utila(x,y)=x-\frac{1}{10}y
\)
% \quad
% \text{and}
and
% \quad
\(
    \utilb(x,y)=y-\frac{1}{3}x
\),
% \]
reflecting a scenario in which $\competitora$ is more sensitive to aiding $\competitorb$ than
vice versa.
Moreover, while the canonical use case of interest involves negative externalities due to competition (i.e. $\utila$ decreases in $\valbj$), 
our model and methods do not preclude positive externalities (i.e. $\utila$ may increase with $\valbj$).




\parahead{Welfare}
% To quantify social benefit due to an exchange, 
To quantify the quality of an exchange,
we introduce a bounded welfare function
\(
    \welfarefunc:\RR^2\rightarrow\RR
\).
We then define the welfare of an exchange $\exchange\in\setexchanges$ to be the random variable
\(
    \rewards
    =
    \welfarefunc\rbr{
        \utilas,\utilbs
    }
\).
For convenience, we will assume
\(
    \welfarefunc\rbr{0,0}=0
\).
Common examples of welfare functions  include the egalitarian welfare~\citep{mas1995microeconomic}
$\welfarefunc(x,y)=\min(x,y)$, the 
Nash welfare~\citep{nash1950bargaining} $W(x,y)=xy$, or the social welfare subject to individual rationality 
$\welfarefunc(x,y)=(x+y) 1_{\cbr{x,y\geq 0}}$.




In the definitions above, and going forward, we assume all functions are measurable unless stated otherwise.
Using the components above we define an exchange environment.

\begin{restatable}{definition}{def:exchange-environment}
    An exchange environment is a tuple
    \(
        \rbr{
            \cbr{
                \valai
            }_{\numdatafora=0}^\numdataa,
            \big\{
                \valbj
            \big\}_{\numdataforb=0}^\numdatab,
            \utila,
            \utilb,
            \welfarefunc
        }
    \),
    where
    \(
        \cbr{
            \valai
        }_{\numdatafora=0}^\numdataa
    \)
    and
    \(
        \big\{
            \valbj
        \big\}_{\numdataforb=0}^\numdatab,
    \)
    are independent sequences drawn from priors $\priora, \priorb$,
    \(
        \utila,\utilb,\welfarefunc : \RR^2\rightarrow\RR
    \)
    are bounded functions, 
    and all variables are normalized to zero as described above.
    % and 
    % \(
    %     \valaii{0},
    %     % =
    %     \valbjj{0},
    %     % =
    %     \utilass{0},
    %     % \utila(0,0),
    %     % =
    %     \utilbss{0},
    %     % \utilb(0,0),
    %     % =
    %     \welfarefunc(0,0),
    %     =0
    % \)
\end{restatable}  

\subsection{Algorithm Definiton and Properties}
\label{sec:modelalgo}

% \parahead{fe}
\parahead{Algorithm description}
Before formally defining an exchange algorithm, we first describe how an algorithm interacts with the agents.
First the algorithm will be published and agents will decide whether they wish to participate.
% in the data sharing process that the algorithm will compute. 
If both agents opt in, then the algorithm
will propose a first exchange 
\(
    \algstatett{1}
    =
    \rbr{
        \numdataforatt{1},
        \numdataforbtt{1}
    }\in\setexchanges 
\)
where $\competitora$ will receive $\numdataforatt{1}$ of $\competitorb$'s resources and 
$\competitorb$ will receive $\numdataforbtt{1}$ of $\competitora$'s resources.
After this exchange both agents truthfully report their observed values for 
each others' resources, 
\(
    \valaii{
        \numdataforatt{1}
    },
    \valbjj{
        \numdataforbtt{1}
    }
\).
Using this feedback, the algorithm will determine whether to terminate or propose another 
exchange 
\(
    \algstatett{2}
    =\rbr{
        \numdataforatt{2},
        \numdataforbtt{2}
    }\in\setexchanges 
\).
The process repeats until the algorithm terminates with a final exchange of 
\(
    \algstatett{
        \algendtime
    }
    =\rbr{
        \numdataforatt{\algendtime},
        \numdataforbtt{\algendtime}
    }\in\setexchanges 
\).
We require that exchanges are increasing\footnote{%
When we write $\rbr{x_1,y_1}\leq\rbr{x_2,y_2}$, we mean
\(x_1\leq x_2\), and \(y_1\leq y_2\).
When we write $\rbr{x_1,y_1}<\rbr{x_2,y_2}$, then at least one of the two inequalities holds strictly.
}, meaning $\algstatett{t} > \algstatett{t-1}$,
to reflect that resources cannot be unshared and that on each round the exchange algorithm does not terminate, at least one agent must receive an additional good.
We now formally define an exchange algorithm.

\begin{restatable}[Algorithm]{definition}{exchange-alg}
    \label{def:exchange-alg}
    \emph{A resource exchange algorithm is a tuple
    \(
        \rbr{
            \cbr{
                \algstatemapt
            % }_{\timeind\geq 1}, 
            }_{\timeind\geq 0}, 
            \cbr{
                \algstopmapt
            }_{\timeind\geq 0}
        }
    \)
    % of sequence of measurable functions
    where
    % The first sequence
    \begin{align*}
        \algstatemapt\rbr{
            \cbr{
                \valaii{
                    k
                }
            }_{k=0}^{
                % \numdataforatt{t-1}
                \numdataforatt{t}
            },
            \cbr{
                \valbjj{
                    k
                }
            }_{k=0}^{
                % \numdataforbtt{t-1}
                \numdataforbtt{t}
            }
        }
        =
        \rbr{
            \numdataforatt{\timeind+1},
            \numdataforbtt{\timeind+1}
        }
        =
        \algstatett{\timeind+1}
        \in\NN^2_{\geq 0}
        \quad
        \text{s.t.}
        \quad
        \algstatet
        <
        \algstatett{t+1}
        \quad
        \rbr{
            \text{here }
            \algstatett{0}
            =
            \rbr{
                \numdataforatt{0},\numdataforbtt{0}
            }
            =
            (0,0)
        }
    \end{align*}     
    % is a sequence of 
    is a sequence of 
    mappings, each from values realized for previous exchanges, to a larger exchange,
    and
    % \(
    %     \algstopmapt
    %     :
    %     % \RR^{2\timeind}
    %     \RR^{\numdataforatt{t}+1}
    %     \times
    %     \RR^{\numdataforbtt{t}+1}
    %     \rightarrow \cbr{0,1}
    % \)
    \[
        \algstopmapt
        \rbr{
            \cbr{
                \valaii{
                    k
                }
            }_{k=0}^{
                \numdataforatt{t}
            },
            \cbr{
                \valbjj{
                    k
                }
            }_{k=0}^{
                \numdataforbtt{t}
            }
        }
        \in\cbr{0,1}
    \]
    is a 
    sequence of decisions on whether to terminate proposing exchanges.
    We define
    % \[
    %     \algendtime
    %     =\inf\cbr{
    %         t\in\NN_{\geq 0} :
    %         \delta_{\timeind}\rbr{
    %             \cbr{
    %                 \valaii{k},
    %                 \valbjj{k}
    %             }_{k=0}^{\timeind-1} 
    %         }=0
    %     }
    % \] 
    \[
        \algendtime
        =\inf\cbr{
            t\in\NN_{\geq 0} :
            \delta_{\timeind}\rbr{
                \cbr{
                    \valaii{k}
                }_{k=0}^{\numdataforatt{t}}
                ,
                \cbr{
                    \valbjj{k}
                }_{k=0}^{\numdataforbtt{t}}
            }=0
        }
    \] 
    to be the time exchanging is discontinued, 
    % and require $\algendtime$ is finite. 
    and require
    \(
        \algstatett{
            \algendtime
        }
        \leq (\numdataa,\numdatab)
    \), i.e. the final exchange does not exceed the resources competitors possess.
    Let $\classalgs$ denote the set of all resource exchange algorithms. }
\end{restatable}

% Going forward, expectations $\EE$ will be with respect to the distribution of realized valuations $\vala_{i_t}, \valb_{j_t}$, welfares $\{\rewards\}_s$, and 
% decisions $\{\algstatett{t}\}_t$ due to the interaction between an algorithm $\alg$ and environment under the priors $\priora, \priorb$.
% If we wish to emphasize the dependence of $\EE$ under an algorithm $\alg$, we will write $\EE_\alg$.
% For a specified environment, the goal of the algorithm designer is to propose an algorithm 
% $\alg\in\classalgs$ to maximize the expected welfare of the final exchange,
% \(
%     \EE\sbr{
%         \rewardss{\algstatett{
%             \algendtime}
%         }
%     }
% \).

\parahead{Algorithm properties}
We now introduce desirable properties than an exchange algorithm can possess. The first such property is optimality. 
We say that an algorithm is optimal if there is no other algorithm with a higher expected welfare.

\begin{restatable}{definition}{welfare-of-alg-and-opt-alg}
    For an algorithm
    \(
        \alg
        =
        \rbr{
            \cbr{
                \algstatemapt
            }_{\timeind\geq 0}, 
            \cbr{
                \algstopmapt
            }_{\timeind\geq 0}
        }
        \in\classalgs
    \)
    we define its expected welfare to be that of the final exchange,
    \(
        \EE\sbr{
            \rewardss{
                \algstatett{\algendtime}
            }
        }
    \).
    We say that an algorithm $\alg$ is optimal if
    \begin{align*}
        % \EE_\alg\sbr{
        \EE\sbr{
            \rewardss{
                \algstatett{\algendtime}
            }
        }
        =
        \sup_{
            \alg'\in\classalgs
        }
        % \EE_{\alg'}\sbr{
        \EE\sbr{
            \rewardss{
                \algstatett{\algendtime'}'
            }
        } .
    \end{align*} 
\end{restatable}  

% In otherwords, an algorithm is optimal if there is no other algorithm with a higher expected 
% welfare. 
% The goal of a benevolent algorithm designer
% is to find an optimal $\alg\in\classalgs$ 
% to maximize the benefit of competitors collaborating
% by maximizing the expected welfare.
% Ideally the algorithm designer would like to propose an optimal exchange algorithm.
% % to maximize the benefit of competitors collaborating
% by maximizing the expected welfare.
% However, it is a priori not obvious if or when there exists an optimal resource exchange algorithm, 
% and if it does, how to go about finding it. 

Our second property for an algorithm is individual rationality and is defined as follows.
% Additionally, we would like to ensure that under the proposed algorithm is individually rational for both agents
% as defined below.

\begin{restatable}{definition}{ir-exchange-alg}
    An algorithm $\alg\in\classalgs$ is individually rational (IR) if 
    $\EE[\utilass{ \algstatett{\algendtime}
            }] \geq 0$
    and
    $\EE[\utilbss{ \algstatett{\algendtime}
            }] \geq 0$.
    % \begin{align*}
    % \(
    %     0
    %     \leq
    %     \EE\sbr{
    %         \utilass{ \algstatett{\algendtime}
    %         }
    %     },
    %     \EE\sbr{
    %         \utilbss{
    %             \algstatett{\algendtime}
    %         }
    %     }
    % \).
\end{restatable}  

%\kkcomment{State $\EE_A$ and that we wil drop goig forward.}

Thus far, we have assumed that once an algorithm is published competitors either opt in or out of the 
algorithm for the entire duration. However, in some real-world use cases, agents may exit the collaboration early. This is a potential issue when one competitor 
benefits significantly from collaboration early on and determines it would only serve their 
competitor by continuing to participate. An algorithm which has the property that competitors are 
never incentivized to leave the collaboration early (prior to the algorithm terminating) is
called incentive compatible (IC). This is made precise in the following definition.

\begin{restatable}{definition}{ic-exchange-alg}
    % We say that an algorithm 
    \label{def:ic-exchange-alg}
    Let
    \(
        \alg
        =
        \rbr{
            \cbr{
                \algstatemapt
            }_{\timeind\geq 0}, 
            \cbr{
                \algstopmapt
            }_{\timeind\geq 0}
        }
        \in\classalgs
    \) be an algorithm, 
    and let
    % is incentive compatible if for any alternate sequence of stopping decisions
    \(
        \cbr{
            \altalgstopmapt
        }_{\timeind\geq 0}
    \)
    be an alternate sequence of stopping decisions, i.e.
    \[
        \altalgstopmapt
        \rbr{
            \cbr{
                \valaii{
                    k
                }
            }_{k=0}^{
                \numdataforatt{t}
            },
            \cbr{
                \valbjj{
                    k
                }
            }_{k=0}^{
                \numdataforbtt{t}
            }
        }
        \in\cbr{0,1}
        \quad
        \text{and}
        \quad
        \altalgendtime
        =\inf\cbr{
            t\in\NN_{\geq 0} :
            \altalgstopmaptt{\timeind}\rbr{
                \cbr{
                    \valaii{k}
                }_{k=0}^{\numdataforatt{t}}
                ,
                \cbr{
                    \valbjj{k}
                }_{k=0}^{\numdataforbtt{t}}
            }=0
        } .
    \]
    We say $\alg$ is incentive compatible if for any choice of 
    \(
        \cbr{
            \altalgstopmapt
        }_{\timeind\geq 0}
    \)
    we have
    \[
        \EE\sbr{
            \utilass{
                \algstatett{
                    \algendtime\wedge\altalgendtime
                }
            }
        }
        \leq
        \EE\sbr{
            \utilass{
                \algstatett{
                    \algendtime
                }
            }
        }
        \quad
        \text{and}
        \quad
        \EE\sbr{
            \utilbss{
                \algstatett{
                    \algendtime\wedge\altalgendtime
                }
            }
        }
        \leq
        \EE\sbr{
            \utilbss{
                \algstatett{
                    \algendtime
                }
            }
        } .
    \] 
    Let $\classicalgs$ denote the set of IC exchange algorithms.
    We say that $\alg\in\classicalgs$ is IC optimal if
    \[
        \EE\sbr{
            \rewardss{
                \algstatett{\algendtime}
            }
        }
        =
        \sup_{
            \altalg\in\classicalgs
        }
        \EE\sbr{
            \rewardss{
                \algstatett{\algendtime'}'
            }
        } .
    \]
\end{restatable}  

In the definition above, any 
\(
    \cbr{
        \altalgstopmapt
    }_{\timeind\geq 0}
\)
can be viewed as a sequence of decisions a competitor makes on whether to exit the collaboration or
continue participating.
With this interpretation,
% In the definition above, 
\(
    \algendtime\wedge\altalgendtime
\)
represents the time collaboration stops, either because the algorithm has terminated on round $T$, or one of 
the participants has decided to exit early at round $T'$.
The definition above says that incentive compatible algorithms are those for which neither competitor
can increase their expected utility by trying to strategically exit the collaboration early.
As we will see, incentive compatibility is a strong requirement for an algorithm and significantly
complicates its design.
\section{An Optimal IR Resource Exchange Algorithm}
\label{sec:opt-non-ic-alg}

% Our first algorithm design goal will be to determine when there exists
% an optimal algorithm and to construct one in this case.
% \mycomment{I don't think the term policy is mentioned anywhere in durrett}
Our first goal is to determine whether an optimal algorithm exists, and to construct one when it does.
To do so,
we will draw on tools from martingale theory to design our exchange algorithm and, accordingly, adopt some of the 
terminology from this literature. In this context, the term \emph{policy} is used as a generalization of stopping 
times to partially ordered sets~\citep{mandelbaum1981optimal,williams1991probability}. We focus on a restricted class 
of such policies, which we term \emph{exchange policies}. These policies satisfy additional structural properties that 
enable the design of exchange algorithms (Definition~\ref{def:exchange-alg}). In Section~\ref{sec:exchangealgotopolicy}, 
we will first show how an exchange algorithm can be reformulated as an exchange policy.
Then, in Section~\ref{subsec:reduction-to-optimal-stopping}, we will show how we can use optimal stopping over the lattice to derive an optimal exchange algorithm.

\subsection{Reformulating an Exchange Algorithm as an Exchange Policy}
\label{sec:exchangealgotopolicy}

%\kkcomment{Can we not call this an 'exchange policy'? }




% We will use tools from martingale theory to design an exchange algorithm, and as such, borrow some verbiage from this literature.
% In this literature, the term 'policy' is used as a generalization of stopping times for partially ordered sets~\citep{mandelbaum1981optimal,williams1991probability,durrett2019probability}.
% We will consider a subset of such policies, which we call `exchange policies', which have additional properties and can be used to design exchange algorithms (Definition~\ref{def:exchange-alg}).
% We will now define these terms formally.

Suppose that an exchange 
\(
    \exchange
    =
    \rbr{
        \numdatafora,
        \numdataforb
    }
\)
has already occurred. In this case $\competitora$ has observed their value for the first 
$\numdatafora$ of $\competitorb$'s resources and 
$\competitorb$ has observed their value for the first $\numdataforb$ of $\competitora$'s resources.
To represent the information
available to an algorithm after $\competitora$ and $\competitorb$ report their values for the 
exchange 
\(
    \exchange=\rbr{
        \numdatafora, \numdataforb
    }\in\setexchanges
\), we define the $\sigma$-algebra
\begin{align*}
    \statefilts
    =
    \sigma\rbr{
        \cbr{
            \valaii{k}
        }_{k=0}^\numdatafora,
        \cbr{
            \valbjj{k}
        }_{k=0}^\numdataforb
    }
    \numberthis
    \label{eq:family-sig-alg}
\end{align*}  
and consider the family
\(
    \cbr{
        \statefilts
    }_{
        \exchange\in\setexchanges
    }
\).
Notice that for two exchanges
\(
    \exchangett{1},\exchangett{2}
    \in\setexchanges
\)
we have 
% \begin{align*}
\(
    \exchangett{1}\leq\exchangett{2}
    \iff 
    \statefiltss{
        \exchangett{1}
    }
    \subseteq
    \statefiltss{
        \exchangett{2}
    }
\).
% \end{align*}  
This reflects that larger exchanges contain more information about the agents' valuation functions.
Additionally, 
\(
    \statefiltss{0}
    =
    \cbr{\varnothing,\Omega}
\)
which captures that no information besides the distributions $\priora,\priorb$ of the valuation functions is known prior to 
$\competitora$ and $\competitorb$ exchanging any data.

From Defintion~\ref{def:exchange-alg}, we see that for any algorithm $\alg\in\classalgs$, there is a correspondence
between $\alg$ and the sequence of $\setexchanges$-valued random variables 
\(
    \rbr{
        \algstatett{0},
        \ldots
        \algstatett{\algendtime}
    }
\)
it induces. Moreover, it is the case that 
\(
    \algstatett{\timeind+1}
\)
is
\(
    \statefiltss{
        \algstatet
    }
\)-measurable, reflecting that at round $\timeind$, the choice of the next exchange
only depends on information revealed at the time of the current exchange.
Intuitively, an algorithm can be thought of as a way of traversing $\setexchanges\subset\NN^2$, starting at $0$, and moving up or right
each round,
and deciding when to stop.
We will shortly see that an algorithm is a special case of 
% what~\citet{mandelbaum1981optimal} defines to be a policy.
a policy, which is the generalization of a stopping time to a partially ordered set.
\begin{restatable}{definition}{policy}
    \label{def:policy} 
    \textup{
        (\citet{mandelbaum1981optimal})
    }
    Let $\osasstates$ be a countable partially ordered set. For $r,s\in S$ denote
    \(
        \sbr{r,\osasstate}
        =
        \cbr{
            u\in \osasstates: r\leq u\leq \osasstate
        } 
    \).
     Let 
    \(
        \osasdirectsuccessorss{
            \osasstate
        }
        =
        \cbr{
            u> \osasstate: \sbr{\osasstate,u}=\cbr{\osasstate,u}
        } 
    \) be the set of direct successors.
    Assume that $\sbr{r,\osasstate}$ is finite and there is a unique minimal element $0\in \osasstates$. 
    Let 
    \(
        \cbr{\statefilts}_{
            \osasstate\in \osasstates
        }
    \)
    be a family of $\sigma$-algebras where $\Fcal_r\subseteq\Fcal_\osasstate$ for $r\leq \osasstate$. A strategy is 
    a sequence of $\osasstates$-valued random variables $\cbr{\sigma_\timeind}_{\timeind\geq 0}$ such that 
    \(
        \sigma_\timeind \leq \sigma_{\timeind+1}
    \)
    and
    \(
        % \quad
        \sigma_{\timeind+1} \text{ is }
        \Fcal_{\sigma_\timeind}\text{-measurable}
    \).
    A policy $\pi$ is a pair
    \(
        \rbr{
            \cbr{
                \osasstrategyt
            }_{\timeind\geq 0},
            \osasstoptime
        }
    \)
    where $\osasstoptime$ is a stopping time with respect to the filtration
    \(
        \cbr{
            \timefiltt
        }_{\timeind\geq 0}
        =
        \cbr{
            \statefilt_{
                \strategyt
            }
        }_{\timeind\geq 0}
    \).
    Denote $\osaspolicies$ as the class of policies such that $\osasstrategytt{0}=0$.
\end{restatable}  

To limit the number of subscripts, we will occassionally write
$\stopstateforpolicy{\apolicy}$ as shorthand for 
$\strategytt{\stoptime}$.
In~\citet{mandelbaum1981optimal}, each state is associated with a reward $\osasrewards$, and the expected 
reward of a policy 
\(
    \rbr{
        \cbr{
            \sigma_t
        }_{t\geq 0},
        \tau
    }
\)
is defined as
\(
    \EE\sbr{
        \osasrewardss{
            \osasstrategytt{
                \osasstoptime
            }
        }
    }
\).
Because of this, when specifying a policy 
\(
    \rbr{
        \cbr{
            \osasstrategyt
        }_{\timeind\geq 0},
        \osasstoptime 
    }
    % \in
    % \osaspolicies
\),
we only need to worry about defining $\strategyt$ for $\timeind\leq\stoptime$.
In our setting, an algorithm can be viewed as a policy that satisfies additional structure. We refer to such policies as \emph{exchange policies} and define them formally below.


% For our problem, an algorithm can be viewed as a policy with two additional requirements, taking $S=\setexchanges$
% (equiped with the standard partially ordering over $\NN^2$), 
% the family of $\sigma$-algebras as given in~\eqref{eq:family-sig-alg},
% and the reward $\osasrewards$ associated to $\exchange\in\setexchanges$ to be $\rewards$. 
% % Thus we have the following equivalent definition of a resource exchange algorithm.
% This is formalized in the following definition.
% \kkcomment{write more clearly. this sounds like there are 3 additional requirements.}


\begin{restatable}{definition}{exchange-policy}
    \label{def:exchange-policy} 
    A resource exchange policy
    % $\alg$
    is a policy 
    \(
        \osaspolicy=
        \rbr{
            \cbr{
                \strategyt
            }_{\timeind\geq 0},
            \stoptime 
        }
        \in
        \osaspolicies
    \),
    as specified in Definition~\ref{def:policy}, taking
    \(
        \osasstates=\setexchanges
    \)
    (which has the unique minimal element $(0,0)$)
    equiped with the family of $\sigma$-algebras 
    \(
        \cbr{
            \statefilts
        }_{\exchange\in\setexchanges}
    \)
    given in~\eqref{eq:family-sig-alg}. We additionally require that
    \(
        \strategyt < \strategytt{\timeind+1}
    \)
    for
    \(
        \strategyt < \rbr{
            \numdataa,\numdatab
        },
    \)
    and that $\stoptime$ is finite.
    Denote $\setexchangepolicies$ as the class of exchange policies.
\end{restatable}  

% Shortly, we will see the benefit of working with Definition~\ref{def:exchange-policy}
% instead of Definition~\ref{def:exchange-alg}.

The 
following result establishes an equivalence between exchange algorithms (Definition~\ref{def:exchange-alg}) and exchange policies (Definition~\ref{def:exchange-policy}).
As we will see, this will enable us to leverage tools from optimal stopping to design a welfare-optimal exchange algorithm.

% , shows that we 
% can move back and forth between talking about exchange algorithms and exchange policies.
% Concretely, given an exchange algorithm we can find an equivalent exchange policy and 
% vice versa.
% given an exchange policy we can find an equivalent exchange algorithm.

\begin{restatable}{theorem}{equiv-alg-defs}
    \label{thm:equiv-alg-defs}
    Fix an exchange algorithm
    \(
        \rbr{
            \cbr{
                \algstatemapt
            }_{\timeind\geq 0}, 
            \cbr{
                \algstopmapt
            }_{\timeind\geq 0}
        }
        \in
        \classalgs
    \).
    % as defined in Definition~\ref{def:exchange-alg}.
    Then, we can construct an exchange policy
    \(
        \rbr{
            \cbr{
                \strategyt
            }_{\timeind\geq 0},
            \stoptime 
        }
        \in
        \setexchangepolicies
    \) 
    such that 
    % satisfying Definition~\ref{def:exchange-alg-as-policy} such that
    % \begin{align*}
    \(
        \EE\sbr{
            \rewardss{
                \algstatett{\algendtime}
            }
        }
        =
        \EE\sbr{
            \rewardss{
                \strategytt{\stoptime}
            }
        }
    \).
    % \end{align*}
    Conversely, if given an exchange policy
    % as defined in Definition~\ref{def:exchange-alg-as-policy}.
    \(
        \rbr{
            \cbr{
                \strategyt
            }_{\timeind\geq 0},
            \stoptime 
        }
        \in
        \setexchangepolicies
    \) 
    we can construct an exchange algorithm
    \(
        \rbr{
            \cbr{
                \algstatemapt
            }_{\timeind\geq 0}, 
            \cbr{
                \algstopmapt
            }_{\timeind\geq 0}
        }
        \in
        \classalgs
    \)
    % satisfying Definition~\ref{def:exchange-alg} 
    such that 
    \(
        \EE\sbr{
            \rewardss{
                \strategytt{\stoptime}
            }
        }
        =
        \EE\sbr{
            \rewardss{
                \algstatett{\algendtime}
            }
        }
    \).
    % In other words, exchange algorithms are equivalent to exchange policies.
\end{restatable}  

\begin{proof}
    We start by constructing an exchange policy  
    given an algorithm
    \(
        \rbr{
            \cbr{
                \algstatemapt
            }_{\timeind\geq 0}, 
            \cbr{
                \algstopmapt
            }_{\timeind\geq 0}
        }
    \).
    Define 
    % the sequence
    \(
        \rbr{
            \cbr{
                \strategyt
            }_{\timeind\geq 0},
            \stoptime
        }
    \) 
    via
    \(
        \strategytt{0} = \zerotuple
    \)
    and
    \begin{align*}
    % \(
        \qquad
        \strategytt{\timeind+1}
        =
        \begin{cases}
            \algstatemapt\rbr{
                \cbr{
                    \valaii{k}
                }_{k=0}^{\numdataforatt{t}}
                ,
                \cbr{
                    \valbjj{k}
                }_{k=0}^{\numdataforbtt{t}}
            }
            & 
            \text{ if } \timeind < \algendtime
            \\
            \rbr{
                \numdataa,\numdatab
            }
            &
            \text{ if } \timeind \geq \algendtime
        \end{cases}  
    \end{align*}  
    along with setting
    \(
    % \begin{align*}
        \stoptime
        =
        \algendtime
        % =
        % \inf\cbr{
        %     t\in\NN_{\geq 0} :
        %     \delta_{\timeind}\rbr{
        %         \cbr{
        %             \valaii{k}
        %         }_{k=0}^{\numdataforatt{t}}
        %         ,
        %         \cbr{
        %             \valbjj{k}
        %         }_{k=0}^{\numdataforbtt{t}}
        %     }=0
        % }
    \).

    % \end{align*}  
    We first argue
    \(
        \rbr{
            \cbr{
                \strategyt
            }_{\timeind\geq 0},
            \stoptime
        }
    \) 
    is a well defined policy for 
    \(
        \osasstates=\setexchanges
    \)
    with the family of $\sigma$-algebras 
    \(
        \cbr{
            \statefilts
        }_{\exchange\in\setexchanges}
    \)
    given in~\eqref{eq:family-sig-alg}.
    After we show 
    \(
        \rbr{
            \cbr{
                \strategyt
            }_{\timeind\geq 0},
            \stoptime
        }
    \) 
    satisfies our additional two requirements to be an exchange policy

    \begin{enumerate}
        
    \item
        ($\strategyt\in\setexchanges$): 
        By definition
        \(
            \strategytt{0}=\zerotuple\in\setexchanges
        \).
        For $\timeind\geq 0$, we want to show
        \(
            \strategytt{\timeind+1}\in\setexchanges
        \).
        If $\timeind\geq \algendtime$ then automatically
        \(
            \strategytt{\timeind+1}
            =
            \rbr{
                \numdataa,\numdatab
            }\in\setexchanges
        \).
        If $\timeind < \algendtime$, then $\timeind+1\leq \algendtime$ and
        \[
            \strategytt{\timeind+1}
            =
            \algstatemapt\rbr{
                \cbr{
                    \valaii{k}
                }_{k=0}^{\numdataforatt{t}}
                ,
                \cbr{
                    \valbjj{k}
                }_{k=0}^{\numdataforbtt{t}}
            }
            =
            \algstatett{\timeind+1} .
            % \rbr{
            %     \numdataforatt{\timeind+1},
            %     \numdataforbtt{\timeind+1}
            % } .
        \]
        From Definition~\ref{def:exchange-alg}, we have
        \(
            \algstatett{k}
            <
            \algstatett{k+1}
            % \rbr{
            %     \numdataforatt{k},
            %     \numdataforbtt{k}
            % }
            % <
            % \rbr{
            %     \numdataforatt{k+1},
            %     \numdataforbtt{k+1}
            % }
        \)
        and 
        \(
            % \rbr{
            %     \numdataforatt{\algendtime},
            %     \numdataforbtt{\algendtime}
            % }
            \algstatett{\algendtime}
            \leq
            \rbr{
                \numdataa,\numdatab
            }
        \). 
        Therefore, 
        \(
            \strategytt{\timeind+1}
            =
            \algstatett{\timeind+1}
            \leq
            \algstatett{\algendtime}
            \leq
            \rbr{
                \numdataa,\numdatab
            }
        \)
        and so 
        \(
            \strategytt{\timeind+1}
            \in\setexchanges
        \).

        \item
        ($\strategyt\leq\strategytt{\timeind+1}$):
        For $\timeind\geq \algendtime$, $\strategytt{\timeind+1}$ is constant so this
        holds automatically. For $\timeind<\algendtime$,
        Definition~\ref{def:exchange-alg} requires
        \(
            \strategyt
            =
            \algstatet
            % \rbr{
            %     \numdataforat,
            %     \numdataforbt
            % }
            <
            % \rbr{
            %     \numdataforatt{\timeind+1},
            %     \numdataforbtt{\timeind+1}
            % }
            \algstatett{\timeind+1}
            =
            \strategytt{\timeind+1}
        \). 
    
        \item 
        ($\strategytt{\timeind+1}$
        is
        \(
            \statefiltss{
                \strategyt
            }
        \)-measurable): 
        Notice that we can rewrite $\strategytt{\timeind+1}$ as
        \[
            \strategytt{\timeind+1}
            =
                1_{
                    \cbr{
                        \timeind < \algendtime
                    }
                }
                % \cdot
                \algstatemapt\rbr{
                    \cbr{
                        \valaii{k}
                    }_{k=0}^{\numdataforatt{t}}
                    ,
                    \cbr{
                        \valbjj{k}
                    }_{k=0}^{\numdataforbtt{t}}
                }
                +
                1_{
                    \cbr{
                        \timeind \geq \algendtime
                    }
                }
                % \cdot
                \rbr{
                    \numdataa,\numdatab
                } .
        \]
        We claim that 
        \(
            \cbr{\timeind<\algendtime}
            \in
            \statefiltss{
                \strategyt
            }
        \). To see this, observe that
        \[
            \cbr{\timeind<\algendtime}
            =
            \bigcap_{\ell=0}^\timeind
            \cbr{
                \algstopmaptt{\ell}\rbr{
                    \cbr{
                        \valaii{k}
                    }_{k=0}^{\numdataforatt{\ell}}
                    ,
                    \cbr{
                        \valbjj{k}
                    }_{k=0}^{\numdataforbtt{\ell}}
                }
                = 1
            } .
        \]
        and from the definition of the filtration in
        ~\eqref{eq:family-sig-alg}
        \[
            \sigma\rbr{
                \cbr{
                    \valaii{k}
                }_{k=0}^{\numdataforatt{\ell}}
                ,
                \cbr{
                    \valbjj{k}
                }_{k=0}^{\numdataforbtt{\ell}}
            }
            =
            \statefiltss{
                \rbr{
                    \numdataforatt{\ell},
                    \numdataforbtt{\ell}
                }
            }
        \]
        so each respective event in the union is 
        \(
            \statefiltss{
                \rbr{
                    \numdataforatt{\ell},
                    \numdataforbtt{\ell}
                }
            }
        \)-measurable. Since $\algstatett{1}\leq\algstatett{2}$ if and only if 
        \(
            \statefiltss{
                \algstatett{1}
            }
            \subseteq
            \statefiltss{
                \algstatett{2}
            }
        \), we have 
        \(
            \statefiltss{
                \rbr{
                    \numdataforatt{\ell},
                    \numdataforbtt{\ell}
                }
            }
            \subseteq
            \statefiltss{
                \rbr{
                    \numdataforat,
                    \numdataforbt
                }
            }
            =
            \statefiltss{
                \strategyt
            }
        \), which implies
        \(
            \cbr{
                % \stoptime\leq\timeind
                \timeind < \algendtime
            }
        \)
        is 
        \(
            \statefiltss{
                \strategyt
            }
        \)-measurable.
        Since 
        \[
            \sigma\rbr{
                \cbr{
                    \valaii{k}
                }_{k=0}^{\numdataforat}
                ,
                \cbr{
                    \valbjj{k}
                }_{k=0}^{\numdataforbt}
            }
            =
            \statefiltss{
                \rbr{
                    \numdataforat,
                    \numdataforbt
                }
            },
        \]
        we get 
        \(
            \algstatemapt\rbr{
                \cbr{
                    \valaii{k}
                }_{k=0}^{\numdataforatt{t}}
                ,
                \cbr{
                    \valbjj{k}
                }_{k=0}^{\numdataforbtt{t}}
            }
        \)
        is 
        \(
            \statefiltss{
                \rbr{
                    \numdataforat,
                    \numdataforbt
                }
            }
        \)-measurable, thus 
        $\strategytt{\timeind+1}$ is too as the composition and product of measurable functions.


        % \mycomment{need to argue that the selection is also measurable}

        % This holds automatically for $\timeind\geq\algendtime$ since 
        % $\strategytt{\timeind+1}$ is constant. For $\timeind<\algendtime$,
        % $\strategytt{\timeind+1}$ is a measurable function of 
        % \(
        %     \rbr{
        %         \cbr{
        %             \valaii{k}
        %         }_{k=0}^{\numdataforatt{t}}
        %         ,
        %         \cbr{
        %             \valbjj{k}
        %         }_{k=0}^{\numdataforbtt{t}}
        %     }
        % \).
        % This means $\strategytt{\timeind+1}$ is 
        % \(
        %     \sigma\rbr{
        %         \cbr{
        %             \valaii{k}
        %         }_{k=0}^{\numdataforatt{t}}
        %         ,
        %         \cbr{
        %             \valbjj{k}
        %         }_{k=0}^{\numdataforbtt{t}}
        %     }
        % \)-measurable.
        % Moreover, from the definition of the filtration in~\eqref{eq:family-sig-alg},
        % \[
        %     \sigma\rbr{
        %         \cbr{
        %             \valaii{k}
        %         }_{k=0}^{\numdataforatt{t}}
        %         ,
        %         \cbr{
        %             \valbjj{k}
        %         }_{k=0}^{\numdataforbtt{t}}
        %     }
        %     =
        %     \statefiltss{
        %         \rbr{
        %             \numdataforat,\numdataforbt
        %         }
        %     }   
        %     =
        %     \statefiltss{
        %         \strategyt
        %     } .
        % \]
        % % so $\strategytt{\timeind+1}$
        % % is
        % % \(
        % %     \statefiltss{
        % %         \strategyt
        % %     }
        % % \)-measurable

        \item
        \big($\stoptime$ is a stopping time w.r.t. the filtration
        \(
            \cbr{
                \timefiltt
            }_{\timeind\geq 0}
            =
            \cbr{
                \statefiltss{
                    \strategyt
                }
            }_{\timeind\geq 0}
        \)\big):
        Start by rewriting
        \[
            \cbr{
                \stoptime\leq\timeind
            }
            =
            \bigcup_{\ell=0}^\timeind
            \cbr{
                \algstopmaptt{\ell}\rbr{
                    \cbr{
                        \valaii{k}
                    }_{k=0}^{\numdataforatt{\ell}}
                    ,
                    \cbr{
                        \valbjj{k}
                    }_{k=0}^{\numdataforbtt{\ell}}
                }
                = 0
            } .
        \]
        Now notice that 
        \[
            \sigma\rbr{
                \cbr{
                    \valaii{k}
                }_{k=0}^{\numdataforatt{\ell}}
                ,
                \cbr{
                    \valbjj{k}
                }_{k=0}^{\numdataforbtt{\ell}}
            }
            =
            \statefiltss{
                \rbr{
                    \numdataforatt{\ell},
                    \numdataforbtt{\ell}
                }
            }
        \]
        so each respective event in the union is 
        \(
            \statefiltss{
                \rbr{
                    \numdataforatt{\ell},
                    \numdataforbtt{\ell}
                }
            }
        \)-measurable. Since $\algstatett{1}\leq\algstatett{2}$ if and only if 
        \(
            \statefiltss{
                \algstatett{1}
            }
            \subseteq
            \statefiltss{
                \algstatett{2}
            }
        \), we have 
        \(
            \statefiltss{
                \rbr{
                    \numdataforatt{\ell},
                    \numdataforbtt{\ell}
                }
            }
            \subseteq
            \statefiltss{
                \rbr{
                    \numdataforat,
                    \numdataforbt
                }
            }
            =
            \statefiltss{
                \strategyt
            }
        \), which implies
        \(
            \cbr{
                \stoptime\leq\timeind
            }
        \)
        is 
        \(
            \statefiltss{
                \strategyt
            }
            =
            \timefiltt
        \)-measurable.
        Therefore, 
        $\stoptime$ is a stopping time with respect to the filtration
        \(
            \cbr{
                \timefiltt
            }_{\timeind\geq 0}
        \).
        
    \end{enumerate}  

    We are left with the two additional requirements of 
    Definition~\ref{def:exchange-policy}. For the first requirement, suppose 
    that
    \(
        \strategyt <
        \rbr{
            \numdataa, \numdatab
        }
    \), then by construction this implies $\timeind<\algendtime=\stoptime$
    so $\strategyt=\algstatet<\algstatett{\timeind+1}=\strategytt{\timeind+1}$
    by Definition~\ref{def:exchange-alg}.
    For the second requirement, we know that $\stoptime=\algendtime$ is finite or else 
    it would contradict the fact that 
    \(
        \algstatett{\algendtime}
        \leq
        \rbr{
            \numdataa,\numdatab
        }
    \), since $\cbr{\algstatet}_{\timeind\geq 0}$ is increasing.
    From our construction, we have that 
    \(
        \rbr{
            \algstatett{0},
            \ldots,
            \algstatett{\algendtime}
        }
        =
        \rbr{
            \strategytt{0},
            \ldots,
            \strategytt{\stoptime}
        }
    \), therefore
    \(
        \EE\sbr{
            \rewardss{
                \strategytt{\stoptime}
            }
        }
        =
        \EE\sbr{
            \rewardss{
                \algstatett{\algendtime}
            }
        }
    \).

    For the reverse direction, we create an exchange algorithm given an
    exchange policy 
    \(
        \rbr{
            \cbr{
                \strategyt
            }_{\timeind\geq 0},
            \stoptime 
        }
    \).
    Let 
    \(
        \rbr{
            \numdataforat,
            \numdataforbt
        }
        =
        \strategyt
    \).
    From Definition~\ref{def:exchange-policy} we know $\stoptime$ is a stopping time,
    thus
    \(
        1_{
            \cbr{
                \stoptime\neq\timeind
            }
        }
    \)
    is
    \[
        \statefiltss{\strategyt}
        =
        \sigma\rbr{
            \cbr{
                \valaii{k}
            }_{k=0}^{\numdataforat}
            ,
            \cbr{
                \valbjj{k}
            }_{k=0}^{\numdataforbt}
        }
    \]measurable.
    % is 
    % \(
    % \)-measurable, 
    Therefore,
    there exists a measurable function $\algstopmapt$ such that 
    \[
        1_{
            \cbr{
                \stoptime\neq\timeind
            }
        }
        =
        \algstopmapt\rbr{
            \cbr{
                \valaii{k}
            }_{k=0}^{\numdataforat}
            ,
            \cbr{
                \valbjj{k}
            }_{k=0}^{\numdataforbt}
        }.
    \]
    Similarly, since $\strategytt{\timeind+1}$ is 
    \(
        \statefiltss{\strategyt}
    \)-measurable,
    there exists a measurable function 
    \(
        \mbleexistsalgstatemapt
    \)
    such that
    \[
        \strategytt{\timeind+1}
        =
        \mbleexistsalgstatemapt
        \rbr{
            \cbr{
                \valaii{k}
            }_{k=0}^{\numdataforat}
            ,
            \cbr{
                \valbjj{k}
            }_{k=0}^{\numdataforbt}
        } .
    \]
    Now define 
    \(
        \cbr{
            \algstatemapt
        }_{\timeind\geq 0}
    \) 
    via
    \[
        \algstatemapt
        =
        \begin{cases}
            \mbleexistsalgstatemapt
            &
            \text{if }~~~
            \algstatet < 
            \rbr{
                \numdataa,\numdatab
            }
            % \forall r\in\cbr{0,\ldots,\timeind}
            % \algstopmapt\rbr{
            %     \cbr{
            %         \valaii{k}
            %     }_{k=0}^{\numdataforatt{r}}
            %     ,
            %     \cbr{
            %         \valbjj{k}
            %     }_{k=0}^{\numdataforbtt{r}}
            % }= 0
            \\
            \algstatet+(1,1)
            & 
            \text{otherwise}
        \end{cases} .
    \]
    By the requirement that
    \(
        \strategyt < \strategytt{\timeind+1}
    \)
    for 
    \(
        \strategyt <
        \rbr{
            \numdataa,\numdatab
        }
    \), we get
    $\algstatet <\algstatett{\timeind+1}$.
    Moreover, by our construction
    \(
        \rbr{
            \algstatett{0},
            \ldots,
            \algstatett{\algendtime}
        }
        =
        \rbr{
            \strategytt{0},
            \ldots,
            \strategytt{\stoptime}
        },
    \)
    % we also have that 
    which implies
    \(
        \algstatett{\algendtime}
        \leq
        \rbr{
            \numdataa,\numdatab
        }
    \)
    and 
    \(
        \EE\sbr{
            \rewardss{
                \strategytt{\stoptime}
            }
        }
        =
        \EE\sbr{
            \rewardss{
                \algstatett{\algendtime}
            }
        }
    \).
\end{proof}  

Therefore, going forward, we will refer to exchange algorithms and exchange policies interchangeably.

% \mycomment{delete below ->}
% \begin{restatable}{definition}{exchange-alg-as-policy2}
%     \label{def:exchange-alg-as-policy2} 
%     A resource exchange algorithm is a tuple 
%     \(
%         \rbr{
%             \cbr{
%                 \strategyt
%             }_{\timeind\geq 0},
%             \stoptime 
%         }
%     \).
%     Here
%     \(
%         \strategyt\in\setexchanges
%     \)
%     is a non-decreasing sequence
%     satisfying
%     \begin{align*}
%         \strategytt{0}=0,
%         \quad
%         \strategytt{\timeind+1}
%         \text{ is }
%         \statefiltss{
%             \strategyt
%         }\text{-measurable},
%         \quad
%         % \text{and }
%         \strategyt < \strategytt{\timeind+1}
%         \text{ for }
%         \strategyt < \rbr{
%             \numdataa,\numdatab
%         },
%     \end{align*}  
%     and $\stoptime$ is a finite stopping time with respect to the filtration 
%     \(
%         \cbr{
%             \timefiltt
%         }_{\timeind\geq 0}
%         =
%         \cbr{
%             \statefilt_{
%                 \strategyt
%             }
%         }_{\timeind\geq 0}
%     \)
%     as given in~\eqref{eq:family-sig-alg}.
% \end{restatable}  
% \mycomment{delete above <-}

% Definition~\ref{def:exchange-alg-as-policy} is exactly the definition of a policy in~\citep{mandelbaum1981optimal}
% where the state space is $\setexchanges$, and the family of $\sigma$-algebras over the state space is 
% \(
%     \cbr{
%         \statefilts
%     }_{\exchange\in\setexchanges}
% \).
% If, in the language of~\citep{mandelbaum1981optimal}, 

\subsection{Reducing Optimal Exchange Algorithm Design to Optimal Stopping}
\label{subsec:reduction-to-optimal-stopping}

We will now show that 
by defining the reward for an exchange $\exchange\in\setexchanges$ to be 
$\rewards$, 
the problem of designing an optimal exchange algorithm $\alg\in\classalgs$,
will reduce to solving an optimal stopping time problem over $\setexchanges$.
Before doing this, we first introduce a 
% few necessary concepts 
necessary concept
from optimal stopping theory.
% and~\citet{mandelbaum1981optimal}.

\begin{restatable}{definition}{partially-ordered-snell-env}
    \textup{
        (Snell envelope, \citet{mandelbaum1981optimal})
    }
    For any $\osasstate\in\osasstates$, let
    \[
        \osassnellenvs
        =
        \esssup_{
            \osaspolicy
            \in
            \osaspoliciess
        }
        \EE\sbr{
            \osasrewardss{
                \osasstrategytt{
                    \osasstoptime
                }
            }
            | \osasstatefilts
        }
    \]
    where 
    $\osaspoliciess$ is the set of policies for which 
    \(
        \osasstrategytt{0} = \osasstate
    \).
    The Snell envelope is defined as the family 
    \(
        \cbr{
            \osassnellenvs
        }_{
            \osasstate\in\osasstates
        }
    \).
\end{restatable}  

Intuitively, each $\osassnellenvs$ represents the best expected reward, given the information available
at $\osasstate$, following policies that start at $\osasstate$. The definition of the Snell envelope 
depends on the type of optimal stopping problem considered; however, it is usually difficult to calculate 
$\osassnellenvs$ directly from the definiton. To remidy this, a dymanic programming equation is often
proved. 
This equation usually takes the following form: $\osassnellenvs$ is the maximum of stopping and accepting
the current reward or the expected reward of ``proceeding optimally''.
Using this, we now present an optimal exchange algorithm.


% The idea is to show that $\osassnellenvs$ is the maximum of the following two choices: stop and accept the 
% current reward or 

% The following is such a result in the case of optimal stopping times over partially ordered sets.

% \mycomment{move this to appendix but cite here, describe high level}
% \kkcomment{do this.}

% \begin{restatable}{theorem}{osas-dp-eq}
%     \label{thm:osas-dp-eq}
%     \textup{
%         (\citet{mandelbaum1981optimal})
%     }
%     Let 
%     \(
%         \osasdirectsuccessorss{
%             \osasstate
%         }
%         =
%         \cbr{
%             u> \osasstate: \sbr{\osasstate,u}=\cbr{\osasstate,u}
%         } 
%     \)
%     be the set of direct successors for $\osasstate$, then
%     \begin{align*}
%         \osassnellenvs
%         =
%         \max
%         \cbr{
%             \osasrewards,
%             \EE\sbr{
%                 \osassnellenvss{u}
%                 | 
%                 \osasstatefilts
%             }
%             : u\in
%             \osasdirectsuccessorss{
%                 \osasstate
%             }
%         }
%         \quad
%         \text{and}
%         \quad
%         \EE\sbr{
%             \osassnellenvss{0}
%         }
%         =
%         \sup_{
%             \osaspolicy\in\osaspolicies
%         }
%         \EE\sbr{
%             \osasrewardss{
%                 \osasstrategytt{
%                     \osasstoptime
%                 }
%             }
%         }
%         .
%     \end{align*}  
% \end{restatable}

% The first part of this result says that the best expected reward of a policy which starts from state
% $\osasstate\in\osasstates$ is the maximum of the reward available from stopping at $\osasstate$,
% $\osasrewards$, and the best 
% expected reward among policies starting at a successor $u$, conditioned on the information 
% available at $\osasstate$. The second part shows that $\osassnellenvss{0}$ characterizes the 
% performance of an optimal policy. Using this, we now present an optimal exchange algorithm,

% \mycomment{remove all the tildes? Unnecessary clutter? I did this with the IC section}

% \kkcomment{Since this is a key result of the paper, you can also move the IR rsult and proof into this theorem. Currently, it looks very disjointed. Also, better to state this in a way where you are presenting the }

\begin{restatable}{theorem}{opt-alg}
    \label{thm:opt-alg}
    % In the setting of~\citep{mandelbaum1981optimal}, 
    Let $\osasstates=\setexchanges$, $\osasrewards=\rewards$,
    and the family 
    \(
        \cbr{
            \statefilts
        }_{\exchange\in\setexchanges}
    \) be given by~\eqref{eq:family-sig-alg}.
    Define
    the policy
    \(
        \optexpolicy=
        \rbr{
            \cbr{
                \optstrategyt
            }_{\timeind\geq 0},
            \optstoptime
        }
    \)
    via
    %  $\strategytt{0}=0$, for $\timeind\geq 1$
    \begin{align*}
        \optstrategytt{0} = (0,0),
        \quad
        % \text{for } \timeind\geq 1
        \optstrategytt{\timeind+1}
        =
        \argmax_{
            \nextexchange\in 
            \directsuccessorss{
                \optstrategyt
            }
        }
        \EE\sbr{
            \snellenvss{\nextexchange}
            | \timefiltt
        }
        ~~~\text{ for }~~~~\timeind\geq 0,
        % \quad
        % \text{and}
        % \quad
        % \stoptime
    \end{align*}  
    % where 
    % \(
    %     \directsuccessorss{
    %         % \optstrategyt
    %         \exchange
    %     }
    %     =
    %     \cbr{
    %         \exchange + (1,0),
    %         \exchange + (0,1)
    %     }
    % \),
    and 
    \(
        \optstoptime
        =
        \inf
        \cbr{
            t\in\NN_{\geq 0} :
            \snellenvss{
                \optstrategyt
            }
            =
            \rewardss{
                \optstrategyt
            }
        }
    \). 
    Then $\optexpolicy\in\setexchangepolicies$ and is optimal, i.e.
    % \begin{align*}
    \(
        \EE\sbr{
            \rewardss{
                \optstrategytt{\optstoptime}
            }
        }
        =
        \sup_{
            % \rbr{
            %     \cbr{\strategyt}_{\timeind\geq 0},
            %     \stoptime
            % }
            % \in\classalgs
            \exchangepolicy
            \in
            \setexchangepolicies
        }
        \EE\sbr{
            \rewardss{
                \strategytt{\stoptime}
            }
        }
    \). 
    Moreover if
    for the choice 
    of welfare function $\welfarefunc$,
    \[
        \EE\sbr{
            \welfarefunc\rbr{
                \utilass{
                    \optstrategytt{\optstoptime}
                },
                \utilbss{
                    \optstrategytt{\optstoptime}
                }
            }
        }
        \geq 0
        ~~~
        \Rightarrow
        ~~~
        \EE\sbr{
            \utilass{
                \optstrategytt{\optstoptime}
            }
        },
        \EE\sbr{
            \utilbss{
                \optstrategytt{\optstoptime}
            }
        }
        \geq 0
    \]
    then $\optexpolicy$ is individually rational.
\end{restatable}  
\begin{proof}
    We first verify that $\optexpolicy$ is a well defined policy.
    Since 
    \(
        \timefiltt
        =
        \statefiltss{
            \optstrategyt
        }
    \),
    \(
        \optstrategytt{\timeind+1}
    \)
    is the argmax of a finite number of
    \(
        \statefiltss{
            \optstrategyt
        }
    \)-measurable random variables, 
    and thus 
    \(
        \optstrategytt{\timeind+1}
    \)
    is
    \(
        \statefiltss{
            \optstrategyt
        }
    \)-measurable.
    From the definition of 
    \(
        \directsuccessorss{
            \optstrategyt
        }
    \), it follows that
    \(
        \optstrategyt
        \leq
        \optstrategytt{\timeind+1}
    \).
    Now observe that 
    % \begin{align*}
    \(
        \cbr{
            \optstoptime\leq \timeind
        }
        =
        \bigcup_{k=0}^\timeind
        \cbr{
            \snellenvss{
                \optstrategytt{k}
            }
            =
            \rewardss{
                \optstrategytt{k}
            }
        }
    \)
    and 
    \(
        \cbr{
            \snellenvss{
                \optstrategytt{k}
            }
            =
            \rewardss{
                \optstrategytt{k}
            }
        }
        \in
        \timefilttt{k}
        \subseteq
        \timefiltt
    \)
    by the definition in~\eqref{eq:family-sig-alg}.
    Thus
    \(
        \cbr{
            \optstoptime\leq \timeind
        }
        \in
        \timefiltt
    \)
    so $\optstoptime$ is a well defined stopping time with respect to 
    \(
        \cbr{
            \timefiltt
        }_{\timeind\geq 0}
    \).

    We now argue $\optexpolicy$ is a valid exchange policy. Since at time $\timeind\geq 0$ the next exchange
    is selected from 
    \(
        \directsuccessorss{
            \optstrategyt
        }
    \),
    we have
    \(
        \optstrategyt < \optstrategytt{\timeind+1} 
    \).
    Now suppose that $\optstoptime$ is not finite. Then, for sufficiently a sufficiently large $T$,
    we would have 
    \(
        \optstrategytt{T}
        =
        \rbr{
            \numdataa,\numdatab
        }
    \).
    Since 
    \(
        \directsuccessorss{
            \optstrategytt{T}
        }
        =
        \varnothing
    \),
    Theorem~\ref{thm:osas-dp-eq}
    implies that 
    \(
        \snellenvss{
            \optstrategytt{T}
        }
        =
        \rewardss{
            \optstrategytt{T}
        }
    \), so $\optstoptime$ is forced to be finite.

    By construction, $\optexpolicy$
    is the natural policy which satisfies conditions~\eqref{eq:osas-thm2-cond1}--\eqref{eq:osas-thm2-cond3}.
    %  (2.5)--(2.7) in~\citep{mandelbaum1981optimal} 
    %  \kkcomment{State the theorems you are using in the appendix and refer to the conditions here.}. 
    Therefore, together with the fact that $\optstoptime$ is finite, we can apply 
    Theorem~\ref{thm:osas-thm2}
    % Theorem 2
    % from~\citep{mandelbaum1981optimal} 
    to conlude $\optexpolicy$ is an optimal policy, i.e.
    \(
        \EE\sbr{
            \rewardss{
                \optstrategytt{\optstoptime}
            }
        }
        =
        \sup_{
            \apolicy\in
            \osaspolicies
        }
        \EE\sbr{
            \rewardss{
                \strategytt{\stoptime}
            }
        }
    \).
    Since $\setexchangepolicies\subseteq\osaspolicies$, and $\optexpolicy\in\setexchangepolicies$ 
    \[
        \sup_{
            \apolicy\in
            % \classalgs
            \setexchangepolicies
        }
        \EE\sbr{
            \rewardss{
                \strategytt{\stoptime}
            }
        }
        \leq
        \sup_{
            \apolicy\in
            \osaspolicies
        }
        \EE\sbr{
            \rewardss{
                \strategytt{\stoptime}
            }
        }
        =
        \EE\sbr{
            \rewardss{
                \optstrategytt{\optstoptime}
            }
        }
        \leq
        \sup_{
            \apolicy\in
            \setexchangepolicies
        }
        \EE\sbr{
            \rewardss{
                \strategytt{\stoptime}
            }
        } .
    \]
    Therefore
    \(
        \EE\sbr{
            \rewardss{
                \optstrategytt{\optstoptime}
            }
        }
        =
        \sup_{
            % \rbr{
            %     \cbr{\strategyt}_{\timeind\geq 0},
            %     \stoptime
            % }
            % \in\classalgs
            \exchangepolicy
            \in
            \setexchangepolicies
        }
        \EE\sbr{
            \rewardss{
                \strategytt{\stoptime}
            }
        }
    \),
    so $\optexpolicy$
    is an optimal exchange policy.

    For the second part of the claim,
    consider the exchange policy which facilitates no exchanges, i.e.
    \(
        \altexchangepolicy
        =
        \rbr{
            \altstrategytt{0} = (0,0),
            \altstoptime = 0
        } 
    \).
    Then by our normalization assumptions we have
    \(
        \valaii{0} = \valbjj{0} = 0
    \)
    which implies 
    \(
        \utilass{0}
        =
        \utilbss{0}
        =
        \welfarefunc\rbr{
            \utilass{0},
            \utilbss{0}
        }
        =0
    \).
    Therefore, since $\optexpolicy$ is optimal,
    \(
        0
        =
        \EE\sbr{
            \rewardss{
                \altstrategytt{\altstoptime}
            }
        }
        \leq
        \EE\sbr{
            \rewardss{
                \optstrategytt{\optstoptime}
            }
        }
        =
        \EE\sbr{
            \welfarefunc\rbr{
                \utilass{
                    \optstrategytt{\optstoptime}
                },
                \utilbss{
                    \optstrategytt{\optstoptime}
                }
            }
        }
    \),
    so our assumption yields individual rationality.
\end{proof}  

It is worth pointing out that there are reasonable choices of welfare functions that satisfy the implication
in the Theorem~\ref{thm:opt-alg}. 
A simple example of such a welfare function is the egalitarian welfare
\(
    \welfarefunc\rbr{x,y}=\min(x,y)
\)
since for any random variables $X$ and $Y$,
\(
    \EE\sbr{
        \min(X,Y)
    }
    \geq 0
    \Rightarrow
    \EE\sbr{X},
    \EE\sbr{Y}
    \geq 0
\).


\subsection{An example}
\label{sec:algoexample}

We now revisit Example~\ref{ex:multi-vs-single-round-alg} to demonstrate how Theorem~\ref{thm:opt-alg}
is gives an optimal algorithm.
We will show that the optimal exchange algorithm is given by
    \(
        \rbr{
            \cbr{\algstatemapt}_{\timeind=0}^{1},
            \cbr{\algstopmap}_{\timeind=0}^{2}
        }
    \)
    where 
    \begin{gather*}
        \algstatemaptt{0}(0,0) = (1,0)
        \qquad
        \algstatemaptt{1}(1,0) = (1,1)
        \\
        \algstopmaptt{0}(0,0) = 1
        \qquad
        \algstopmaptt{1}(1,0) 
        =
        \begin{cases}
            1 & \text{if } ~\valaii{1} = 1
            \\
            0 & \text{if } ~\valaii{1} = 0
        \end{cases}  
        \qquad
        \algstopmaptt{2}(1,1) 
        = 0 . 
    \end{gather*}  
    and moreover that,
    \(
        \EE\sbr{
            \rewardss{
                \algstatett{\algendtime}
            }
        }
        =
        \frac{1}{8}
    \).

% \begin{restatable}{claim}{two-bern-ex-opt-alg}
%     \label{ex:two-bern-ex-opt-alg}
%     % Let $\numdataa=\numdatab=1$.
%     % Suppose that each competitor's value for the other's good is 
%     % $\valaii{1},\valbjj{1}\drawniid\Bern(1/2)$ and assume $\valaii{0}=\valbjj{0}=0$.
%     % Let the utilities and welfare be given by
%     % \begin{align*}
%     %     \utilas
%     %     =
%     %     \valai-\frac{1}{4}\valbj,
%     %     \quad 
%     %     \utilbs
%     %     =
%     %     \valbj-\frac{1}{4}\valai,
%     %     \quad
%     %     \text{and}
%     %     \quad
%     %     \welfarefunc(x,y)=\min(x,y) .
%     % \end{align*}  
%     Recall the environment from Example~\ref{ex:two-bern-ex-opt-alg}.
%     The optimal exchange algorithm is given by
%     \(
%         \rbr{
%             \cbr{\algstatemapt}_{\timeind=0}^{1},
%             \cbr{\algstopmap}_{\timeind=0}^{2}
%         }
%     \)
%     where 
%     \begin{gather*}
%         \algstatemaptt{0}(0,0) = (1,0)
%         \qquad
%         \algstatemaptt{1}(1,0) = (1,1)
%         \\
%         \algstopmaptt{0}(0,0) = 0
%         \qquad
%         \algstopmaptt{1}(1,0) 
%         =
%         \begin{cases}
%             1 & \text{if } ~\valaii{1} = 0
%             \\
%             0 & \text{if } ~\valaii{1} = 1
%         \end{cases}  
%         \qquad
%         \algstopmaptt{2}(1,1) 
%         = 1 . 
%     \end{gather*}  
%     and 
%     \(
%         \EE\sbr{
%             \rewardss{
%                 \algstatett{\algendtime}
%             }
%         }
%         =
%         \frac{1}{8}
%     \).
% \end{restatable}  

\begin{proof}
    We will use Theorem~\ref{thm:opt-alg} to determine the optimal exchange policy which will correspond to 
    the optimal exchange algorithm above. 
    % \mycomment{add and thus is the opt alg by the corollary to add?}
    To determine the optimal exchange policy, we first calculate the Snell envelope
    \(
        \cbr{
            \snellenvs
        }_{\exchange\in\setexchanges}
    \) working backwards.

    From Theorem~\ref{thm:osas-dp-eq} we have
    % \begin{align*}
    \(
        \snellenvss{(1,1)}
        =
        \rewardss{(1,1)}
    \)
    % \end{align*}  
    which gives
    \begin{align*}
        \snellenvss{(1,0)}
        =
        \max\rbr{
            \rewardss{(1,0)},
            \EE\sbr{
                \snellenvss{(1,1)}
                | \statefiltss{(1,0)}
            }
        }
        =
        \max\rbr{
            \rewardss{(1,0)},
            \EE\sbr{
                \rewardss{(1,1)}
                | \valaii{1}
            }
        }
        \\
        \snellenvss{(0,1)}
        =
        \max\rbr{
            \rewardss{(0,1)},
            \EE\sbr{
                \snellenvss{(1,1)}
                | \statefiltss{(0,1)}
            }
        }
        =
        \max\rbr{
            \rewardss{(0,1)},
            \EE\sbr{
                \rewardss{(1,1)}
                | \valbjj{1}
            }
        }
    \end{align*}  
    and thus by symmetry,
    \begin{align*}
        \snellenvss{(0,0)}
        &=
        \max\rbr{
            \rewardss{(0,0)},
            \EE\sbr{
                \snellenvss{(1,0)}
                | \statefiltss{(0,0)}
            },
            \EE\sbr{
                \snellenvss{(0,1)}
                | \statefiltss{(0,0)}
            }
        }
        \\
        &=
        \max\rbr{
            \rewardss{(0,0)},
            \EE\sbr{
                \snellenvss{(1,0)}
            },
            \EE\sbr{
                \snellenvss{(0,1)}
            }
        }
        \\
        &=
        \max\rbr{
            0,
            \EE\sbr{
                \snellenvss{(1,0)}
            }
        } .
    \end{align*}  
    To determine
    \(
        \EE\sbr{
            \rewardss{(1,1)}
            | \valaii{1}
        }
    \), consider two cases.

        If $\valaii{1}=0$ then $\rewardss{(1,1)}=0$ when $\valbjj{1}=0$, and 
        $\rewardss{(1,1)}=-\frac{1}{4}$ when $\valbjj{1}=1$ so
        \[
            \EE\sbr{
                \rewardss{(1,1)}
                | \valaii{1}=0
            }
            =
            \frac{1}{2}\cdot 0
            +
            \frac{1}{2}\cdot\rbr{-\frac{1}{4}} 
            =
            -\frac{1}{8}
        \]

        If $\valaii{1}=1$ then $\rewardss{(1,1)}=-\frac{1}{4}$ when $\valbjj{1}=0$, and 
        $\rewardss{(1,1)}=\frac{3}{4}$ when $\valbjj{1}=1$ so
        \[
            \EE\sbr{
                \rewardss{(1,1)}
                | \valaii{1}=1
            }
            =
            \frac{1}{2}\cdot \rbr{-\frac{1}{4}}
            +
            \frac{1}{2}\cdot\rbr{\frac{3}{4}} 
            =
            \frac{1}{4} .
        \]

    \noindent
    Therefore, 
    \(
        \EE\sbr{
            \rewardss{(1,1)}
            | \valaii{1}
        }
        =
        \frac{3}{8}\valaii{1}
        -
        \frac{1}{8}
    \).
    We use this to calculate
    \begin{align*}
        \EE\sbr{
            \snellenvss{(1,0)}
        }
        &=
        \max\rbr{
            \rewardss{(1,0)},
            \frac{3}{8}\valaii{1}
            -
            \frac{1}{8}
        }
        \\
        &=
        \frac{1}{4}
        \max\rbr{
            \min\rbr{
                0,0
            },
            -
            \frac{1}{8}
        }
        +
        \frac{1}{4}
        \max\rbr{
            \min\rbr{
                1, -\frac{1}{4}
            },
            \frac{3}{8}-\frac{1}{8}
        }
        \\
        &~+
        \frac{1}{4}
        \max\rbr{
            \min(0,0),
            -\frac{1}{8}
        }
        +
        \frac{1}{4}
        \max\rbr{
            \min\rbr{
                1, -\frac{1}{4}
            },
            \frac{3}{8}-\frac{1}{8}
        }
        \\
        &= 
        \frac{1}{8}
    \end{align*}  
    and thus 
    \(
        \snellenvss{(0,0)}
        =
        \max\rbr{0,\frac{1}{8}}
        =
        \frac{1}{8}
    \).

    By definition, the policy starts at
    \(
        \optstrategytt{0}
        =
        (0,0)
    \).
    Because we define 
    \(
        \optstoptime
        =
        \inf
        \cbr{
            t\in\NN_{\geq 0} :
            \snellenvss{
                \optstrategyt
            }
            =
            \rewardss{
                \optstrategyt
            }
        }
    \),
    we have $\optstoptime>0$, i.e. the policy does not stop immediately and since 
    \[
        \optstrategytt{1}
        =
        \argmax_{
            \nextexchange\in 
            \cbr{
                (1,0),(0,1)
            }
        }
        \EE\sbr{
            \snellenvss{\nextexchange}
        }
    \]
    it proceeds to either $(1,0)$ $(0,1)$. Because of symmetry 
    \(
        \EE\sbr{
            \snellenvss{(1,0)}
        }
        =
        \EE\sbr{
            \snellenvss{(0,1)}
        }
    \)
    so without loss of generality $\optstrategytt{1}=(1,0)$.
    At this point
    \(
        \optstoptime
        =
        1
    \)
    exactly when
    \[
        \rewardss{(1,0)}
        \geq
        \EE\sbr{
            \rewardss{(1,1)}
            | \valaii{1}
        }
        \iff
        \valaii{1}=0,
    \]
    so the policy stops if $\valaii{1}=0$.
    Otherwise, if $\valaii{1}=1$ the policy continues to
    \[
        \optstrategytt{2}
        =
        \argmax_{
            \nextexchange\in 
            \cbr{
                (1,1)
            }
        }
        \EE\sbr{
            \snellenvss{\nextexchange}
        }
        =
        (1,1)
    \]
    as the only remaining state.
    Finally, if the policy continued, 
    then $\optstoptime=2$
    since $\snellenvss{(1,1)}=\rewardss{(1,1)}$.
    Therefore, the $\optexpolicy$ corresponds to the exchange algorithm above.
\end{proof}  

For more complex models it may be difficult to analytically compute the condition expectations
used to determine the sequence of exchanges
\(
    \cbr{
        \optstrategyt
    }_{\timeind\geq 0}
\)
and when to terminate the exchanges $\optstoptime$. This is a common phenomena in optimal 
stopping, in which case numerical approximation methods such as Monte Carlo based techniques are 
used~\citep{longstaff2001valuing, tsitsiklis2001regression, broadie1997pricing}.


% \mycomment{add corollary that we have an opt exchange alg?}

% \subsection{An IR Sufficient Condition}

% \kkcomment{Move up to main theorem.}

% In addition to being optimal, the following result shows that for a simple condition on the welfare function
% $\welfarefunc(x,y)$, the algorithm above is also individually rational.

% \begin{restatable}{proposition}{ir-welfare-condition}
%     Let $\optalg\in\classalgs$ be the algorithm given in Theorem~\ref{thm:opt-alg}. 
%     If for the choice 
%     of welfare function $\welfarefunc$,
%     \[
%         % \EE\sbr{
%         %     \rewardss{
%         %         \optstrategytt{\optstoptime}
%         %     }
%         % }
%         \EE\sbr{
%             \welfarefunc\rbr{
%                 \utilass{
%                     \optstrategytt{\optstoptime}
%                 },
%                 \utilbss{
%                     \optstrategytt{\optstoptime}
%                 }
%             }
%         }
%         \geq 0
%         ~~~
%         \Rightarrow
%         ~~~
%         \EE\sbr{
%             \utilass{
%                 \optstrategytt{\optstoptime}
%             }
%         },
%         \EE\sbr{
%             \utilbss{
%                 \optstrategytt{\optstoptime}
%             }
%         }
%         \geq 0
%     \]
%     then $\optalg$ is individually rational.
% \end{restatable}  
% \begin{proof}
%     Consider the algorithm which facilitates no exchanges, i.e.
%     \(
%         \alg
%         =
%         \rbr{
%             \strategytt{0} = (0,0),
%             \stoptime = 0
%         } 
%     \).
%     Then by our normalization assumptions we have
%     \(
%         \valaii{0} = \valbjj{0} = 0
%     \)
%     which implies 
%     \(
%         \utilass{0}
%         =
%         \utilbss{0}
%         =
%         \welfarefunc\rbr{
%             \utilass{0},
%             \utilbss{0}
%         }
%         =0
%     \).
%     Therefore, since $\optalg$ is optimal,
%     \(
%         0
%         =
%         \EE\sbr{
%             \rewardss{
%                 \strategytt{\stoptime}
%             }
%         }
%         \leq
%         \EE\sbr{
%             \rewardss{
%                 \optstrategytt{\optstoptime}
%             }
%         }
%         =
%         \EE\sbr{
%             \welfarefunc\rbr{
%                 \utilass{
%                     \optstrategytt{\optstoptime}
%                 },
%                 \utilbss{
%                     \optstrategytt{\optstoptime}
%                 }
%             }
%         }
%     \)
%     ,
%     so our assumption yields individual rationality.
% \end{proof}  

% A simple example of such a welfare function is 
% \(
%     \welfarefunc\rbr{x,y}=\min(x,y)
% \)
% since for any random variables $X$ and $Y$,
% \(
%     \EE\sbr{
%         \min(X,Y)
%     }
%     \geq 0
%     \Rightarrow
%     \EE\sbr{X},
%     \EE\sbr{Y}
%     \geq 0
% \).
% \kkcomment{When you move this up, state that the other welfares Nash s.t IR or social welfare s.t IR also work.}
% \mycomment{not necessarily true}

\subsection{Impossibility Under Strategic Reporting}

Thus far we have assumed that the agents report their valuations truthfully or that they can be publicly 
verified. This assumption is necessary to design a performant exchange algorithm; otherwise, competitors 
could misreport their valuations to manipulate the algorithm, increasing their own utility at the expense of overall welfare.
The following result demonstrates that the gap between the optimal algorithm when competitors are truthful and the optimal
algorithm when competitors are untruthful can be arbitrarily large.

\begin{restatable}{proposition}{optTruthfulVsOptUntruthful}
    \label{prop:opt-truthful-vs-opt-untruthful}
    For any $k\geq 0$ there exists an environment such that  
    \(
        \EE\sbr{
            \rewardss{
                \optstrategytt{
                    \optstoptime
                }
            }
        }
        -
        \EE\sbr{
            \rewardss{
                \opttruthstrategytt{
                    \opttruthstoptime
                }
            }
        }
        \geq k
    \),
    where 
    \(
        \optexpolicy
        =
        \rbr{
            \cbr{
                \optstrategyt
            }_{\timeind\geq 0},
            \optstoptime
        }
    \)
    and 
    \(
        \opttruthexpolicy
        =
        \rbr{
            \cbr{
            \opttruthstrategyt
            }_{\timeind\geq 0},
            \opttruthstoptime
        }
    \)
    are the optimal policies assuming and not assuming truthful reporting respectively.
    % \kkcomment{This should be stated more clearly.
    % You have defined $\optexpolicy$
    % and 
    % $\opttruthexpolicy$, but not used them. On the other hand, you haven't defined $\EE\sbr{
    %         \rewardss{
    %             \opttruthstrategytt{
    %                 \opttruthstoptime
    %             }
    %         }
    %     }$. }
    % \mycomment{put =}
\end{restatable}  
\begin{proof}
    Let $\numdataa=\numdatab=1$.
    Consider the valuations
    \begin{align*}
        \valaii{0}=\valbjj{0}=0
        \qquad
        \valaii{1}
        =
        \begin{cases}
            4k
            &
            \text{w.p. } \frac{1}{2}
            \\
            % \frac{k}{2}
            2k
            &
            \text{w.p. } \frac{1}{2}
        \end{cases}  
        \qquad
        \valbjj{1} = 4k
    \end{align*}  
    and let the competitors' utilities and the welfare function be given by
    \begin{align*}
        \utila(x,y) 
        =
        x-y
        \qquad
        \utilb(x,y)
        =
        \begin{cases}
            0 
            &
            \text{if }
            (y-x)\geq -2k
            %  k\geq x-y
            \\
            2(y-x)
            &
            \text{otherwise}
        \end{cases}  
        \qquad
        \welfarefunc(x,y)=x+y .
    \end{align*}
    First consider when competitors are truthful. 
    Under this model, the optimal policy 
    \(
        \optexpolicy
    \)
    first gives $\competitora$ competitor $\competitorb$'s point,
    if 
    $\valaii{1}=2k$ 
    then sharing terminates, otherwise $\competitorb$ is given
    $\competitora$'s point to ensure $\competitorb$'s utility does not lead to a negative welfare.
    Thus the expected welfare under this policy is 
    \[
        \EE\sbr{
            \rewardss{
                \stopstateforpolicy{
                    \optexpolicy_k
                }
            }
        }
        =
        \frac{1}{2}
        \welfarefunc\rbr{
            \utilass{
                (1,0)
            },
            \utilbss{
                (1,0)
            }
        }
        +
        \frac{1}{2}
        \welfarefunc\rbr{
            \utilass{
                (1,1)
            },
            \utilbss{
                (1,1)
            }
        }
        =
        \frac{1}{2}
        \rbr{
            % \frac{k}{2}
            2k
            + 0 
        }
        +
        \frac{1}{2}
        \rbr{
            0 + 0
        } 
        % =\frac{k}{4}
        =
        k.
    \] 
    Now consider when competitors are untruthful. 
    In this setting, the optimal truthful policy will
    not always 
    give $\competitorb$ competitor
    $\competitora$'s point. This is because doing so will always lead to a welfare of no more than 0,
    thus such a policy is no better than one which exchanges zero resources.
    Now consider a policy which only gives $\competitorb$ the point based on what is reported as $\valaii{1}$.
    In this case $\competitora$ is incentivized to always report the value for $\valaii{1}$ for which $\competitorb$
    does not receive the point as it increases their utility. However, this leads to a negative welfare.

    Therefore, the optimal policy when competitors are untruthul recommends no exchange so 
    \(
        \EE\sbr{
            \rewardss{
                \stopstateforpolicy{
                    \opttruthexpolicy_k
                }
            }
        }
        =0
    \)
    and thus
    \(
        \EE\sbr{
            \rewardss{
                \optstrategytt{
                    \optstoptime
                }
            }
        }
        -
        \EE\sbr{
            \rewardss{
                \opttruthstrategytt{
                    \opttruthstoptime
                }
            }
        }
        =k
    \).
\end{proof}  

For any $k\geq 0$ the proof of Proposition~\ref{prop:opt-truthful-vs-opt-untruthful} constructs an environment
where $\numdataa=\numdatab=1$ and $\valaii{1}$ is random but $\valbjj{1}$ is deterministic. When 
competitors are truthful, an optimal algorithm
relies on learning the realization of $\valaii{1}$ to determine whether to give $\competitorb$ competitor $\competitora$'s
resource. However, when competitors may be untruthful,
$\competitora$ is always incentivized to report whichever value leads to $\competitorb$ not obtaining
the resource. In this case the optimal exchange algorithm is to not recommend any exchanges.
\section{An Optimal IR and IC Resource Exchange Algorithm}
\label{sec:opt-ic-alg}

The previous section proves that,
under the assumption that competitors do not exit the sharing algorithm early,
we can construct an optimal exchange algorithm. Moreover, we have demonstrated a simple sufficient condition
on the welfare function $\welfarefunc$ for which this optimal algorithm is also IR. A natural next question is 
whether the algorithm is also automatically IC?
The answer to this question turns out to be no and is illustrated by revisiting Example~\ref{ex:multi-vs-single-round-alg}. 

\begin{restatable}{example}{opt-alg-not-always-ic}
    Recall the environment from Example~\ref{ex:multi-vs-single-round-alg}. In~\S\ref{sec:algoexample}
    %Example~\ref{ex:two-bern-ex-opt-alg}
    we showed that the optimal exchange algorithm 
    \(
        \optalg
        =
        \rbr{
            \cbr{\algstatemapt}_{\timeind=0}^{1},
            \cbr{\algstopmap}_{\timeind=0}^{2}
        }
    \)
    first gives $\competitora$ competitor $\competitorb$'s point,
    and if $\valaii{1}=1$, competitor $\competitorb$ receives $\competitora$'s points. However, consider when
    $\competitora$ exits sharing after receiving $\competitorb$ but before $\competitorb$ receives
    $\competitora$'s point. This corresponds to the exit sequence 
    \(
        \rbr{
            \altalgstopmaptt{0} = 0,
            \altalgstopmaptt{1} = 1
        }
    \) so 
    $\altalgendtime=1$ (see Definition~\ref{def:ic-exchange-alg}). But notice that since 
    $\algendtime\geq 1$, 
    % when $\competitora$ exits at $\timeind=1$
    $\competitora$'s utility when exiting early is
    \begin{align*}
        \EE\sbr{
            \utilass{
                \algstatett{
                    \algendtime\wedge\altalgendtime
                }
            }
        }
        =
        \EE\sbr{
            \utilass{
                \algstatett{
                    1
                }
            }
        }
        =
        \EE\sbr{
            \utilass{
                (1,0)
            }
        }
        =
        \frac{1}{2}\cdot 1 + \frac{1}{2} \cdot 0
        =
        \frac{1}{2}
    \end{align*}  
    whereas if $\competitora$ does not exit early his utility is
    \begin{align*}
        \EE\sbr{
            \utilass{
                \algstatett{
                    \algendtime
                }
            }
        }
        =
        \frac{1}{2}\utila\rbr{0,0}
        +
        \frac{1}{2}\rbr{
            \frac{1}{2}\utila\rbr{1,0}
            +
            \frac{1}{2}\utila\rbr{1,1}
        }
        =
        0 + \frac{1}{4}+\frac{3}{16}
        =
        \frac{7}{16} .
    \end{align*}  
    Since $\competitora$ obtains a strictly higher expected utilty by exiting early, $\optalg$ is not IC 
    according to Definition~\ref{def:ic-exchange-alg}
\end{restatable}  

\subsection{Reformulating IC as a Stopping Time Condition}

% In order to analyze IC algorithms in the language of optimal stopping, we define 
Previously, we simplified algorithm design by reformulating our problem in the language of optimal 
stopping and showing a natural correspondence between exchange algorithms and exchange policies. 
% To analyze IC algorithms, we now take the same approach; the following definition 
We now take a similar approach and start by defining what it means for an exchange policy to be IC.
% from which we can
% to which there is a corresponding IC exchange algorithm.

\begin{restatable}{definition}{ic-exchange-policy}
    \label{def:ic-exchange-policy} 
    Let 
    \(
        \exchangepolicy
        =
        \rbr{
            \cbr{
                \strategyt
            }_{\timeind\geq 0},
            \stoptime
        }
        \in\setexchangepolicies
    \),
    we say that $\exchangepolicy$ is IC if 
    for any stopping time $\altstoptime$
    with respect to the filtration
    \(
        \cbr{
            \timefiltt
        }_{\timeind\geq 0}
        =
        \cbr{
            \statefiltss{
                \strategyt
            }
        }_{\timeind\geq 0}
    \),
    \[
        \EE\sbr{
            \utilass{
                \strategytt{
                    \stoptime\wedge\altstoptime
                }
            }
        }
        \leq
        \EE\sbr{
            \utilass{
                \strategytt{
                    \stoptime
                }
            }
        }
        \quad
        \text{and}
        \quad
        \EE\sbr{
            \utilbss{
                \strategytt{
                    \stoptime\wedge\altstoptime
                }
            }
        }
        \leq
        \EE\sbr{
            \utilbss{
                \strategytt{
                    \stoptime
                }
            }
        }.
    \]
    Let $\seticexchangepolicies$ denote the set of IC exchange policies. 
    We say that $\exchangepolicy\in\seticexchangepolicies$
    is IC optimal if 
    \[
        \EE\sbr{
            \rewardss{
                \strategytt{\stoptime}
            }
        }
        =
        \sup_{
            \altexchangepolicy\in\seticexchangepolicies
        }
        \EE\sbr{
            \rewardss{
                \strategytt{\stoptime'}'
            }
        } .
    \]
\end{restatable}  

The definition above is the stopping time analogue of Definition~\ref{def:ic-exchange-alg}. At a high level
it says that no matter what early exit strategy either competitor employs, they cannot improve their expected
utility by stopping collaboration early. By repeating the arguments in Theorem~\ref{thm:equiv-alg-defs} we can show that 
for any IC algorithm there is a corresponding IC policy with the same expected welfare and vice versa. Thus, 
we will work with IC policies going forward.

\subsection{Closure Under Pasting}

% A common theme in optimal stopping is that 
A useful property in optimal stopping problems is what we will refer to as closure under pasting. This property
describes the fact that given any two stopping times $\tau$ and $\sigma$ with respect to a filtration $\cbr{\Fcal_n}_{n\geq 0}$,
and an event $A\in\Fcal_\tau$ we can construct a new stopping time $\theta=1_A \tau+1_{A^c}\sigma$. Similarly,
the results in~\citet{mandelbaum1981optimal} rely on an analogue of this property for policies, that a third policy can be created
by pasting together two policies. More precisely we mean that for 
$\exchange\in\setexchanges$, 
$A\in\statefilts$,
and 
\(
    \apolicy^i
    =
    \Big( 
        \cbr{
            \strategyt^i
        }_{\timeind\geq 0}
        \stoptime^i
    \Big)
    \in\osaspoliciesss{\exchange}
\),
$i\in\cbr{1,2}$,
we can define 
\(
    \apolicy^3
    =
    \Big( 
        \cbr{
            \strategyt^3
        }_{\timeind\geq 0}
        \stoptime^3
    \Big)
    \in\osaspoliciesss{\exchange}
\)
via
\[
    \strategyt^3
    =
    \strategyt^1 1_A
    +
    \strategyt^2 1_{A^c}
    \qquad
    \stoptime^3
    =
    \stoptime^1 1_A
    +
    \stoptime^2 1_{A^c} .
\]

However, the following simple pathological example illustrates that we cannot have an analogous condition for the set
of IC policies for every choice of environment.

\begin{restatable}{example}{ic-pasting}
    Let $\numdataa=3$ and $\numdatab=1$. Suppose the valuations are given by
    \(
        \valbjj{0}=0
    \),
    and
    \begin{gather*}
        \valaii{0}=0
        \quad
        \valaii{1}=1
        \quad
        \valaii{2}=2
        % \begin{cases}
        %     2
        %     & 
        %     \text{w.p. } \frac{1}{2}
        %     \\
        %     0
        %     & 
        %     \text{w.p. } \frac{1}{2}
        % \end{cases}
        \quad
        \valaii{3}
        =
        % \abr{
        %     \valaii{2}-2
        % }
        0
        % \hspace{0.4cm}
        \qquad
        \valbjj{0}=0
        \quad
        \valbjj{1}=
        \begin{cases}
            1
            & 
            \text{w.p. } \frac{1}{2}
            \\
            0
            & 
            \text{w.p. } \frac{1}{2}
        \end{cases}
        % \\
    \end{gather*}
    and consider the utility/welfare functions 
    \(
        \utila(x,y)=x
    \),
    \(
        \utilb(x,y)=y
    \),
    \(
        \welfarefunc(x,y)
        =
        -x
    \).
    % \kkcomment{Reviewers will complain that this example is pathological since the welfare is decreasing with $x$. If you can redo a version with a more sensible welfare function, it would be more powerful.}
    % and the welfare function
    Now consider the following two IC policies 
    \(
        \exchangepolicy^1, \exchangepolicy^2\in\seticexchangepolicies
    \)
    given by
    \begin{align*}
        \strategytt{0}^1 
        = 
        \strategytt{0}^2 
        =
        (0,0)
        \qquad
        \strategytt{1}^1
        =
        \strategytt{1}^2
        = (1,1)
        \qquad
        \strategytt{2}^1 
        =
        \begin{cases}
            (2,1)
            &
            \text{if } \valbjj{1}=1
            \\
            (3,1)
            &
            \text{if } \valbjj{1}=0
        \end{cases}  
        \qquad
        \strategytt{2}^2
        =
        \begin{cases}
            (2,1)
            &
            \text{if } \valbjj{1}=0
            \\
            (3,1)
            &
            \text{if } \valbjj{1}=1
        \end{cases}  .
    \end{align*}  
    and $\stoptime^1=\stoptime^2=2$.
    Now consider 
    \(
        \exchangepolicy^3\in\setexchangepolicies
    \)
    given by
    \[
        \strategytt{0}^3
        = 
        (0,0)
        \qquad
        \strategytt{1}^3
        = (1,1)
        \qquad
        \strategytt{2}^3
        =
        \strategytt{2}^1
        1_{
            \cbr{
                \valbjj{1}= 0
            }
        }
        +
        \strategytt{2}^2
        1_{
            \cbr{
                \valbjj{1}= 1
            }
        }
        =
        (3,1)
    \]
    and 
    \(
        \stoptime^3
        =
        \stoptime^1
        1_{
            \cbr{
                \valbjj{1}= 0
            }
        }
        +
        \stoptime^2
        1_{
            \cbr{
                \valbjj{1}= 1
            }
        }
        = 2.
    \)
    The issue is now that using the exit strategy $\altstoptime=1$, i.e. leaving at $\timeind=1$, yields
    a higher expected welfare for $\competitora$ since
    \[
        1
        =
        \EE\sbr{
            \utilass{
                (1,1)
            }
        }
        =
        \EE\sbr{
            \utilass{
                \strategytt{
                    \stoptime^3\wedge\altstoptime
                }^3
            }
        }
        >
        \EE\sbr{
            \utilass{
                \strategytt{
                    \stoptime^3
                }^3
            }
        }
        =
        \EE\sbr{
            \utilass{
                (3,1)
            }
        }
        =0 .
    \]
    Thus $\exchangepolicy^3\not\in\seticexchangepolicies$.
\end{restatable}  


% If $\pi_1,\pi_2\in\Pi_s$ are policies starting at the same $s$, and $A\in\Fcal_s$ then
% \[
%     \alpha(\pi)=\alpha(\pi_1)1_A + \alpha(\pi_2)1_{A^c}\in\Pi_s
%     \quad
%     \tau=\tau_1 1_A + \tau_2 1_{A^c}
%     \quad
%     \sigma_t = \sigma^1_t 1_A + \sigma^2_t 1_{A^c}
% \]

In order have hope of designing an optimal IC policy, we need to make an assumption on the environment
that rules out this type of phenomena.
% \kkcomment{Consider emphasizing that such a phenomenon occurs only with pathological welfares, so the assumption is really ruling out such welfares. Is this true?}

% \mycomment{fix formatting later?}
\begin{restatable}{assumption}{ic-pasting}
    \label{assum:ic-pasting}
    Let 
    \(
        \fixedpath
        =
        \rbr{
            \exchangett{0},
            \ldots,
            \exchanget
        }
    \)
    be a fixed increasing path of points in $\setexchanges$ starting at $(0,0)$.
    Let 
    \(
        \exchangepolicy^1,\exchangepolicy^2\in\seticexchangepolicies 
    \)
    and define the policy $\exchangepolicy^3$ to be whichever of $\exchangepolicy^1$ and $\exchangepolicy^2$
    follows $\fixedpath$ for longer. In the case that both policies reach $\exchanget$ before terminating,
    continue with policy
    \(
        \argmax_{
            i\in\cbr{1,2}
        }
        \EE\sbr{
            \rewardss{
                \strategytt{
                    \stoptime^i
                }^i
            }
            |
            \statefiltss{
                \exchanget
            }
        }
    \),
    i.e.
    whichever has the higher expected reward going forward.
    We assume the environment is such that $\exchangepolicy^3\in\seticexchangepolicies$.
% \mycomment{say this is spelled out in more detail in the appendix?}
\end{restatable}  

\subsection{A Conditionally Optimal IC Algorithm}

Under the assumption from the previous section we will construct a policy whose expected welfare
is 
\(
    \sup_{
        \exchangepolicy\in\seticexchangepolicies
    }
    \EE\sbr{
        \rewardss{
            \strategytt{\stoptime}
        }
    }
\). Thus when our constructed policy is IC, it will be IC optimal.
Before providing this construction, we introduce some necessary results and notation,
analogous to those given in Section~\ref{subsec:reduction-to-optimal-stopping}.



\begin{restatable}{definition}{replace-later}
    % constrained snell env?
    Let 
    \(
        \fixedpath
        =
        \rbr{
            \exchangett{0},
            \ldots,
            \exchanget
        }
    \)
    be a fixed increasing path of points in $\setexchanges$ starting at $(0,0)$, and 
    $\rewardbound > \abr{\rewards}$ be a bound on the rewards. Define
    \begin{align*}
        \icsnellenvp
        =
        \esssup_{
            \exchangepolicy\in\seticexchangepolicies
        }
        \EE\sbr{
            1_{
                \policyfollowspath{\fixedpath}{\exchangepolicy}
            }
            \rewardss{
                \strategytt{\stoptime}
            }
            -
            1_{
                \policyfollowspath{\fixedpath}{\exchangepolicy}^c
            }
            \rewardbound
            ~
            |
            ~
            \statefiltss{\exchanget}
        }
    \end{align*}  
    where 
    \(
        \policyfollowspath{\fixedpath}{\exchangepolicy}
        =
        \cbr{
            \strategyindexrange{0}{\timeind}
            =
            \fixedpath,
            % \exchangeindexrange{0}{\timeind},
            \timeind\leq \stoptime
        }
    \)
    is the event that $\exchangepolicy$ follows $\fixedpath$ prior to stopping and
    \(
        \strategyindexrange{0}{\timeind}
    \) is shorthand for 
    \(
        \rbr{
            \strategytt{0},\ldots,\strategytt{\timeind} 
        }
    \).
\end{restatable}  

Intuitively, each $\icsnellenvp$ represents the best expected reward for policies that have followed $\fixedpath$, 
given the information available at the end of the path. 
We will see that for $\icsnellenvp$ there is an analogous dynamic programming equation to
the one in
% In order to prove an analogue to 
Theorem~\ref{thm:osas-dp-eq}.
This result will be needed for our construction. 
However, before proving it, we first need the following lemma.
% we introduce the following dynamic programming equation.

\begin{restatable}{lemma}{directed-upwards}
    \label{lem:directed-upwards}
    Suppose that Assumption~\ref{assum:ic-pasting} holds
    and let 
    \(
        \fixedpath
        =
        \rbr{
            \exchangett{0},
            \ldots,
            \exchanget
        }
    \)
    be a fixed increasing path of points in $\setexchanges$ starting at $(0,0)$.
    The family of random variables 
    \begin{align*}
        \cbr{
            1_{
                \policyfollowspath{\fixedpath}{\exchangepolicy}
            }
            \EE\sbr{
                \rewardss{
                    \strategytt{\stoptime}
                }
                | \statefiltss{\exchanget}
            }
            -
            1_{
                \policyfollowspath{\fixedpath}{\exchangepolicy}^c
                \rewardbound
            }
            :
            \exchangepolicy\in\seticexchangepolicies
        }
    \end{align*}  
    is closed under maxima. 
\end{restatable}  
\begin{proof}
    Our goal is to show that given any 
    \(
        \exchangepolicy_1,\exchangepolicy_2\in\seticexchangepolicies
    \),
    there exists
    \(
        \exchangepolicy_3\in\seticexchangepolicies
    \)
    such that
    \begin{gather*}
        \rbr{
            1_{
                \policyfollowspath{\fixedpath}{\exchangepolicy_1}
            }
            \EE\sbr{
                \rewardss{
                    \stopstateforpolicy{
                        \exchangepolicy_1
                    }
                }
                | \statefiltss{\exchanget}
            }
            -
            1_{
                \policyfollowspath{\fixedpath}{\exchangepolicy_1}^c
            }
            \rewardbound
        }
        \vee
        \rbr{
            1_{
                \policyfollowspath{\fixedpath}{\exchangepolicy_2}
            }
            \EE\sbr{
                \rewardss{
                    \stopstateforpolicy{
                        \exchangepolicy_2
                    }
                }
                | \statefiltss{\exchanget}
            }
            -
            1_{
                \policyfollowspath{\fixedpath}{\exchangepolicy_2}^c
            }
            \rewardbound
        }
        \\
        =
        \rbr{
            1_{
                \policyfollowspath{\fixedpath}{\exchangepolicy_3}
            }
            \EE\sbr{
                \rewardss{
                    \stopstateforpolicy{
                        \exchangepolicy_3
                    }
                }
                | \statefiltss{\exchanget}
            }
            -
            1_{
                \policyfollowspath{\fixedpath}{\exchangepolicy_3}^c
            }
            \rewardbound
        } .
        \numberthis
        \label{eq:closed-under-maxima}
    \end{gather*}  
    We will define
    \(
        \exchangepolicy_3
        =
        \rbr{
            \cbr{
                \strategyt^3
            }_{\timeind\geq 0},
            \stoptime^3
        }
    \)
    to be whichever of policy 
    \(
        \exchangepolicy_1
        =
        \rbr{
            \cbr{
                \strategyt^1
            }_{\timeind\geq 0},
            \stoptime^1
        }
    \)
    or
    \(
        \exchangepolicy_2
        =
        \rbr{
            \cbr{
                \strategyt^2
            }_{\timeind\geq 0},
            \stoptime^2
        }
    \)
    adheres to $\fixedpath$ for longer. In the case that 
    \(
        \strategyindexrange{0}{\timeind}^1
        =
        \strategyindexrange{0}{\timeind}^2
        =
        \fixedpath
    \)
    and $\stoptime^1,\stoptime^2\geq\timeind$, we select the policy
    given by 
    \(
        \argmax_{
            i\in\cbr{1,2}
        }
        \EE\sbr{
            \rewardss{
                \stopstateforpolicy{\pi_i}
            }
            | \statefiltss{\exchanget}
        }
    \). 
    In otherwords, if both policies follow $\fixedpath$, then we 
    take the one which gives a better outlook for the reward. In the 
    case that both policies do not complete following $\fixedpath$ at the
    same time, we will set 
    \(
        \exchangepolicy_3
        =
        \exchangepolicy_1
    \)
    without loss of generality.

    Formally, let 
    \begin{align*}
        \lastfollowtimei
        =
        \inf\cbr{
            k\in\cbr{
                0,\ldots,\timeind-1
            }
            :
            \strategytt{k+1}^i
            \neq
            \exchangett{k+1}
        }
        \wedge
        \stoptime^i
    \end{align*}  
    be the last time at which $\exchangepolicy_i$ is following $\fixedpath$
    (here we let $\inf\varnothing=\timeind$). 
    Thus 
    \(
        \lastcommonfollowtime
        =
        \lastfollowtimeii{1}
        \wedge
        \lastfollowtimeii{2}
    \)
    is the last time at which both policies are following $\fixedpath$.
    Now let 
    \begin{align*}
        \chosenpath
        =
        \begin{cases}
            1
            & 
            \text{if } \lastfollowtimeii{1} > \lastfollowtimeii{2}
            \\
            2
            & 
            \text{if } \lastfollowtimeii{1} < \lastfollowtimeii{2}
            \\
            \argmax_{
                i\in\cbr{1,2}
            }
            \EE\sbr{
                \rewardss{
                    \stopstateforpolicy{\pi_i}
                }
                | \statefiltss{\exchanget}
            }
            & 
            \text{if } \lastfollowtimeii{1} = \lastfollowtimeii{2} = \timeind
            \\
            1
            & 
            \text{if } \lastfollowtimeii{1} = \lastfollowtimeii{2} <\timeind
        \end{cases}  
    \end{align*}  
    be which policy is selected and define $\exchangepolicy_3$ via
    the strategy
    \begin{align*}
        \strategytt{k+1}^3
        &=
        1_{
            \cbr{
                k<\lastcommonfollowtime
            }
        }
        \exchangett{k+1}
        +
        1_{
            \cbr{
                k\geq\lastcommonfollowtime
            }
        }
        \strategytt{k+1}^\chosenpath
        \\
        &=
        1_{
            \cbr{
                \chosenpath=1
            }
        }
        \strategytt{k+1}^1
        +
        1_{
            \cbr{
                \chosenpath=2
            }
        }
        \strategytt{k+1}^2
    \end{align*}  
    and the stopping time 
    \(
        \stoptime^3
        =
        1_{
            \cbr{
                \chosenpath=1
            }
        }
        \stoptime^1
        +
        1_{
            \cbr{
                \chosenpath=2
            }
        }
        \stoptime^2
    \). 
    Intuitively, $\exchangepolicy_3$ follows $\fixedpath$ up until
    $\lastcommonfollowtime$, at which point we can determine which of 
    $\exchangepolicy_1$ or $\exchangepolicy_2$ to follow.
    From the construction above, we have that 
    \(
        \cbr{
            \chosenpath=1
        }
        =
        \cbr{
            \exchangepolicy_3
            =
            \exchangepolicy_1
        }
    \)
    and 
    \(
        \cbr{
            \chosenpath=2
        }
        =
        \cbr{
            \exchangepolicy_3
            =
            \exchangepolicy_2
        }
    \).

    For brevity let 
    \(
        A=
        \policyfollowspath{\fixedpath}{\exchangepolicy_1}
    \)
    and 
    \(
        B=
        \policyfollowspath{\fixedpath}{\exchangepolicy_2}
    \).
    Since $\exchangepolicy_3$ fails to follow $\fixedpath$ exactly when
    both $\exchangepolicy_1$ and $\exchangepolicy_2$ fail to follow 
    $\fixedpath$, we can write~\eqref{eq:closed-under-maxima} as 
    \begin{align*}
        1_{
            \policyfollowspath{\fixedpath}{\exchangepolicy_3}
        }
        \EE\sbr{
            \rewardss{
                \stopstateforpolicy{
                    \exchangepolicy_3
                }
            }
            | \statefiltss{\exchanget}
        }
        -
        1_{
            \policyfollowspath{\fixedpath}{\exchangepolicy_3}^c
        }
        \rewardbound
        =
        \rbr{
            1_{
                A\cap B
            }
            +
            1_{
                A\cap B^c
            }
            +
            1_{
                A^c\cap B
            }
        }
        \EE\sbr{
            \rewardss{
                \stopstateforpolicy{
                    \exchangepolicy_3
                }
            }
            | \statefiltss{\exchanget}
        }
        -
        1_{
            A^c\cap B^c
        }
        \rewardbound .
    \end{align*}   
    By construction of $\exchangepolicy_3$, the right hand side becomes
    \begin{gather*}
        1_{
            A\cap B
        }
        \rbr{
            \EE\sbr{
                \rewardss{
                    \stopstateforpolicy{
                        \exchangepolicy_1
                    }
                }
                | \statefiltss{\exchanget}
            }
            \vee
            \EE\sbr{
                \rewardss{
                    \stopstateforpolicy{
                        \exchangepolicy_2
                    }
                }
                | \statefiltss{\exchanget}
            }
        }
        +
        1_{
            A\cap B^c
        }
        \EE\sbr{
            \rewardss{
                \stopstateforpolicy{
                    \exchangepolicy_1
                }
            }
            | \statefiltss{\exchanget}
        }
        +
        1_{
            A^c\cap B
        }
        \EE\sbr{
            \rewardss{
                \stopstateforpolicy{
                    \exchangepolicy_2
                }
            }
            | \statefiltss{\exchanget}
        }
        -
        1_{
            A^c\cap B^c
        }
        \rewardbound
        \\
        \rbr{
            1_{
                \policyfollowspath{\fixedpath}{\exchangepolicy_1}
            }
            \EE\sbr{
                \rewardss{
                    \stopstateforpolicy{
                        \exchangepolicy_1
                    }
                }
                | \statefiltss{\exchanget}
            }
            -
            1_{
                \policyfollowspath{\fixedpath}{\exchangepolicy_1}^c
            }
            \rewardbound
        }
        \vee
        \rbr{
            1_{
                \policyfollowspath{\fixedpath}{\exchangepolicy_2}
            }
            \EE\sbr{
                \rewardss{
                    \stopstateforpolicy{
                        \exchangepolicy_2
                    }
                }
                | \statefiltss{\exchanget}
            }
            -
            1_{
                \policyfollowspath{\fixedpath}{\exchangepolicy_2}^c
            }
            \rewardbound
        } 
    \end{gather*}  
    as desired. Finally, by Assumption~\ref{assum:ic-pasting}
    we have $\exchangepolicy\in\seticexchangepolicies$ which concludes the 
    proof.
\end{proof}  

The point of Lemma~\ref{lem:directed-upwards} is to show that under Assumption~\ref{assum:ic-pasting}, 
two policies can be pasted together to form a new policy which either follows the path for longer
or has a higher expected reward.
We now use this to prove the following theorem which our constructed policy will be based on.

\begin{restatable}{theorem}{icDpEq}
    \label{thm:ic-dp-eq}
    Suppose that Assumption~\ref{assum:ic-pasting} holds
    and $\seticexchangepolicies$ is countable. Let 
    \(
        \existspolicymatchingpath{\fixedpath}
        =
        \cbr{
            \exists \exchangepolicy\in\seticexchangepolicies :
            \strategyindexrange{0}{\stoptime}=\fixedpath
        }
    \)
    be the event that there is an IC exchange policy which follows $\fixedpath$ and then stops. Then,
    \[
        \icsnellenvp
        =
        \max\cbr{
            \rbr{
                1_{
                    \existspolicymatchingpath{\fixedpath}
                }
                \rewardss{\exchanget}
                -
                1_{
                    \existspolicymatchingpath{\fixedpath}^c
                }
                \rewardbound
            },
            \EE\sbr{
                \icsnellenvpp{
                    \rbr{
                        \fixedpath, \nextexchange
                    }
                }
                |
                \statefiltss{\exchanget}
            }
            \nextexchange
            > \exchanget 
            % \in
            % \validicextensionspp{\fixedpath}
        }
    \]
    where 
    \(
        \rbr{
            \fixedpath, \nextexchange
        }
    \)
    represents the path $\fixedpath$ extended to $\nextexchange$.
    % , and 
    % \(
    %     \validicextensionspp{\fixedpath}
    %     =
    %     \cbr{
    %         \nextexchange\in\setexchanges :
    %         \nextexchange > \exchanget,
    %         \exists 
    %     }
    % \)
    % is the set of
\end{restatable}  

% We now give a sketch of this proof with the full details available in Appendix~\ref{app:omitted-proofs}.
% At a high level, this result says that the 
\begin{proof}
    % [Proof sketch]
    Fix
    \(
        \fixedpath
        =
        \rbr{
            \exchangett{0},
            \ldots,
            \exchanget
        }
    \).
    % We will follow the same high level strategy that~\citep{mandelbaum1981optimal}
    % uses to prove Theorem~\ref{thm:osas-dp-eq}, showing the left hand side is both greater than or equal and less than
    % or equal to the right hand side.
    % KK: I took this off since it seemed to minimize your contributons. If you want to put it back in, then highlight how your proof differs from theirs. Ideally you would do this if you had time.
    We start by proving
    \begin{align*}
        \icsnellenvp
        \geq
        \max\cbr{
            \rbr{
                1_{
                    \existspolicymatchingpath{\fixedpath}
                }
                \rewardss{\exchanget}
                -
                1_{
                    \existspolicymatchingpath{\fixedpath}^c
                }
                \rewardbound
            },
            \EE\sbr{
                \icsnellenvpp{
                    \rbr{
                        \fixedpath, \nextexchange
                    }
                }
                |
                \statefiltss{\exchanget}
            }
            \nextexchange
            > \exchanget 
        } .
    \end{align*}  
    Let 
    \(
        \exchanget<\nextexchange\in\setexchanges
    \)
    and consider the path $\fixedpath$ extended to $\nextexchange$, 
    \(
        \rbr{
            \fixedpath,\nextexchange
        }
    \).
    By Lemma~\ref{lem:directed-upwards} the family of random variables
    \begin{align*}
        \cbr{
            1_{
                \policyfollowspath{
                    \fixedpath,\nextexchange
                }{\exchangepolicy}
            }
            \EE\sbr{
                \rewardss{
                    \strategytt{\stoptime}
                }
                |
                \statefiltss{\nextexchange}
            }
            -
            1_{
                \policyfollowspath{
                    \fixedpath,\nextexchange
                }{\exchangepolicy}^c
            }
            \rewardbound 
            :
            \exchangepolicy\in\seticexchangepolicies
        }
    \end{align*}  
    is directed upwards.
    Therefore, there exists a sequence of exchange policies 
    \(
        \exchangepolicy_k 
        \in\seticexchangepolicies
    \)
    such that 
    \begin{align*}
        \snellenvss{
            \rbr{
                \fixedpath,\nextexchange
            }
        }
        =
        \lim_{k\rightarrow\infty}
        \rbr{
            1_{
                \policyfollowspath{
                    \fixedpath,\nextexchange
                }{\exchangepolicy_k}
            }
            \EE\sbr{
                \rewardss{
                    \strategytt{\stoptime}
                }
                |
                \statefiltss{\nextexchange}
            }
            -
            1_{
                \policyfollowspath{
                    \fixedpath,\nextexchange
                }{\exchangepolicy_k}^c
            }
            \rewardbound 
        } .
    \end{align*}  
    Since $\rewards$ is bounded, we can apply the dominated convergence theorem for conditional expectations, which
    together with the tower property, yields
    \begin{align*}
        \EE\sbr{
            \snellenvss{
                \rbr{
                    \fixedpath,\nextexchange
                }
            }
            | 
            \statefiltss{\exchanget}
        }
        &=
        \EE\sbr{
            \lim_{k\rightarrow\infty}
            \rbr{
                1_{
                    \policyfollowspath{
                        \fixedpath,\nextexchange
                    }{\exchangepolicy_k}
                }
                \EE\sbr{
                    \rewardss{
                        \stopstateforpolicy{
                            \exchangepolicy_k
                        }
                    }
                    |
                    \statefiltss{\nextexchange}
                }
                -
                1_{
                    \policyfollowspath{
                        \fixedpath,\nextexchange
                    }{\exchangepolicy_k}^c
                }
                \rewardbound 
            }
            ~|~ 
            \statefiltss{\exchanget}
        }
        \\
        &=
        \lim_{k\rightarrow\infty}
        \rbr{
            1_{
                \policyfollowspath{
                    \fixedpath,\nextexchange
                }{\exchangepolicy_k}
            }
            \EE\sbr{
                \rewardss{
                    \stopstateforpolicy{
                        \exchangepolicy_k
                    }
                }
                | 
                \statefiltss{\exchanget}
            }
            -
            1_{
                \policyfollowspath{
                    \fixedpath,\nextexchange
                }{\exchangepolicy_k}^c
            }
            \rewardbound 
        } .
    \end{align*}  
    Now notice that 
    \(
        \policyfollowspath{
            \fixedpath,\nextexchange
        }{\exchangepolicy}
        \subseteq
        \policyfollowspath{
            \fixedpath
        }{\exchangepolicy}
    \)
    implies
    \(
        \policyfollowspath{
            \fixedpath,\nextexchange
        }{\exchangepolicy}^c
        \supseteq
        \policyfollowspath{
            \fixedpath
        }{\exchangepolicy}^c
    \).
    This, together with the definiton of $\icsnellenvp$, gives
    \begin{align*}
        \lim_{k\rightarrow\infty}
        \rbr{
            1_{
                \policyfollowspath{
                    \fixedpath,\nextexchange
                }{\exchangepolicy_k}
            }
            \EE\sbr{
                \rewardss{
                    \stopstateforpolicy{
                        \exchangepolicy_k
                    }
                }
                | 
                \statefiltss{\exchanget}
            }
            -
            1_{
                \policyfollowspath{
                    \fixedpath,\nextexchange
                }{\exchangepolicy_k}^c
            }
            \rewardbound 
        }
        % \\
        &\leq
        % \leq
        \lim_{k\rightarrow\infty}
        \rbr{
            1_{
                \policyfollowspath{
                    \fixedpath
                }{\exchangepolicy_k}
            }
            \EE\sbr{
                \rewardss{
                    \stopstateforpolicy{
                        \exchangepolicy_k
                    }
                }
                | 
                \statefiltss{\exchanget}
            }
            -
            1_{
                \policyfollowspath{
                    \fixedpath
                }{\exchangepolicy_k}^c
            }
            \rewardbound 
        }
        \\
        &\leq
        % \leq
        \lim_{k\rightarrow\infty}
        \icsnellenvp
        =
        \icsnellenvp .
    \end{align*}  
    Thus, we have show for any $\nextexchange>\exchanget$,
    \(
        \icsnellenvp
        \geq
        \EE\sbr{
            \snellenvss{
                \rbr{
                    \fixedpath,\nextexchange
                }
            }
            | 
            \statefiltss{\exchanget}
        }
    \).
    To show 
    \(
        \icsnellenvp
        \geq
        1_{
            \existspolicymatchingpath{\fixedpath}
        }
        \rewardss{
            \exchanget
        }
        -
        1_{
            \existspolicymatchingpath{\fixedpath}^c
        }
        \rewardbound
    \), 
    it sufficies to prove
    \(
        1_{
            \existspolicymatchingpath{\fixedpath}
        }
        \icsnellenvp
        \geq
        1_{
            \existspolicymatchingpath{\fixedpath}
        }
        \rewardss{
            \exchanget
        }
    \) 
    since 
    $-\rewardbound$ lower bounds $\icsnellenvp$.
    By the definition, on the event
    $\existspolicymatchingpath{\fixedpath}$ there is a policy in $\seticexchangepolicies$ matching $\fixedpath$.
    Since $\seticexchangepolicies$, the essential supremum can be replaced with a supremum. Therefore,
    \begin{align*}
        1_{
            \existspolicymatchingpath{\fixedpath}
        }
        \icsnellenvp
        &=
        1_{
            \existspolicymatchingpath{\fixedpath}
        }
        \esssup_{
            \exchangepolicy\in\seticexchangepolicies
        }
        \EE\sbr{
            1_{
                \policyfollowspath{\fixedpath}{\exchangepolicy}
            }
            \rewardss{
                \strategytt{\stoptime}
            }
            -
            1_{
                \policyfollowspath{\fixedpath}{\exchangepolicy}^c
            }
            \rewardbound
            ~
            |
            ~
            \statefiltss{\exchanget}
        }
        \\
        &=
        1_{
            \existspolicymatchingpath{\fixedpath}
        }
        \sup_{
            \exchangepolicy\in\seticexchangepolicies
        }
        \EE\sbr{
            1_{
                \policyfollowspath{\fixedpath}{\exchangepolicy}
            }
            \rewardss{
                \strategytt{\stoptime}
            }
            -
            1_{
                \policyfollowspath{\fixedpath}{\exchangepolicy}^c
            }
            \rewardbound
            ~
            |
            ~
            \statefiltss{\exchanget}
        }
        \\
        &\geq
        1_{
            \existspolicymatchingpath{\fixedpath}
        }
        \rewardss{
            \strategyt
        } .
    \end{align*} 
    % \mycomment{maybe go into a bit more detail on this, its correct but maybe not extremely clear}
    Putting everything together we conclude
    \begin{align*}
        \icsnellenvp
        \geq
        \max\cbr{
            \rbr{
                1_{
                    \existspolicymatchingpath{\fixedpath}
                }
                \rewardss{\exchanget}
                -
                1_{
                    \existspolicymatchingpath{\fixedpath}^c
                }
                \rewardbound
            },
            \EE\sbr{
                \icsnellenvpp{
                    \rbr{
                        \fixedpath, \nextexchange
                    }
                }
                |
                \statefiltss{\exchanget}
            }
            \nextexchange
            > \exchanget 
        } .
    \end{align*}  
    For the other inequality we will show that for any $\exchangepolicy\in\seticexchangepolicies$,
    \begin{align}
        1_{
            \policyfollowspath{\fixedpath}{\exchangepolicy}
        }
        \EE\sbr{
            \rewardss{
                \strategytt{\stoptime}
            }
            | \statefiltss{\strategyt}
        }
        -
        1_{
            \policyfollowspath{\fixedpath}{\exchangepolicy}^c
        }
        \rewardbound
        \leq 
        \max\cbr{
            \rbr{
                1_{
                    \existspolicymatchingpath{\fixedpath}
                }
                \rewardss{\exchanget}
                -
                1_{
                    \existspolicymatchingpath{\fixedpath}^c
                }
                \rewardbound
            },
            \EE\sbr{
                \icsnellenvpp{
                    \rbr{
                        \fixedpath, \nextexchange
                    }
                }
                |
                \statefiltss{\exchanget}
            }
            \nextexchange
            > \exchanget 
        }
        \label{eq:ic-dp-proof}
    \end{align}  
    from which, by taking the essential supremum over both sides, it will follow that 
    \begin{align*}
        \icsnellenvp
        \leq
        \max\cbr{
            \rbr{
                1_{
                    \existspolicymatchingpath{\fixedpath}
                }
                \rewardss{\exchanget}
                -
                1_{
                    \existspolicymatchingpath{\fixedpath}^c
                }
                \rewardbound
            },
            \EE\sbr{
                \icsnellenvpp{
                    \rbr{
                        \fixedpath, \nextexchange
                    }
                }
                |
                \statefiltss{\exchanget}
            }
            \nextexchange
            > \exchanget 
        } .
    \end{align*}  
    To prove~\eqref{eq:ic-dp-proof}, it suffices to show that
    \begin{align*}
        1_{
            \policyfollowspath{\fixedpath}{\exchangepolicy}
        }
        \EE\sbr{
            \rewardss{
                \strategytt{\stoptime}
            }
            | \statefiltss{\exchanget}
        }
        \leq 
        1_{
            \policyfollowspath{\fixedpath}{\exchangepolicy}
        }
        \max\cbr{
            \rbr{
                1_{
                    \existspolicymatchingpath{\fixedpath}
                }
                \rewardss{\exchanget}
                -
                1_{
                    \existspolicymatchingpath{\fixedpath}^c
                }
                \rewardbound
            },
            \EE\sbr{
                \icsnellenvpp{
                    \rbr{
                        \fixedpath, \nextexchange
                    }
                }
                |
                \statefiltss{\exchanget}
            }
            \nextexchange
            > \exchanget 
        }
    \end{align*}  
    since $-\rewardbound$ is a lower bound for $\rewards$. Notice that on the event
    $\policyfollowspath{\fixedpath}{\exchangepolicy}$, $\exchangepolicy$ either stops at time 
    $\timeind$ or continues, moving to a state $\nextexchange>\exchanget$. Thus
    \begin{align*}
        1_{
            \policyfollowspath{\fixedpath}{\exchangepolicy}
        }
        \EE\sbr{
            \rewardss{
                \strategytt{\stoptime}
            }
            | \statefiltss{\exchanget}
        }
        =
        1_{
            \policyfollowspath{\fixedpath}{\exchangepolicy}
        }
        \rbr{
            1_{
                \cbr{\stoptime=\timeind}
            }
            \rewardss{\exchanget}
            +
            \sum_{
                \nextexchange>\exchanget
            }
            1_{
                \cbr{
                    \substack{
                        \strategytt{\timeind+1}=\nextexchange,
                        \\
                        \stoptime\geq\timeind+1
                    }
                }
            }
            \EE\sbr{
                \rewardss{
                    \strategytt{\stoptime}
                }
                | \statefiltss{\exchanget}
            }
        } .
    \end{align*}  
    Because 
    \(
        \policyfollowspath{\fixedpath}{\exchangepolicy}
        \cap
        \cbr{
            \stoptime=\timeind
        }
        \subseteq
        \existspolicymatchingpath{\fixedpath}
    \),
    we have
    \(
        \policyfollowspath{\fixedpath}{\exchangepolicy}
        \cap
        \cbr{
            \stoptime=\timeind
        }
        \cap
        \existspolicymatchingpath{\fixedpath}^c
        =
        \varnothing
    \). Therefore,
    the first term can be rewritten as 
    \begin{align*}
        1_{
            \policyfollowspath{\fixedpath}{\exchangepolicy}
        }
        1_{
            \cbr{\stoptime=\timeind}
        }
        \rewardss{\exchanget}
        =
        1_{
            \policyfollowspath{\fixedpath}{\exchangepolicy}
        }
        1_{
            \cbr{\stoptime=\timeind}
        }
        \rbr{
            1_{
                \existspolicymatchingpath{\fixedpath}
            }
            \rewardss{\exchanget}
            -
            1_{
                \existspolicymatchingpath{\fixedpath}^c
            }
            \rewardbound
        } .
    \end{align*}  
    For the second term,
    \(
        \policyfollowspath{\fixedpath}{\exchangepolicy}
        \cap
        \cbr{
            \substack{
                \strategytt{\timeind+1}=\nextexchange,
                \\
                \stoptime\geq\timeind+1
            }
        }
        =
        \policyfollowspath{\fixedpath,\nextexchange}{\exchangepolicy}
    \)
    and so
    \begin{align*}
        1_{
            \policyfollowspath{\fixedpath}{\exchangepolicy}
        }
        \sum_{
            \nextexchange>\exchanget
        }
        1_{
            \cbr{
                \substack{
                    \strategytt{\timeind+1}=\nextexchange,
                    \\
                    \stoptime\geq\timeind+1
                }
            }
        }
        \EE\sbr{
            \rewardss{
                \strategytt{\stoptime}
            }
            | \statefiltss{\exchanget}
        }
        &=
        1_{
            \policyfollowspath{\fixedpath}{\exchangepolicy}
        }
        \sum_{
            \nextexchange>\exchanget
        }
        1_{
            \cbr{
                \substack{
                    \strategytt{\timeind+1}=\nextexchange,
                    \\
                    \stoptime\geq\timeind+1
                }
            }
        }
        \EE\sbr{
            \EE\sbr{
                \rewardss{
                    \strategytt{\stoptime}
                }
                | \statefiltss{\nextexchange}
            }
            | \statefiltss{\exchanget}
        }
        \\
        &=
        1_{
            \policyfollowspath{\fixedpath}{\exchangepolicy}
        }
        \sum_{
            \nextexchange>\exchanget
        }
        1_{
            \cbr{
                \substack{
                    \strategytt{\timeind+1}=\nextexchange,
                    \\
                    \stoptime\geq\timeind+1
                }
            }
        }
        \EE\sbr{
            1_{
                \policyfollowspath{\fixedpath,\nextexchange}{\exchangepolicy}
            }
            \EE\sbr{
                \rewardss{
                    \strategytt{\stoptime}
                }
                | \statefiltss{\nextexchange}
            }
            -
            1_{
                \policyfollowspath{\fixedpath,\nextexchange}{\exchangepolicy}^c
            }
            \rewardbound
            ~|~ \statefiltss{\exchanget}
        }
        \\
        &\leq
        1_{
            \policyfollowspath{\fixedpath}{\exchangepolicy}
        }
        \sum_{
            \nextexchange>\exchanget
        }
        1_{
            \cbr{
                \substack{
                    \strategytt{\timeind+1}=\nextexchange,
                    \\
                    \stoptime\geq\timeind+1
                }
            }
        }
        \EE\sbr{
            \icsnellenvpp{
                \fixedpath,\nextexchange
            }
            | \statefiltss{\exchanget}
        } .
    \end{align*}  
    Putting together both pieces gives
    \begin{align*}
        &\hspace{0.4cm}
        1_{
            \policyfollowspath{\fixedpath}{\exchangepolicy}
        }
        \rbr{
            1_{
                \cbr{\stoptime=\timeind}
            }
            \rewardss{\exchanget}
            +
            \sum_{
                \nextexchange>\exchanget
            }
            1_{
                \cbr{
                    \substack{
                        \strategytt{\timeind+1}=\nextexchange,
                        \\
                        \stoptime\geq\timeind+1
                    }
                }
            }
            \EE\sbr{
                \rewardss{
                    \strategytt{\stoptime}
                }
                | \statefiltss{\exchanget}
            }
        }
        \\
        &=
        1_{
            \policyfollowspath{\fixedpath}{\exchangepolicy}
        }
        \rbr{
            1_{
                \cbr{\stoptime=\timeind}
            }
            \rbr{
                1_{
                    \existspolicymatchingpath{\fixedpath}
                }
                \rewardss{\exchanget}
                -
                1_{
                    \existspolicymatchingpath{\fixedpath}^c
                }
                \rewardbound
            }        
            +
            \sum_{
                \nextexchange>\exchanget
            }
            1_{
                \cbr{
                    \substack{
                        \strategytt{\timeind+1}=\nextexchange,
                        \\
                        \stoptime\geq\timeind+1
                    }
                }
            }
            \EE\sbr{
                \rewardss{
                    \strategytt{\stoptime}
                }
                | \statefiltss{\exchanget}
            }
        }
        \\
        &\leq
        1_{
            \policyfollowspath{\fixedpath}{\exchangepolicy}
        }
        \max\cbr{
            \rbr{
                1_{
                    \existspolicymatchingpath{\fixedpath}
                }
                \rewardss{\exchanget}
                -
                1_{
                    \existspolicymatchingpath{\fixedpath}^c
                }
                \rewardbound
            },
            \EE\sbr{
                \icsnellenvpp{
                    \rbr{
                        \fixedpath, \nextexchange
                    }
                }
                |
                \statefiltss{\exchanget}
            }
            \nextexchange
            > \exchanget 
        }
    \end{align*}  
    which concludes the proof.
\end{proof}  

We are now ready to state the main IC result which is derived from Theorem~\ref{thm:ic-dp-eq}.
% \kkcomment{You probably need to explain what it means for the set of IC policies to be countable. Give some examples where this might be true.}

\begin{restatable}{theorem}{policyMatchingIcOptReward}
    \label{thm:policy-matching-ic-opt-reward}
    Suppose that Assumption~\ref{assum:ic-pasting} holds and $\seticexchangepolicies$ is countable.
    Let 
    \(
        \expolicymatchicopt
        =
        \rbr{
            \cbr{
                \strategyt
            }_{\timeind\geq 0},
            \stoptime
        }
    \)
    be the exchange policy given by
    \begin{align*}
        \strategytt{0}=(0,0),
        \quad
        \strategytt{\timeind+1}
        =
        \argmax_{
            \nextexchange>\exchange
        }
        \EE\sbr{
            \icsnellenvpp{
                \rbr{
                    \strategyindexrange{0}{\timeind},\nextexchange
                }
            }
            | \timefiltt
        }
    \end{align*}  
    and
    \(
        \stoptime
        =
        \inf\cbr{
            \timeind\in\NN_{\geq 0} :
            \icsnellenvpp{
                \strategyindexrange{0}{\timeind}
            }
            =
            \rbr{
                1_{
                    \existspolicymatchingpath{\fixedpath}
                }
                \rewardss{\exchanget}
                -
                1_{
                    \existspolicymatchingpath{\fixedpath}^c
                }
                \rewardbound
            }
        }
    \).
    % then ...
    If 
    $\expolicymatchicopt\in\seticexchangepolicies$,
    then $\expolicymatchicopt$ is IC optimal. 
    % Then,
    % \(
    %     \EE\sbr{
    %         \rewardss{
    %             \strategytt{
    %                 \stoptime
    %             }
    %         }
    %     }
    %     =
    %     \sup_{
    %         \altexchangepolicy\in\seticexchangepolicies
    %     }
    %     \EE\sbr{
    %         \rewardss{
    %             \altstrategytt{
    %                 \altstoptime
    %             }
    %         }
    %     }
    % \).
\end{restatable}  
\begin{proof}
    By the definition $\expolicymatchicopt$
    and Theorem~\ref{thm:ic-dp-eq},
    we have that on $\cbr{\timeind <\stoptime}$,
    \(
        \icsnellenvpp{
            \strategyindexrange{0}{\timeind}
        }
        =
        \EE\sbr{
            \icsnellenvpp{
                \strategyindexrange{0}{\timeind+1}
            }
            | \timefiltt
        }
    \). 
    This tells us that the sequence
    \(
        M_\timeind
        =
        \icsnellenvpp{
            \strategyindexrange{0}{\timeind\wedge\stoptime}
        }
    \)
    is a martingale. Therefore,
    \(
        \EE\sbr{
            \icsnellenvpp{
                \strategytt{0}
            }
        }
        =
        \EE\sbr{M_0}
        =
        \EE\sbr{M_\timeind}
    \).
    From Theorem~\ref{thm:ic-dp-eq}, the fact that $\strategyt$ is increasing, and $\setexchanges$ is finite
    we know that $\stoptime$ is finite. Thus applying the Dominated Convergence Theorem gives 
    \[
        \EE\sbr{
            \icsnellenvpp{
                \strategytt{0}
            }
        }
        =
        \lim_{\timeind\rightarrow\infty}
        \EE\sbr{
            M_\timeind
        }
        =
        \EE\sbr{
            \lim_{\timeind\rightarrow\infty}
            M_\timeind
        }
        =
        \EE\sbr{
            \icsnellenvpp{
                \strategyindexrange{0}{\stoptime}
            }
        }
    \]
    % \mycomment{check}
    % \mycomment{bit odd state notation below}
    Since $\expolicymatchicopt\in\seticexchangepolicies$,
    by Theorem~\ref{thm:equiv-alg-defs} and the definition of $\stoptime$,
    \(
        \icsnellenvpp{
            \strategyindexrange{0}{\stoptime}
        }
        =
        1_{
            \existspolicymatchingpath{
                \strategyindexrange{0}{\stoptime}
            }
        }
        \rewardss{\stoptime}
        -
        1_{
            \existspolicymatchingpath{
                \strategyindexrange{0}{\stoptime}
            }^c
        }
        \rewardbound
        =
        \rewardss{\stoptime}
    \).
    Therefore, by the definitions of $\icsnellenvp$ and $\seticexchangepolicies$,
    \begin{align*}
        \EE\sbr{
            \icsnellenvpp{\strategytt{0}}
        }
        &=
        \EE\sbr{
            \esssup_{
                \exchangepolicy\in\seticexchangepolicies
            }
            \EE\sbr{
                1_{
                    \policyfollowspath{\strategytt{0}}{\exchangepolicy}
                }
                \rewardss{
                    \strategytt{\stoptime}
                }
                -
                1_{
                    \policyfollowspath{\strategytt{0}}{\exchangepolicy}^c
                }
                \rewardbound
                ~
                |
                ~
                \statefiltss{\exchangett{0}}
            }
        }
        \\
        &=
        \EE\sbr{
            \esssup_{
                \exchangepolicy\in\seticexchangepolicies
            }
            \EE\sbr{
                \rewardss{
                    \strategytt{\stoptime}
                }
            }
        }
        \\
        &=
        \EE\sbr{
            \sup_{
                \exchangepolicy\in\seticexchangepolicies
            }
            \EE\sbr{
                \rewardss{
                    \strategytt{\stoptime}
                }
            }
        }
        \\
        &=
        \sup_{
            \exchangepolicy\in\seticexchangepolicies
        }
        \EE\sbr{
            \rewardss{
                \strategytt{\stoptime}
            }
        } .
    \end{align*}  
    Therefore,
    \(
        \EE\sbr{
            \rewardss{
                \strategytt{\stoptime}
            }
        }
        =
        \EE\sbr{
            \icsnellenvpp{
                \strategyindexrange{0}{\stoptime}
            }
        }
        =
        \EE\sbr{
            \icsnellenvpp{
                \strategytt{0}
            }
        }
        =
        \sup_{
            \exchangepolicy\in\seticexchangepolicies
        }
        \EE\sbr{
            \rewardss{
                \strategytt{\stoptime}
            }
        }
    \), so $\expolicymatchicopt$ is IC optimal.
\end{proof}  

For technical reasons, Theorem~\ref{thm:policy-matching-ic-opt-reward} requires that 
$\seticexchangepolicies$ is countable.
However, we wish to emphasize that there a number of models in which this is the case.
For example, when the valuations $\cbr{\valai}_{\numdatafora=0}^\numdataa$,  
$\cbr{\valbj}_{\numdataforb=0}^\numdatab$ have finite support, there are finitely many policies
and thus finitely many IC policies.

% \mycomment{what about external randomness?}

% \begin{restatable}{corollary}{ic-opt-alg}
%     Suppose that Assumption~\ref{assum:ic-pasting} holds.
%     Let $\expolicymatchicopt$ be the policy given in Theorem~\ref{thm:policy-matching-ic-opt-reward}.
%     If $\expolicymatchicopt\in\seticexchangepolicies$ then $\expolicymatchicopt$ is IC optimal.
% \end{restatable}  



\section{Conclusion}
\label{sec:conclusion}

We studied the problem of incentivizing resource sharing between competing agents, without money, in settings where valuations are uncertain, 
revealed only through interaction, and sharing is irreversible. We formalized a novel model for incremental resource sharing, 
where a central planner proposes resource exchanges over multiple rounds. Our results show that carefully designed multi-round exchange 
algorithms can substantially outperform static policies by learning agents' valuations over multiple rounds. 

We characterized 
such optimal policies via a reduction to an optimal stopping problem on the lattice $\NN^2_{\geq 0}$. 
We then considered the problem of designing exchange algorithms when participants can leave early, and provided examples 
demonstrating why it is fundamentally more challenging.
These examples motivated necessary model conditions for which
optimal stopping and martingale methods can be used to characterize an optimal exchange algorithm.
% At the same time, we 
% identified fundamental limitations 
% when agents can strategically withdraw mid-policy, \kkcomment{expand}
Two open questions arise from our work. 
The first is whether or not an IC optimal algorithm is always guaranteed to exist, beyond the sufficient condition
we provide.
% While we have provided necessary conditions for optimal stopping and martingale methods
% to yield an IC algorithm, 
The second is how can our model and algorithms be extended to accommodate more than two competitors when there are complex interactions
between participants' utility functions, and thus collaboration incentives.

% Two open questions arising from our work is the  design of  \kkcomment{?} IC policies, \kkcomment{expand}, and the extension of our model 
% and methods to multiple agents.

% \section{IC outline}

% \begin{enumerate}
%     \item Section 4 outline
%     \item 
%     Say that the idea with optimal stopping is that once you reach a state, it doesn't matter how you got there, just where you are,
%     aka its kind of markovian. In these settings you can proceed greedily by pasting policies together.
%     However, this is not the case with IC policies because 
    
%     \item Give the 1/2 D example where IC pasting fails.

%     \item Argue that intuitively this represents a misalignment between welare and utility optimization.

%     \item In order to generalize the previous optimal stopping time algorithm approach to search for an IC policy, we need to make 
%     an assumption that IC policies can be combined in a gredy way. 

%     \item Assumption block

%     \item Under this assumption we can give a dp equation analog. First introduce new ess sup, State the proposition, say (-C) is a lower bound.
%     \item Proof sketch.

%     \item Explain that under this assumption we can contruct a policy that matches the performance of an optimal IC
    
%     \item Explain idea of theory, if we are allowed to stop vs not allowed to. Give theorem

%     \item Add remark that when this policy is IC, it is IC optimal.
    
%     \item Give proof citing proposition
    
%     \item \mycomment{NOte to self, revisit the stick note book marked, non trivially IC policy, see if it satisfies pasting 
%     assumptions? make sure it matches our policy?}

%     \item Add note that this really isn't practical to run. On going line of work to see if there are assumptions or special 
%     classes of welfare functions for which pasting is satisfied and its easier to run the alg

%     \item add corollary after main theorem saying if construction is IC then it is IC optimal
% \end{enumerate}  

% \begin{enumerate}
%     \item \mycomment{In the intro and techniques section probably needs a brief summary of what 
%     \citep{mandelbaum1981optimal} does so that the connection is not as abrupt}

%     \item 
%     say we're gonna start by reframing the problem. Say once an exchange of $(i,j)$ has occured, a know the value of the 
%     first i points and b knows the value of the first j points. To represent the information available after an exchange,
%     introduce sigma algebras for each of the states and define the family. And makke the observation that they are nested.
%     say intuitively this represents that you know more after a superset of the data has been exchanged. Initially, we know nothing
%     so $\Fcal_0=\cbr{\varnothing,\Omega}$

%     \item
%     % Say that given any algorithm, the corresponding sequence of states generated are a 
%     Say that corresponding to any alg there is the sequence (replace the $s_t$s with sigmas)
%     $\sigma_0,\ldots,\sigma_\tau$ which are a sequence of $\Scal$ 
%     valued random variables. More over, it is the case that $\sigma_{t+1}$ is ... measurable. This reflect that the choice 
%     the next exchange is only based on values observed from the previous exchange. 

%     \item Short proposition showing correspondence between algorithms and policies? 
%     \mycomment{feels unnecessary to make into prop but maybe reviewers would like to see more props?}

%     \item 
%     Give explanation saying that intuitively finding the optimal algorithm reduces to traversing 
%     the grid moving up/right at each step and determining when to stop. Therefore, we can use the 
%     language of (insert stopping paper) to formalize and analyze this problem.
%     The snell env is a classic concept from optimal stopping times 
%     We start with the snell env for partially ordered sets

%     \item Define snell envelope $X_s$ for $s\in\Scal$ (citing paper). Explain what it means 
%     intuitively, saying $\Pi_s$ is algorithms starting at state $s$.

%     \item Say in optimal stopping time theory problems, because $X_s$ is difficult to compute direclty,
%     you often prove a dynamic programming / recursive type of result. (cite paper) gives a formula
%     for paritially ordered sets, specializing to 0...m x 0...n gives (insert result for our
%     $N^2$ case). Then give proposition saying our optimal stopping algorithm?
%     % \item 
%     Have an algorithm block? Write it in algorithm form or policy form?
%     maybe instead of alg block, just have a theorem saying the specific algorithm is optimal 

%     \item 
%     Add remarks about how practically you would use techniques like MC methods to evaluate things 

%     \item Add note about how based on choice of welfare function you can make this indivisually
%     rational
%     \mycomment{add somewhere in assumptions that utility is 0 with no data exchanged.
%     without this IR being defined to be >0 doesn't make sense?
%     I think you only need to require $U^a_0=U^b_0=0$, don't worry about other states}
% \end{enumerate}  


% \subsection{An Optimal Algorithm when Participants cannot opt out}
% \paragraph{Reducting to an optimal stopping problem over N2 to solve}



% \begin{enumerate}
    
%     \item
%     Say how a algorithm can be viewed as a sequence of stopping points with respect to the 
%     sigma algebra generated by the values of data exchanged so far. Say that the stop decision
%     can be viewed as a stopping time w.r.t. the $\Gcal_t$.

%     \item
%     Define rewards for stopping as the welfare of the utilities.

%     \item
%     Therefore, finding the optimal data sharing algorithm reduces to finding the optimal stopping time
%     algorithm.

%     \item
%     Write out the algorithm to clairfy the interaction procedure

%     \item
%     Write a theorem saying that this is optimal. Proof involves using that rewards are bounded
%     and the stopping time is a.s. finite thus we can apply theorems 1 and 2 from the paper. 

%     \citet{mandelbaum1981optimal}

% \end{enumerate}  

% \paragraph{IR guarantees}

% \begin{enumerate}
%     \item Show that we automatically have ex ante IR for a sutible welfare function choice
%     under the optimal algorithm

%     \item 
%     Say we can also give a high probability ex post IR guarantee assuming we use the optimal
%     algorithm

%     \item Write and prove a proposition about this 
% \end{enumerate}  


% \section{Optimal algorithm when agents can opt}

% \begin{enumerate}
%     \item The model above assumes that agents either commit to following the data
%     sharing algorithm or do not, i.e. once they buy in they can't leave early.

%     \item But what happens if we assume agents can leave early. You might choose this model
%     for a real world setting if it is difficult to enforce agents to continue participating
%     or avoid agents who don't want to continue from submitting bad data
% \end{enumerate}  

% \subsection{Model changes when agents can opt out}

% \begin{enumerate}
%     \item Before agents made a single decision, whether or not to participate.
%     Now at each round they make a decision of whether or not to continue.
%     If either agent or the mechanism stops the sharing its over.
%     (Someone might ask what happens if both agents want to continue sharing but
%     the mechanism says to stop? In this case you can update the welfare function
%     so this won't happen?)

%     \item Define agent strategies as sequence of decisions

%     \item Remark that each agent's strategy corresponds to a stopping time
    
%     \item Use this to write out the expected welfare under a strategy profile in 
%     terms of a policy and agent strategy stopping times

%     \item Define a NE of policies under an algorithm

%     \item Define the set of IC policies using the min trick

%     \item State the problem of trying to find the optimal policy in this set,
%     i.e. the optimal IC policy.

% \end{enumerate}  

% \subsection{An Optimal Algorithm when Participants can opt out}


% \subsection{The opt IC algorithm}

% \begin{enumerate}
%     \item TBD
% \end{enumerate}  



% \section{Impossibility Results}

% \begin{enumerate}
%     \item 
%     Thus far we've assumed that agents report their observed values at each round
%     truthfully. The following result shows that if agents don't report truthfully
%     no algorithm can obtain the optimal expected welfare.

%     \item State proposition and give proof using two-flip.
    
% \end{enumerate}  


% \section{Running questions}

% \begin{enumerate}
%     \item I feel like there is a blind spot that I have for prior work in 
%     exchange of goods.

%     \item Should we consider other definitions of utility? It would make computations more annoying
%     but the results more general. (YES)

%     \item In the valuation functions example do I need to make sure everything is bounded? Technically to apply our
%     methods this is necessary but it would make the example cleaner if I didn't need to say $D_a,D_b$ are bounded and I
%     could write $V^a_0=V^b_0=-\infty$.

%     \item I don't even think independence of the value sequences are used anywhere?

% \end{enumerate}  

% \subsection{Small questions}

% \begin{enumerate}
%     \item Maybe give more examples of how value of data may not be known util it is used?

%     \item maybe show the calculations more in ex 1?

%     \item the general utility functions footnote?
    
%     \item should we specify probability space or require that the valuation functions are 
%     markov chains?

%     \item In the sharing alg def how much should we allude to the stopping time notation

%     \item Intellectual / property rights as another example of goods?

%     \item \mycomment{Do we need utilities to start at 0?}

%     \item
%     \mycomment{might definition to satisfy the property: to algs are = iff their dist over states is the same}

%     \item technically the paper gives a vague def for snell env as supermartingale ... and shows its equiv to 
%     what we say is the def

%     \item Weird that exchange policies is $\Pi^{\text{ex}}$ but IC policies, which is a subset, is $\Pi^{\text{IC}}$

%     \item fix stacking subscripts in pasting examples

%     % \item Is it ok that I chopen part of thm 1 from their paper?


% \end{enumerate}  

% \section{TODO}

% \begin{enumerate}
%     \item \mycomment{add assumption somewhere that tie breaking always chooses smaller or lexigraphically
%     smaller choice so that policies are consistent}
%     \item Resolve why the two flip calculation doesnt match the optimal policy calculation
%     \item cut down the descriptive text in section 2 
%     \item In existing results prove conditional dominated convergence,
%     proof is online
%     \item
%     \mycomment{add note about we must assume truthful reporting and its impossible to guarantee
%     this type of algorithm does better than doing nothing.
%     when agents are not truthful, welfare can be arbitrarily worse under the optimal algorithm } ?

%     \item Probably worth mentioning the connection of IC policies to solving for NE
%     \item the tildes and subscripts are hard to read, ex thm 3.5, fix this 
%     \item make sure not to take credit from paper in thm 3.5
%     \item I think the $\Ucal$ condition is probably redundant but harmless, maybe take this out later

%     \item I added the $\alpha$ policy notation in the truthfulness proof, make sure this is introduced somewhere

%     \item em dash creep

%     \item What should we call the sequence of exit strategies?

%     \item
%     \mycomment{In 4.2 to make more pasting more comparable, maybe make the 1d stopping time version not random and have it be stopping times which stop no 
%     earlier than fixed $t$}

%      \item \mycomment{we say that when random variables have finite support there are finitely many policies?
    %     is this only true assuming there is no external randomness to a policy?}
% \end{enumerate}  

% \section{To beef up submission}

% \begin{enumerate}
%     \item 
%     Add an equivalence proof of two definitions of algorithms

%     \item
%     Add the unbounded difference example

%     \item Work through the double flip example for the non-IC setting. write example followed 
%     by a proof block.

%     \item The IC alg might be more paletable if we can remove the $\Ucal$ part
    
%     \item 
%     Work through 4 flips example for the IC setting?

%     \item
%     Does pasting hold for reasonable welfare functions like min?


% \end{enumerate}  

\newpage
\bibliographystyle{ACM-Reference-Format} % EC formatting
\bibliography{references.bib}

\appendix
% \section{fill}

% 
% \begin{restatable*}{proposition}{optTruthfulVsOptUntruthful}
%     For any $k\geq 0$ there exists an environment such that  
%     \(
%         \EE\sbr{
%             \rewardss{
%                 \optstrategytt{
%                     \optstoptime
%                 }
%             }
%         }
%         -
%         \EE\sbr{
%             \rewardss{
%                 \opttruthstrategytt{
%                     \opttruthstoptime
%                 }
%             }
%         }
%         \geq k
%     \),
%     where 
%     $\optexpolicy$
%     and 
%     $\opttruthexpolicy$
%     are the optimal policies assuming and not assuming truthful reporting respectively.
% \end{restatable*}  


% \begin{restatable*}{proposition}{opt-truthful-vs-opt-untruthful2}
% % \label{prop:untruthfulreporting}
%     There exist a sequence of model instances such that 
%     \mycomment{typo, remove $k$, just one problem instance}
%     \[
%         \EE\sbr{
%             \rewardss{
%                 \optstrategytt{
%                     \optstoptime
%                 }
%             }
%         }
%         -
%         \EE\sbr{
%             \rewardss{
%                 \opttruthstrategytt{
%                     \opttruthstoptime
%                 }
%             }
%         }
%         \xrightarrow{k\rightarrow\infty}\infty
%     \]
%     where 
%     \(
%         \cbr{
%             \optexpolicy_k
%         }_{k\geq 0}
%         =
%         \cbr{
%             \rbr{
%                 \cbr{
%                     \optstrategyt
%                 }_{\timeind\geq 0},
%                 \optstoptime
%             }_k
%         }_{k\geq 0}
%     \)
%     and 
%     \(
%         \cbr{
%             \opttruthexpolicy_k
%         }_{k\geq 0}
%         =
%         \cbr{
%             \rbr{
%                 \cbr{
%                     \opttruthstrategyt
%                 }_{\timeind\geq 0},
%                 \opttruthstoptime
%             }_k
%         }_{k\geq 0}
%     \)
%     are sequences of optimal and optimal truthful policies for each model.
% \end{restatable*}  


% \begin{proof}
%     Let $\numdataa=\numdatab=1$.
%     For a particular $k\in\NN$, consider the valuations
%     % utilties, and welfare function,
%     % \(
%     % \)
%     \begin{align*}
%         \valaii{0}=\valbjj{0}=0
%         \qquad
%         \valaii{1}
%         =
%         \begin{cases}
%             2k
%             &
%             \text{w.p. } \frac{1}{2}
%             \\
%             % \frac{k}{2}
%             k
%             &
%             \text{w.p. } \frac{1}{2}
%         \end{cases}  
%         \qquad
%         \valbjj{1} = 2k
%     \end{align*}  
%     and let the competitors' utilities and the welfare function be given by
%     \begin{align*}
%         \utila(x,y) 
%         =
%         x-y
%         \qquad
%         \utilb(x,y)
%         =
%         \begin{cases}
%             0 
%             &
%             \text{if }
%             (y-x)\geq -k
%             %  k\geq x-y
%             \\
%             2(y-x)
%             &
%             \text{otherwise}
%         \end{cases}  
%         \qquad
%         \welfarefunc(x,y)=x+y .
%     \end{align*}
%     First consider when competitors are truthful. 
%     Under this model, the optimal policy 
%     \(
%         \optexpolicy_k
%         % =
%         % \rbr{
%         %     \cbr{
%         %         \optstrategyt
%         %     }_{\timeind\geq 0},
%         %     \optstoptime
%         % }_k
%     \)
%     first gives $\competitora$ competitor $\competitorb$'s point,
%     if 
%     % $\valaii{1}=\frac{k}{2}$ 
%     $\valaii{1}=k$ 
%     then sharing terminates, otherwise $\competitorb$ is given
%     $\competitora$'s point to ensure $\competitorb$'s utility does not lead to a negative welfare.
%     Thus the expected welfare under this policy is 
%     \[
%         \EE\sbr{
%             \rewardss{
%                 % \optstrategytt{
%                 %     \optstoptime_k
%                 % }^k
%                 \stopstateforpolicy{
%                     \optexpolicy_k
%                 }
%             }
%         }
%         =
%         \frac{1}{2}
%         \welfarefunc\rbr{
%             \utilass{
%                 (1,0)
%             },
%             \utilbss{
%                 (1,0)
%             }
%         }
%         +
%         \frac{1}{2}
%         \welfarefunc\rbr{
%             \utilass{
%                 (1,1)
%             },
%             \utilbss{
%                 (1,1)
%             }
%         }
%         =
%         \frac{1}{2}
%         \rbr{
%             % \frac{k}{2}
%             k
%             + 0 
%         }
%         +
%         \frac{1}{2}
%         \rbr{
%             0 + 0
%         } 
%         % =\frac{k}{4}
%         =\frac{k}{2}
%         .
%     \] 
%     Now consider when competitors are untruthful. 
%     In this setting, the optimal truthful policy will
%     not always 
%     give $\competitorb$ competitor
%     $\competitora$'s point. This is because doing so will always lead to a welfare of no more than 0,
%     thus such a policy is no better than one which exchanges zero resources.
%     Now consider a policy which only gives $\competitorb$ the point based on what is reported as $\valaii{1}$.
%     In this case $\competitora$ is incentivized to always report the value for $\valaii{1}$ for which $\competitorb$
%     does not receive the point as it increases their utility. However, this leads to a negative welfare.

%     Therefore, the optimal policy when competitors are untruthul recommends no exchange so 
%     \(
%         \EE\sbr{
%             \rewardss{
%                 \stopstateforpolicy{
%                     \opttruthexpolicy_k
%                 }
%             }
%         }
%         =0
%     \)
%     and thus
%     \[
%         \EE\sbr{
%             \rewardss{
%                 \optstrategytt{
%                     \optstoptime
%                 }
%             }
%         }
%         -
%         \EE\sbr{
%             \rewardss{
%                 \opttruthstrategytt{
%                     \opttruthstoptime
%                 }
%             }
%         }
%         \xrightarrow{k\rightarrow\infty}\infty .
%     \]
% \end{proof}  

\mycomment{before stating this result and when introducing truthfulness,
spend time explaining how the opt alg isn't truthful. If we put this in a section
4, then I can reference the example we walked through to claim optimality}

\mycomment{could add a figure of 2D stopping times which could
take up some space}

% \section{Omitted Proofs}
\label{app:omitted-proofs}

\mycomment{I don't think we need countability anywhere, I think the replacing trick with $B(S)$ actually works for $\geq$,
actually i think its needed}
\mycomment{I changed the IC theorem a big for the Y=equality}

\mycomment{get rid of $\Ucal$ if we have time, change to $u>s$}
\mycomment{I removed $\Ucal$, might need to add back if there is a problem later}

% \icDpEq*
% \begin{proof}
%     \mycomment{fill}
% \end{proof}  
\optTruthfulVsOptUntruthful*
\begin{proof}
    Let $\numdataa=\numdatab=1$.
    Consider the valuations
    \begin{align*}
        \valaii{0}=\valbjj{0}=0
        \qquad
        \valaii{1}
        =
        \begin{cases}
            4k
            &
            \text{w.p. } \frac{1}{2}
            \\
            % \frac{k}{2}
            2k
            &
            \text{w.p. } \frac{1}{2}
        \end{cases}  
        \qquad
        \valbjj{1} = 4k
    \end{align*}  
    and let the competitors' utilities and the welfare function be given by
    \begin{align*}
        \utila(x,y) 
        =
        x-y
        \qquad
        \utilb(x,y)
        =
        \begin{cases}
            0 
            &
            \text{if }
            (y-x)\geq -2k
            %  k\geq x-y
            \\
            2(y-x)
            &
            \text{otherwise}
        \end{cases}  
        \qquad
        \welfarefunc(x,y)=x+y .
    \end{align*}
    First consider when competitors are truthful. 
    Under this model, the optimal policy 
    \(
        \optexpolicy
    \)
    first gives $\competitora$ competitor $\competitorb$'s point,
    if 
    $\valaii{1}=2k$ 
    then sharing terminates, otherwise $\competitorb$ is given
    $\competitora$'s point to ensure $\competitorb$'s utility does not lead to a negative welfare.
    Thus the expected welfare under this policy is 
    \[
        \EE\sbr{
            \rewardss{
                \stopstateforpolicy{
                    \optexpolicy_k
                }
            }
        }
        =
        \frac{1}{2}
        \welfarefunc\rbr{
            \utilass{
                (1,0)
            },
            \utilbss{
                (1,0)
            }
        }
        +
        \frac{1}{2}
        \welfarefunc\rbr{
            \utilass{
                (1,1)
            },
            \utilbss{
                (1,1)
            }
        }
        =
        \frac{1}{2}
        \rbr{
            % \frac{k}{2}
            2k
            + 0 
        }
        +
        \frac{1}{2}
        \rbr{
            0 + 0
        } 
        % =\frac{k}{4}
        =
        k.
    \] 
    Now consider when competitors are untruthful. 
    In this setting, the optimal truthful policy will
    not always 
    give $\competitorb$ competitor
    $\competitora$'s point. This is because doing so will always lead to a welfare of no more than 0,
    thus such a policy is no better than one which exchanges zero resources.
    Now consider a policy which only gives $\competitorb$ the point based on what is reported as $\valaii{1}$.
    In this case $\competitora$ is incentivized to always report the value for $\valaii{1}$ for which $\competitorb$
    does not receive the point as it increases their utility. However, this leads to a negative welfare.

    Therefore, the optimal policy when competitors are untruthul recommends no exchange so 
    \(
        \EE\sbr{
            \rewardss{
                \stopstateforpolicy{
                    \opttruthexpolicy_k
                }
            }
        }
        =0
    \)
    and thus
    \(
        \EE\sbr{
            \rewardss{
                \optstrategytt{
                    \optstoptime
                }
            }
        }
        -
        \EE\sbr{
            \rewardss{
                \opttruthstrategytt{
                    \opttruthstoptime
                }
            }
        }
        =k
    \).
\end{proof}  
% \section{Borrowed and Known Results}
\section{Borrowed Results}

\begin{restatable}{theorem}{osasThm1}
    \label{thm:osas-dp-eq}
    \textup{
        (\citet{mandelbaum1981optimal})
    }
    % Let 
    % \(
    %     \osasdirectsuccessorss{
    %         \osasstate
    %     }
    %     =
    %     \cbr{
    %         u> \osasstate: \sbr{\osasstate,u}=\cbr{\osasstate,u}
    %     } 
    % \)
    % be the set of direct successors for $\osasstate$, then
    Snell's envelope for $\osassnellenvs$ exists and satisfies
    \begin{gather*}
        \osassnellenvs
        =
        \max
        \cbr{
            \osasrewards,
            \EE\sbr{
                \osassnellenvss{u}
                | 
                \osasstatefilts
            }
            : u\in
            \osasdirectsuccessorss{
                \osasstate
            }
        },
        \\
        % \quad
        % \text{and} \quad
        \EE\sbr{
            \osassnellenvss{0}
        }
        =
        \sup_{
            \osaspolicy\in\osaspolicies
        }
        \EE\sbr{
            \osasrewardss{
                \osasstrategytt{
                    \osasstoptime
                }
            }
        }
        .
    \end{gather*}  
\end{restatable}

\begin{restatable}{theorem}{osasThm2}
    \label{thm:osas-thm2}
    If
    \(
        \osaspolicy
        =
        \rbr{
            \cbr{
                \osasstrategyt
            }_{\timeind\geq 0},
            \osasstoptime
        }
    \)
    is a policy with the properties
    \begin{align}
        &
        \osasstrategytt{\timeind+1} \text{ is a direct successor of } \osasstrategyt
        \label{eq:osas-thm2-cond1}
        \\
        &
        \osasstoptime
        =
        \inf\cbr{
            t : 
            \osassnellenvss{
                \strategyt
            }
            =
            \osasrewardss{
                \strategyt
            }
        }
        \label{eq:osas-thm2-cond2}
        \\
        &
        \EE\sbr{
            \osassnellenvss{
                \osasstrategytt{
                    \timeind+1
                }
            }
            | \timefiltt
        }
        =
        \osassnellenvss{
            \osasstrategyt
        }
        \quad
        \text{on}
        \quad
        \cbr{\osasstoptime<\timeind}
        \label{eq:osas-thm2-cond3}
    \end{align}  
    and $\stoptime<\infty$ a.s., then $\pi$ is optimal. If there exists 
    some optimal policy then there exists an one with the above properties which is optimal. 
\end{restatable}  

\end{document}